% Les échanges inter-régionaux au Cameroun — Géographie Tle A — Cours signé OnBuch+
% Programme officiel MINESEC. Compiler avec : tectonic N09-echanges-inter-regionaux-cameroun.tex
\newcommand{\DOCMATIERE}{Géographie}
\newcommand{\DOCNIVEAU}{Tle A}
\newcommand{\DOCMODULE}{Module 1 : Le Cameroun}
\newcommand{\DOCLECON}{N09}
\newcommand{\DOCTITRE}{Les échanges inter-régionaux au Cameroun}
% OnBuch+ — préambule « cours signés » ENRICHI (compilé avec tectonic / XeLaTeX)
% Système visuel « Pop » d'OnBuch+ : fond crème (jamais blanc), bordures encre
% épaisses, ombres « dures » décalées, titres Archivo Black en capitales.
% Version MATHÉMATIQUES : graphes (pgfplots/tikz), boîtes pédagogiques
% (définition, propriété, méthode, attention, le savais-tu, exercice résolu,
% exercice, TP, bilan), page de garde et signature OnBuch+ ancrée en bas.
%
% Macros attendues dans main.tex AVANT \input{preamble} :
%   \DOCMATIERE  \DOCNIVEAU  \DOCMODULE  \DOCLECON (numéro)  \DOCTITRE
\documentclass[11pt]{article}

\usepackage{fontspec}
\usepackage[french]{babel}
\usepackage[a4paper,margin=2cm,top=3.4cm,bottom=2.8cm,headheight=1.4cm,footskip=1.3cm]{geometry}
\usepackage{amsmath,amssymb,amsfonts}
\usepackage{mathtools} % \xmapsto, \coloneqq... souvent utilisés par les modèles
\usepackage{cancel} % simplifications barrées (bilans de forces)
% \abs{z} (module, valeur absolue) et \norm{u} (norme), utilisés par les modèles
\providecommand{\abs}{}\let\abs\relax\DeclarePairedDelimiter{\abs}{\lvert}{\rvert}
\providecommand{\norm}{}\let\norm\relax\DeclarePairedDelimiter{\norm}{\lVert}{\rVert}
\usepackage[table]{xcolor} % [table] : fournit aussi \rowcolors (alternance de lignes)
\usepackage{pagecolor}
\usepackage{tikz}
\usetikzlibrary{arrows.meta,positioning,calc,shapes.geometric,shapes.misc,shapes.arrows,decorations.pathmorphing,decorations.pathreplacing,decorations.markings,patterns,shadows,mindmap,trees,fit,backgrounds,matrix,intersections,angles,quotes}
\usepackage{pgfplots}
\usepackage[version=4]{mhchem} % équations nucléaires \ce{^{235}_{92}U}
\usepackage[european,siunitx]{circuitikz} % schémas électriques
\pgfplotsset{compat=1.18}
\usepgfplotslibrary{fillbetween,groupplots} % aires entre courbes (calcul intégral)
\usepackage{enumitem}
\usepackage{fancyhdr}
% [most] charge toutes les bibliothèques tcolorbox (listings, minted, raster,
% documentation...) et gonfle inutilement la mémoire TeX sur un document long
% (~300 boîtes) : on ne charge que ce qui est réellement utilisé.
\usepackage[skins,breakable]{tcolorbox}
\usepackage{graphicx}
\graphicspath{{/home/user/t-t/geo-onbuch/images/}}
\usepackage{colortbl}
\usepackage{booktabs}
\usepackage{tabularx}
\usepackage{multirow}
\usepackage{diagbox} % case d'en-tête diagonale des tableaux à double entrée
% Colonnes extensibles centrée / gauche / droite (C, L, R), souvent utilisées par les modèles
\newcolumntype{C}{>{\centering\arraybackslash}X}
\newcolumntype{L}{>{\raggedright\arraybackslash}X}
\newcolumntype{R}{>{\raggedleft\arraybackslash}X}
\usepackage{array}
\usepackage{multicol}
\usepackage{siunitx}
% parse-numbers=false : accepte aussi \SI{2,5 \times 10^{-3}}{...}
\sisetup{locale=FR,output-decimal-marker={,},group-minimum-digits=4,parse-numbers=false}
\DeclareSIUnit\ohm{\ensuremath{\symup{\Omega}}} % U+2126 absent de la police
\DeclareSIUnit\division{div} % réglages d'oscilloscope (ms/div, V/div)
\DeclareSIUnit\tr{tr} % tours (vitesse de rotation, ex. \SI{1500}{\tr\per\minute})
\DeclareSIUnit\molar{M} % concentration molaire (ex. \SI{3}{\molar}, \SI{1}{\micro\molar})
\DeclareSIUnit\an{an} % année (le modèle confond parfois avec \year, un compteur TeX) — ex. \SI{5}{\kilo\meter\per\an}
\DeclareSIUnit\FCFA{FCFA} % franc CFA (ex. \SI{500}{\FCFA\per\kilo\gram})
\DeclareSIUnit\degreeBrix{\textdegree Brix} % teneur en sucre d'un sirop/jus (réfractométrie)
\DeclareSIUnit\month{mois} % le modèle utilise parfois \month, un compteur TeX, comme unité de durée
\DeclareSIUnit\microliter{\micro\liter} % le modèle écrit parfois \microliter en un seul token
\DeclareSIUnit\million{millions} % dénombrement (ex. \SI{15}{\million\per\mL} spermatozoïdes)
\usepackage{qrcode}
% \degree : commande fréquemment utilisée par les modèles pour un angle ou une
% température en dehors de \SI{}{} (non fournie par les paquets chargés).
\providecommand{\degree}{^{\circ}}
% \qed (fin de démonstration), utilisé par les modèles sans amsthm
\providecommand{\qed}{\hfill$\square$}
% \barre : double barre des valeurs interdites dans un tableau de variations (tkz-tab)
\providecommand{\barre}{\ensuremath{\|}}
% environnement proof (amsthm non chargé) utilisé par les modèles
\ifdefined\proof\else\newenvironment{proof}[1][Démonstration]{\par\noindent\textit{#1.}\ }{\qed\par}\fi
\usepackage{lastpage}
\usepackage{hyperref}
\hypersetup{hidelinks}

% ============================================================
% Palette « Pop » OnBuch+ — mêmes hex que la classe P (plus_ui.dart)
% ============================================================
\definecolor{poporange}{HTML}{F59321}
\definecolor{popdark}{HTML}{E07A0C}
\definecolor{popcream}{HTML}{FFF6E8}
\definecolor{popink}{HTML}{1C1714}
\definecolor{poppurple}{HTML}{7A5AE0}
\definecolor{popgreen}{HTML}{2E9E6A}
\definecolor{poppink}{HTML}{E0457B}
\definecolor{popblue}{HTML}{2F7DE1}
\definecolor{popgold}{HTML}{C98A12}
\definecolor{popmuted}{HTML}{8A7F76}
\definecolor{popline}{HTML}{EADFD2}
\definecolor{popgray}{HTML}{8A7F76}
\colorlet{popgrey}{popgray}
% Alias tolérants (le modèle nomme parfois les couleurs différemment)
\colorlet{popwhite}{white}
\colorlet{popblack}{popink}
\colorlet{popred}{poppink}
\colorlet{popyellow}{popgold}
% Teintes claires (fonds de tableaux/encadrés) — le modèle nomme parfois
% les couleurs avec un suffixe L (ex. popblueL) sans les avoir définies.
\colorlet{poporangeL}{poporange!20}
\colorlet{popdarkL}{popdark!20}
\colorlet{poppurpleL}{poppurple!20}
\colorlet{popgreenL}{popgreen!20}
\colorlet{poppinkL}{poppink!20}
\colorlet{popblueL}{popblue!20}
\colorlet{popgoldL}{popgold!20}
\colorlet{popredL}{popred!20}
\colorlet{popgrayL}{popgray!20}
% Teintes claires pour fonds de tableaux / zones
\colorlet{popblueL}{popblue!10!white}
\colorlet{poporangeL}{poporange!14!white}
\colorlet{popgreenL}{popgreen!12!white}
\colorlet{poppurpleL}{poppurple!12!white}
\colorlet{poppinkL}{poppink!12!white}
\colorlet{popgoldL}{popgold!14!white}

\pagecolor{popcream}

% ============================================================
% Typographie
% ============================================================
\defaultfontfeatures{Ligatures=TeX}
\setmainfont{PlusJakartaSans-Medium.ttf}[Path=../../../fonts/,BoldFont=PlusJakartaSans-Bold.ttf]
% Maths en sans-serif (Fira Math) avec chiffres et lettres droites de Plus
% Jakarta, pour que les formules se fondent dans le texte.
\usepackage[math-style=ISO,bold-style=ISO]{unicode-math}
\setmathfont{FiraMath-Regular.otf}[Path=../../../fonts/]
\setmathfont{PlusJakartaSans-Medium.ttf}[Path=../../../fonts/,range={up/num,up/{Latin,latin},\mathup}]
\usepackage{icomma} % virgule décimale sans espace en mode math
% Caractères mathématiques Unicode absents des polices de texte (Plus Jakarta,
% Archivo Black) : rendus par leur équivalent mathématique, en texte comme en
% formule — sinon ils s'affichent en carré vide (ex. « Arithmétique dans ℤ »).
\usepackage{newunicodechar}
\newunicodechar{ℝ}{\ensuremath{\mathbb{R}}}
\newunicodechar{ℤ}{\ensuremath{\mathbb{Z}}}
\newunicodechar{ℂ}{\ensuremath{\mathbb{C}}}
\newunicodechar{ℕ}{\ensuremath{\mathbb{N}}}
\newunicodechar{ℚ}{\ensuremath{\mathbb{Q}}}
\newunicodechar{𝔻}{\ensuremath{\mathbb{D}}}
\newunicodechar{α}{\ensuremath{\alpha}}
\newunicodechar{β}{\ensuremath{\beta}}
\newunicodechar{γ}{\ensuremath{\gamma}}
\newunicodechar{δ}{\ensuremath{\delta}}
\newunicodechar{ε}{\ensuremath{\varepsilon}}
\newunicodechar{θ}{\ensuremath{\theta}}
\newunicodechar{λ}{\ensuremath{\lambda}}
\newunicodechar{μ}{\ensuremath{\mu}}
\newunicodechar{σ}{\ensuremath{\sigma}}
\newunicodechar{τ}{\ensuremath{\tau}}
\newunicodechar{φ}{\ensuremath{\varphi}}
\newunicodechar{ψ}{\ensuremath{\psi}}
\newunicodechar{ω}{\ensuremath{\omega}}
\newunicodechar{Σ}{\ensuremath{\Sigma}}
% majuscules produites par \MakeUppercase dans les titres (α-aminés -> Α)
\newunicodechar{Α}{\ensuremath{\alpha}}
\newunicodechar{Β}{\ensuremath{\beta}}
\newunicodechar{Γ}{\ensuremath{\Gamma}}
\newunicodechar{Δ}{\ensuremath{\Delta}}
\newunicodechar{Ε}{\ensuremath{\varepsilon}}
\newunicodechar{Θ}{\ensuremath{\Theta}}
\newunicodechar{Λ}{\ensuremath{\Lambda}}
\newunicodechar{Μ}{\ensuremath{\mu}}
\newunicodechar{Τ}{\ensuremath{\tau}}
\newunicodechar{Φ}{\ensuremath{\Phi}}
\newunicodechar{Ψ}{\ensuremath{\Psi}}
\newunicodechar{Ω}{\ensuremath{\Omega}}
\newunicodechar{↦}{\ensuremath{\mapsto}}
\newunicodechar{∈}{\ensuremath{\in}}
\newunicodechar{∉}{\ensuremath{\notin}}
\newunicodechar{∧}{\ensuremath{\wedge}}
\newunicodechar{⇔}{\ensuremath{\Leftrightarrow}}
\newunicodechar{⟺}{\ensuremath{\Longleftrightarrow}}
\newunicodechar{⇒}{\ensuremath{\Rightarrow}}
\newunicodechar{⟹}{\ensuremath{\Longrightarrow}}
\newunicodechar{∘}{\ensuremath{\circ}}
\newunicodechar{∪}{\ensuremath{\cup}}
\newunicodechar{∩}{\ensuremath{\cap}}
\newunicodechar{≡}{\ensuremath{\equiv}}
\newunicodechar{⊂}{\ensuremath{\subset}}
\newunicodechar{⊥}{\ensuremath{\perp}}
\newunicodechar{∃}{\ensuremath{\exists}}
\newunicodechar{∀}{\ensuremath{\forall}}
\newunicodechar{∛}{\ensuremath{\sqrt[3]{\,}}}
\newunicodechar{∜}{\ensuremath{\sqrt[4]{\,}}}
\newunicodechar{⃗}{\ensuremath{{}^{\to}}}
\newunicodechar{ⁿ}{\ifmmode^{n}\else\textsuperscript{\ensuremath{n}}\fi}
\newunicodechar{ˣ}{\ifmmode^{x}\else\textsuperscript{\ensuremath{x}}\fi}
\newunicodechar{ᵇ}{\ifmmode^{b}\else\textsuperscript{\ensuremath{b}}\fi}
\newunicodechar{ʳ}{\ifmmode^{r}\else\textsuperscript{\ensuremath{r}}\fi}
\newunicodechar{ᵘ}{\ifmmode^{u}\else\textsuperscript{\ensuremath{u}}\fi}
\newunicodechar{ʸ}{\ifmmode^{y}\else\textsuperscript{\ensuremath{y}}\fi}
\newunicodechar{ᵃ}{\ifmmode^{a}\else\textsuperscript{\ensuremath{a}}\fi}
\newunicodechar{ᵉ}{\ifmmode^{e}\else\textsuperscript{\ensuremath{e}}\fi}
\newunicodechar{ᵏ}{\ifmmode^{k}\else\textsuperscript{\ensuremath{k}}\fi}
\newunicodechar{ᵖ}{\ifmmode^{p}\else\textsuperscript{\ensuremath{p}}\fi}
\newunicodechar{ᴺ}{\ifmmode^{N}\else\textsuperscript{\ensuremath{N}}\fi}
\newunicodechar{ᶜ}{\ifmmode^{c}\else\textsuperscript{\ensuremath{c}}\fi}
\newunicodechar{ʰ}{\ifmmode^{h}\else\textsuperscript{\ensuremath{h}}\fi}
\newunicodechar{ᵠ}{\ifmmode^{q}\else\textsuperscript{\ensuremath{q}}\fi}
\newunicodechar{ₙ}{\ifmmode_{n}\else\textsubscript{\ensuremath{n}}\fi}
\newunicodechar{ₐ}{\ifmmode_{a}\else\textsubscript{\ensuremath{a}}\fi}
\newunicodechar{ᵢ}{\ifmmode_{i}\else\textsubscript{\ensuremath{i}}\fi}
\newunicodechar{ᵣ}{\ifmmode_{r}\else\textsubscript{\ensuremath{r}}\fi}
\newunicodechar{ₖ}{\ifmmode_{k}\else\textsubscript{\ensuremath{k}}\fi}
\newunicodechar{ᵦ}{\ifmmode_{\beta}\else\textsubscript{\ensuremath{\beta}}\fi}
\newunicodechar{ₓ}{\ifmmode_{x}\else\textsubscript{\ensuremath{x}}\fi}
\newfontfamily\popdisplay{ArchivoBlack-Regular.ttf}[Path=../../../fonts/]
\newfontfamily\poplogo{SpaceGrotesk-Bold.ttf}[Path=../../../fonts/]
\newfontfamily\popsemi{PlusJakartaSans-SemiBold.ttf}[Path=../../../fonts/]
\newfontfamily\popxbold{PlusJakartaSans-ExtraBold.ttf}[Path=../../../fonts/]
\newfontfamily\popxboldls{PlusJakartaSans-ExtraBold.ttf}[Path=../../../fonts/,LetterSpace=3.0]

\setlist[enumerate]{leftmargin=1.7em,itemsep=5pt,topsep=4pt}
\setlist[itemize]{leftmargin=1.7em,itemsep=4pt,topsep=4pt,label={\color{poporange}\textbullet}}
\setlength{\parindent}{0pt}
\setlength{\parskip}{5pt}
\renewcommand{\arraystretch}{1.25}

% pgfplots : style Pop par défaut
\pgfplotsset{
  every axis/.append style={axis line style={popink,line width=1pt},tick style={popink},
    grid style={popline,line width=0.5pt},label style={font=\small},tick label style={font=\footnotesize}},
  popcurve/.style={poporange,line width=1.8pt,no markers,smooth},
}
% Même style disponible dans un \draw TikZ simple (hors pgfplots)
\tikzset{popcurve/.style={poporange,line width=1.8pt}}
\tikzset{popgrid/.style={popline,very thin}}
\colorlet{popcurve}{poporange} % le nom du style est parfois utilisé comme couleur

% ============================================================
% Pill / squiggle
% ============================================================
\newcommand{\poppill}[3]{%
  \tikz[baseline=-0.55ex]\node[draw=popink,line width=1.2pt,rounded corners=2.6pt,%
    fill=#2,inner xsep=6.5pt,inner ysep=3pt]{%
    \popxboldls\fontsize{7}{7}\selectfont\color{#3}\MakeUppercase{#1}};%
}
\newcommand{\popsquiggle}[1]{%
  \tikz[baseline=-0.4ex]{
    \draw[#1,line width=2.6pt,line cap=round]
      plot [smooth] coordinates {(0,0.09) (0.34,0.01) (0.68,0.21) (1.02,0.01) (1.36,0.10)};
  }%
}
% Mot-clé surligné (marqueur orange sous le texte)
\newcommand{\cle}[1]{{\popxbold\color{popdark}#1}}

% ============================================================
% En-tête / pied
% ============================================================
\pagestyle{fancy}
\fancyhf{}
\renewcommand{\headrulewidth}{0pt}
\renewcommand{\footrulewidth}{0pt}
\lhead{\raisebox{-0.15ex}{\poplogo\fontsize{13}{13}\selectfont\color{popink}OnBuch\color{poporange}+}}
\rhead{\poppill{\DOCMATIERE\ \textbullet\ \DOCNIVEAU\ \textbullet\ Leçon \DOCLECON}{white}{popink}}
\lfoot{\popxbold\fontsize{7.6}{7.6}\selectfont\color{popmuted}\MakeUppercase{Cours signé OnBuch+}\ \textbullet\ {\popsemi\DOCTITRE}}
\rfoot{\tikz[baseline=-0.6ex]\node[rounded corners=8.5pt,fill=popink,minimum height=17pt,minimum width=17pt,inner xsep=5pt,inner sep=0]{\popxbold\fontsize{8}{8}\selectfont\color{popcream}\thepage\,/\,\pageref*{LastPage}};}

% ============================================================
% Titres
% ============================================================
\newcommand{\chaptitre}[2]{%
  \begin{tcolorbox}[enhanced,colback=poporange,colframe=popink,boxrule=2.6pt,arc=11pt,
      drop shadow=popink,left=15pt,right=15pt,top=12pt,bottom=11pt,before skip=3pt,after skip=18pt]
    \poppill{#2}{popink}{popcream}\par\vspace{9pt}
    {\raggedright\hyphenpenalty=10000\exhyphenpenalty=10000\popdisplay\fontsize{22}{24}\selectfont\color{popcream}\MakeUppercase{#1}\par}
  \end{tcolorbox}
}
% Section numérotée : pastille orange avec le numéro + titre Archivo
\newcounter{popsec}
\newcommand{\coursec}[1]{%
  \stepcounter{popsec}\setcounter{popsub}{0}%
  \par\vspace{16pt}\noindent
  \begin{minipage}{\linewidth}
  \tikz[baseline=-0.7ex]\node[draw=popink,line width=1.6pt,fill=poporange,rounded corners=4pt,minimum size=22pt,inner sep=0]{\popdisplay\fontsize{12}{12}\selectfont\color{popcream}\thepopsec};%
  \hspace{8pt}{\popdisplay\fontsize{15}{17}\selectfont\color{popink}\MakeUppercase{#1}}\par
  \vspace{4pt}\noindent{\color{poporange}\rule{36pt}{3.6pt}}
  \end{minipage}\par\nopagebreak\vspace{8pt}
}
\newcounter{popsub}
\newcommand{\courssub}[1]{%
  \stepcounter{popsub}%
  \par\vspace{9pt}\noindent{\popxbold\fontsize{11.5}{13}\selectfont\color{popink}{\color{poporange}\thepopsec.\thepopsub}\hspace{6pt}#1}\par\nopagebreak\vspace{4pt}
}

% ============================================================
% Boîtes pédagogiques — même recette : fond blanc, bordure épaisse colorée,
% ombre dure encre, badge plein. Titre optionnel [#1].
% ============================================================
\tcbset{popbase/.style={breakable,enhanced,colback=white,boxrule=2.3pt,arc=9pt,
  shadow={3pt}{-3pt}{0pt}{popink},left=14pt,right=14pt,top=11pt,bottom=11pt,before skip=11pt,after skip=15pt,
  fontupper=\popsemi\fontsize{10.6}{14.5}\selectfont\color{popink},
  fontlower=\popsemi\fontsize{10.6}{14.5}\selectfont\color{popink}}}
% Coche dessinée (le glyphe ✓ n'existe pas dans les polices du document)
\newcommand{\popcheck}[1]{\tikz[baseline=-0.5ex]\draw[#1,line width=1.6pt,line cap=round,line join=round](-0.13,0.0)--(-0.04,-0.09)--(0.13,0.1);}
\providecommand{\coche}{\popcheck{popgreen}}
\providecommand{\resultat}[1]{\par\smallskip\noindent\fcolorbox{popgreen}{popgreenL}{\parbox{\dimexpr\linewidth-2\fboxsep-2\fboxrule\relax}{\textbf{Résultat.} #1}}\par\smallskip}
\newcommand{\popboxhead}[3]{\poppill{#1}{#2}{white}\ifx\relax#3\relax\else\hspace{6pt}{\popxbold\fontsize{10.6}{12}\selectfont\color{popink}#3}\fi\par\vspace{7pt}}

\newtcolorbox{aretenir}[1][]{popbase,colframe=popgreen,before upper={\popboxhead{À retenir}{popgreen}{#1}}}
\newtcolorbox{exemplebox}[1][]{popbase,colframe=poporange,before upper={\popboxhead{Exemple}{poporange}{#1}}}
\newtcolorbox{definition}[1][]{popbase,colframe=popblue,before upper={\popboxhead{Définition}{popblue}{#1}}}
\newtcolorbox{propriete}[1][]{popbase,colframe=poppurple,before upper={\popboxhead{Propriété}{poppurple}{#1}}}
\newtcolorbox{methode}[1][]{popbase,colframe=popgold,before upper={\popboxhead{Méthode}{popgold}{#1}}}
\newtcolorbox{attention}[1][]{popbase,colframe=poppink,before upper={\popboxhead{Attention}{poppink}{#1}}}
\newtcolorbox{experience}[1][]{popbase,colframe=popblue,colback=popblueL,before upper={\popboxhead{Expérience}{popblue}{#1}}}
% « Le savais-tu ? » : carte violette pleine (contexte, culture, Cameroun)
\newtcolorbox{savaistu}[1][]{popbase,colback=poppurple,colframe=popink,
  fontupper=\popsemi\fontsize{10.4}{14.2}\selectfont\color{white},
  before upper={\poppill{Le savais-tu\,?}{white}{poppurple}\ifx\relax#1\relax\else\hspace{6pt}{\popxbold\color{white}#1}\fi\par\vspace{7pt}}}
% Exercice résolu : énoncé en haut, \tcblower puis la solution sur fond crème
\newcounter{popexo}\newcounter{popexr}
\newtcolorbox{exoresolu}[1][]{popbase,colframe=poporange,colbacklower=popcream,
  segmentation style={popink,line width=1.2pt,dashed},
  before upper={\stepcounter{popexr}\popboxhead{Exercice résolu \thepopexr}{poporange}{#1}},
  before lower={\poppill{Solution}{popink}{popcream}\par\vspace{6pt}}}
% Exercice (non résolu en place) — la correction est dans la partie Corrigés
\newtcolorbox{exercice}[1][]{popbase,colframe=popink,
  before upper={\stepcounter{popexo}\popboxhead{Exercice \thepopexo}{popink}{#1}}}
% Corrigé
\newtcolorbox{corrige}[1][]{popbase,colframe=popgreen,colback=popgreenL,
  before upper={\popboxhead{Corrigé}{popgreen}{#1}}}
% Objectifs / prérequis / bilan
\newtcolorbox{objectifs}{popbase,colframe=popink,colback=poporangeL,
  before upper={\popboxhead{Objectifs de la leçon}{popink}{}}}
\newtcolorbox{prerequis}{popbase,colframe=popink,colback=popblueL,
  before upper={\popboxhead{Prérequis}{popblue}{}}}
\newtcolorbox{bilan}{popbase,colframe=popink,colback=white,boxrule=2.8pt,
  before upper={\popboxhead{Fiche bilan}{poporange}{L'essentiel à connaître pour l'examen}}}
% Figure encadrée avec légende
\newtcolorbox{popfigure}[1][]{enhanced,breakable=false,colback=white,colframe=popline,boxrule=1.4pt,arc=8pt,
  left=10pt,right=10pt,top=10pt,bottom=8pt,before skip=10pt,after skip=12pt,halign=center,
  lower separated=false,halign lower=center,
  fontlower=\popsemi\fontsize{9}{11.5}\selectfont\color{popmuted},
  #1}
\newcommand{\legende}[1]{\par\vspace{6pt}{\popsemi\fontsize{9}{11.5}\selectfont\color{popmuted}#1}}

% ============================================================
% Signature OnBuch+
% ============================================================
\newcommand{\poplogoinline}[1]{{\poplogo\fontsize{#1}{#1}\selectfont\color{popink}OnBuch\color{poporange}+}}
\newcommand{\signatureonbuch}{%
  \begin{tcolorbox}[enhanced,colback=popink,colframe=popink,boxrule=2.2pt,arc=10pt,
      drop shadow=poporange,left=14pt,right=14pt,top=10pt,bottom=10pt,before skip=0pt,after skip=0pt]
    \tikz[baseline=-0.7ex]\node[circle,fill=popgreen,draw=popcream,line width=1.3pt,minimum size=22pt,inner sep=0]{\popcheck{white}};%
    \hspace{9pt}\begin{minipage}[c]{0.62\linewidth}
      {\popxbold\fontsize{12}{13}\selectfont\color{popcream}Signé\ \textbullet\ {\poplogo OnBuch\color{poporange}+}}\par\vspace{2pt}
      {\raggedright\popsemi\fontsize{8}{10.5}\selectfont\color{popline}\MakeUppercase{Équipe pédagogique OnBuch+}\\\MakeUppercase{Conforme au programme officiel MINESEC}\par}
    \end{minipage}\hfill
    {\poplogo\fontsize{22}{22}\selectfont\color{popcream}OnBuch\color{poporange}+}
  \end{tcolorbox}%
}

% ============================================================
% Page de garde
% #1 titre · #2 matière · #3 niveau · #4 module · #5 n° leçon · #6 durée ·
% #7 description / objectifs (texte ou itemize)
% ============================================================
\newcommand{\pagegarde}[7]{%
  \thispagestyle{empty}
  \newgeometry{a4paper,margin=1.7cm,top=1.6cm,bottom=1.4cm}%
  \begin{tikzpicture}[remember picture,overlay]
    \fill[poporange!18] ([xshift=-3.2cm,yshift=-3.6cm]current page.north east) circle (4.6cm);
    \fill[poppurple!14] ([xshift=2.4cm,yshift=7.6cm]current page.south west) circle (3.8cm);
  \end{tikzpicture}%
  {\poplogo\fontsize{30}{30}\selectfont\color{popink}OnBuch\color{poporange}+}\hfill
  \raisebox{0.6ex}{\poppill{Cours signé \textbullet\ Premium}{popink}{popcream}}\par
  \vspace{1pt}\popsquiggle{poppurple}\par
  \vspace{8pt}
  \begin{tcolorbox}[enhanced,colback=poporange,colframe=popink,boxrule=2.8pt,arc=13pt,
      drop shadow=popink,left=18pt,right=18pt,top=12pt,bottom=13pt,after skip=10pt]
    \poppill{#2 \textbullet\ #3}{popink}{popcream}\hspace{5pt}\poppill{#4}{white}{popink}\par\vspace{8pt}
    {\popdisplay\fontsize{14}{14}\selectfont\color{popink}LEÇON #5}\par\vspace{4pt}
    {\raggedright\hyphenpenalty=10000\exhyphenpenalty=10000\popdisplay\fontsize{25}{27}\selectfont\color{popcream}\MakeUppercase{#1}\par}
    \vspace{8pt}
    {\popxbold\fontsize{9.5}{9.5}\selectfont\color{popink}\MakeUppercase{Durée conseillée : #6}}
  \end{tcolorbox}
  \begin{tcolorbox}[enhanced,colback=white,colframe=popink,boxrule=2.2pt,arc=10pt,
      drop shadow=popink,left=16pt,right=16pt,top=10pt,bottom=10pt,after skip=8pt]
    \poppill{Dans cette leçon}{poppurple}{white}\par\vspace{5pt}
    \popsemi\fontsize{9.3}{12.6}\selectfont\color{popink}#7
  \end{tcolorbox}
  \noindent\begin{minipage}[t]{0.487\linewidth}
    \begin{tcolorbox}[enhanced,colback=popgreen,colframe=popink,boxrule=2.2pt,arc=10pt,
        drop shadow=popink,left=11pt,right=11pt,top=10pt,bottom=10pt,height=3.1cm,valign=center]
      \tcbox[colback=white,colframe=popink,boxrule=1.3pt,arc=5pt,boxsep=0pt,nobeforeafter,
        left=3pt,right=3pt,top=3pt,bottom=3pt,valign=center]{\qrcode[height=1.9cm]{https://play.google.com/store/apps/details?id=com.luvvix.onbuch&hl=fr}}%
      \hspace{8pt}\begin{minipage}[c]{0.5\linewidth}
        {\popdisplay\fontsize{11}{12.5}\selectfont\color{white}\MakeUppercase{Télécharge OnBuch}}\par\vspace{4pt}
        {\raggedright\popsemi\fontsize{8.4}{11}\selectfont\color{white}Gratuit \textbullet\ Google Play\\Tuteur IA, annales, concours, résultats\par}
      \end{minipage}
    \end{tcolorbox}
  \end{minipage}\hfill
  \begin{minipage}[t]{0.487\linewidth}
    \begin{tcolorbox}[enhanced,colback=poppurple,colframe=popink,boxrule=2.2pt,arc=10pt,
        drop shadow=popink,left=11pt,right=11pt,top=10pt,bottom=10pt,height=3.1cm,valign=center]
      \tikz[baseline=-0.7ex]\node[circle,fill=white,draw=popink,line width=1.3pt,minimum size=1.6cm,inner sep=0]{\popdisplay\fontsize{24}{24}\selectfont\color{poppurple}+};%
      \hspace{8pt}\begin{minipage}[c]{0.6\linewidth}
        {\popdisplay\fontsize{11}{12.5}\selectfont\color{white}\MakeUppercase{Passe à OnBuch+}}\par\vspace{4pt}
        {\raggedright\popsemi\fontsize{8.4}{11}\selectfont\color{white}Tous les cours signés, Tuteur IA illimité, fiches PDF — onglet OnBuch+ de l'app\par}
      \end{minipage}
    \end{tcolorbox}
  \end{minipage}
  \vspace{8pt}
  \popcontenu
  \vfill
  \signatureonbuch
  \restoregeometry
  \newpage
}

% Rangée « Ce cours contient » (page de garde)
\newcommand{\poptile}[3]{% #1 couleur #2 gros texte #3 légende
  \begin{tcolorbox}[enhanced,colback=#1,colframe=popink,boxrule=2pt,arc=9pt,shadow={2.5pt}{-2.5pt}{0pt}{popink},
      left=6pt,right=6pt,top=9pt,bottom=9pt,halign=center,valign=center,height=1.75cm,width=\linewidth,nobeforeafter]
    {\popdisplay\fontsize{12.5}{13}\selectfont\color{white}\MakeUppercase{#2}}\par\vspace{3pt}
    {\centering\popsemi\fontsize{7.8}{9.5}\selectfont\color{white}#3\par}
  \end{tcolorbox}}
\newcommand{\popcontenu}{%
  \par\noindent{\popxboldls\fontsize{7.5}{7.5}\selectfont\color{popmuted}CE COURS SIGNÉ CONTIENT}\par\vspace{7pt}
  \noindent\begin{minipage}[t]{0.235\linewidth}\poptile{popblue}{Cours}{complet, illustré, pas à pas}\end{minipage}\hfill
  \begin{minipage}[t]{0.235\linewidth}\poptile{poporange}{Méthodes}{et exercices résolus}\end{minipage}\hfill
  \begin{minipage}[t]{0.235\linewidth}\poptile{poppink}{Exercices}{type examen + corrigés}\end{minipage}\hfill
  \begin{minipage}[t]{0.235\linewidth}\poptile{popgreen}{Fiche bilan}{l'essentiel à retenir}\end{minipage}\par}

% Sommaire
\newtcolorbox{sommaire}{enhanced,colback=white,colframe=popink,boxrule=2.2pt,arc=10pt,
  drop shadow=popink,left=18pt,right=18pt,top=13pt,bottom=13pt,before skip=6pt,after skip=20pt,
  before upper={\poppill{Sommaire}{poppurple}{white}\par\vspace{9pt}\popsemi\fontsize{10.6}{14.5}\selectfont\color{popink}}}

% Signature de fin de cours — poussée tout en bas de la dernière page
\newcommand{\findecours}{%
  \par\vfill\nopagebreak
  \begin{center}\popsquiggle{poporange}\end{center}\vspace{4pt}
  \signatureonbuch
}
\AtBeginDocument{\def\llbracket{\mathopen{[\mkern-3.5mu[}}\def\rrbracket{\mathclose{]\mkern-3.5mu]}}} % intervalles d'entiers (glyphe ⟦ absent de la police)
% Glyphes absents de Fira Math : redéfinis à partir de glyphes disponibles
\AtBeginDocument{%
  \def\setminus{\mathbin{\backslash}}\let\smallsetminus\setminus
  \def\vdots{\mathord{\vbox{\baselineskip3.2pt\lineskiplimit0pt\kern4pt\hbox{.}\hbox{.}\hbox{.}}}}%
  \def\ddots{\mathinner{\mkern1mu\raise6.5pt\hbox{.}\mkern2mu\raise3.6pt\hbox{.}\mkern2mu\raise0.7pt\hbox{.}\mkern1mu}}%
  \def\triangle{\mathord{\scalebox{0.9}{$\symup{\Delta}$}}}%
  \def\nabla{\mathord{\raisebox{\depth}{\scalebox{1}[-1]{$\symup{\Delta}$}}}}%
  \def\checkmark{\popcheck{popdark}}%
  \def\star{\mathbin{\ast}}\def\bigstar{\mathord{\ast}}%
  \def\diamond{\mathbin{\scalebox{0.6}{$\Diamond$}}}%
  \def\bigcup{\mathop{\mathchoice{\raisebox{-0.3em}{\scalebox{1.7}{$\cup$}}}{\scalebox{1.25}{$\cup$}}{\cup}{\cup}}\displaylimits}%
  \def\bigcap{\mathop{\mathchoice{\raisebox{-0.3em}{\scalebox{1.7}{$\cap$}}}{\scalebox{1.25}{$\cap$}}{\cap}{\cap}}\displaylimits}%
  \def\lesseqqgtr{\mathrel{\lessgtr}}\def\bigoplus{\mathop{\oplus}\displaylimits}%
}
\providecommand{\ssi}{si et seulement si} % abréviation fréquente
\AtBeginDocument{\providecommand{\wideparen}{\overparen}} % arc de cercle

%% FIGFIX-BEGIN v5 — correctifs globaux des figures (docs/onbuch-generation/tools/figfix)
\usepackage{adjustbox}\usepackage{etoolbox}\tcbuselibrary{raster}
% toute figure TikZ/pgfplots plus large que la ligne est réduite à la largeur disponible
\BeforeBeginEnvironment{tikzpicture}{\begin{adjustbox}{max width=\linewidth}}
\AfterEndEnvironment{tikzpicture}{\end{adjustbox}}
% légendes pgfplots lisibles par-dessus les courbes
\pgfplotsset{every axis/.append style={legend style={fill=white,fill opacity=0.92,text opacity=1,draw=black!15,inner sep=3pt}}}
% étiquettes d'arêtes (syntaxe edge["..."]) avec fond blanc
\tikzset{every edge quotes/.append style={fill=white,inner sep=1.2pt}}
%% FIGFIX-END
\begin{document}

\pagegarde{Les échanges inter-régionaux au Cameroun}{Géographie}{Tle A}{Module 1 : Le Cameroun}{N09}{2 h + TP 4 (2 h)}{Cette leçon étudie les échanges inter-régionaux au Cameroun : leurs facteurs de développement (infrastructures, interdépendance des régions), l'organisation du commerce intérieur (marchés périodiques, bayam selam, vendeurs à la sauvette, centres commerciaux) et ses difficultés. Elle s'achève par la construction d'une carte des flux commerciaux inter-régionaux (TP 4).
\vspace{4pt}
\begin{multicols}{2}\small
\begin{itemize}
\item À la fin de cette leçon, je sais définir les notions d'échanges inter-régionaux, de commerce intérieur et d'intégration nationale
\item À la fin de cette leçon, je sais identifier et expliquer les facteurs de développement des échanges inter-régionaux au Cameroun
\item À la fin de cette leçon, je sais évaluer l'état des infrastructures de transport (routes, chemin de fer, aéroports, ports) et leur rôle dans les échanges
\item À la fin de cette leçon, je sais montrer l'interdépendance entre les régions du Cameroun à partir d'exemples de produits échangés
\item À la fin de cette leçon, je sais décrire l'organisation du commerce intérieur : marchés périodiques, bayam selam, vendeurs à la sauvette, centres commerciaux
\item À la fin de cette leçon, je sais inventorier et expliquer les difficultés du commerce intérieur camerounais
\end{itemize}\end{multicols}}

\begin{sommaire}
\begin{multicols}{2}
\begin{enumerate}
\item Introduction : notions clés et situation-problème
\item Les facteurs de développement des échanges : les infrastructures
\item L'interdépendance entre les régions
\item L'organisation du commerce intérieur
\item Les difficultés du commerce intérieur
\item TP 4 : Carte des flux commerciaux inter-régionaux au Cameroun
\item Activité d'intégration
\item Méthodes et exercices résolus
\item Exercices
\item Corrigés des exercices
\item Fiche bilan
\end{enumerate}
\end{multicols}
\end{sommaire}

\begin{objectifs}
\begin{itemize}
\item À la fin de cette leçon, je sais définir les notions d'échanges inter-régionaux, de commerce intérieur et d'intégration nationale
\item À la fin de cette leçon, je sais identifier et expliquer les facteurs de développement des échanges inter-régionaux au Cameroun
\item À la fin de cette leçon, je sais évaluer l'état des infrastructures de transport (routes, chemin de fer, aéroports, ports) et leur rôle dans les échanges
\item À la fin de cette leçon, je sais montrer l'interdépendance entre les régions du Cameroun à partir d'exemples de produits échangés
\item À la fin de cette leçon, je sais décrire l'organisation du commerce intérieur : marchés périodiques, bayam selam, vendeurs à la sauvette, centres commerciaux
\item À la fin de cette leçon, je sais inventorier et expliquer les difficultés du commerce intérieur camerounais
\item À la fin de cette leçon, je sais observer, localiser et décrire des phénènes géographiques à partir de cartes, photographies et tableaux statistiques
\item À la fin de cette leçon, je sais construire et légender une carte des flux commerciaux inter-régionaux (produits, zones de production, zones de commercialisation et de consommation)
\item À la fin de cette leçon, je sais construire un croquis ou un diagramme représentant les flux commerciaux entre régions
\end{itemize}
\end{objectifs}

\begin{prerequis}
\begin{itemize}
\item La diversité des milieux physiques et des productions agricoles du Cameroun (classe de Première)
\item Les régions et les grandes villes du Cameroun (localisation)
\item Les notions de population, d'urbanisation et de consommation urbaine
\item Les infrastructures de transport du Cameroun vues en début d'année
\item Les techniques de lecture et de légendage d'une carte (acquis de Seconde et Première)
\end{itemize}
\end{prerequis}

\newpage

\coursec{Introduction : notions clés et situation-problème}

\courssub{Une situation d'accroche : le panier de la ménagère au marché Mokolo}

Il est 7 heures du matin au marché Mokolo, l'un des plus grands marchés de Yaoundé. Une ménagère remplit son panier : des pommes de terre venues de l'Ouest, du poisson fumé pêché dans le Sud, des oignons séchés de l'Extrême-Nord, du riz étuvé du Nord. En quelques minutes, elle a réuni dans son cabas les produits de quatre régions du Cameroun, dont certaines se trouvent à plus de \SI{800}{km} de la capitale. Rien d'extraordinaire, semble-t-il : c'est le quotidien de millions de Camerounais. Et pourtant, ce simple achat pose une véritable question de géographie.

Comment ces produits parcourent-ils des centaines de kilomètres malgré des routes dégradées et un chemin de fer vétuste ? Qui sont ces commerçants, les fameux \cle{bayam selam}, qui sillonnent le pays pour approvisionner nos villes ? Pourquoi le commerce intérieur camerounais reste-t-il à la fois dynamique et fragile ? Cette leçon explore les échanges inter-régionaux, moteur discret mais essentiel de l'intégration nationale.

\begin{attention}[La problématique de la leçon]
Toute la leçon répond à une question centrale qu'il faut retenir dès maintenant : \emph{comment s'organisent les échanges inter-régionaux au Cameroun, quels facteurs les favorisent et quelles difficultés les freinent ?} Au Baccalauréat, une problématique clairement formulée est la colonne vertébrale de toute composition de géographie : apprends à la repérer et à la mobiliser.
\end{attention}

\courssub{Les notions clés : échanges inter-régionaux et commerce intérieur}

\textbf{De l'intuition à la définition.} Avant de donner des définitions rigoureuses, partons d'une observation simple. Le Cameroun compte dix régions aux milieux très différents : l'Ouest, aux sols volcaniques fertiles, produit massivement des pommes de terre, des choux et des carottes ; l'Extrême-Nord, de climat soudanien, est le pays de l'oignon, du niébé et du coton ; le Littoral fournit le poisson de mer, le plantain et le manioc ; le Nord et l'Adamaoua fournissent le bétail et le riz. Aucune région ne peut donc vivre en autarcie : chacune a besoin des produits des autres. Il en résulte des \cle{flux} permanents de marchandises, mais aussi de personnes (commerçants, transporteurs, consommateurs) et de capitaux (l'argent qui circule en sens inverse des marchandises). C'est exactement cela, les échanges inter-régionaux.

\begin{definition}[Les échanges inter-régionaux]
Les \cle{échanges inter-régionaux} désignent l'ensemble des mouvements de biens (marchandises), de personnes et de capitaux entre les régions d'un même pays. Au Cameroun, ils relient les zones de production (par exemple les hautes terres de l'Ouest pour les pommes de terre) aux zones de commercialisation et de consommation (par exemple les marchés de Yaoundé et de Douala).
\end{definition}

Ces échanges prennent une forme précise : le commerce. Lorsqu'un bayam selam achète un sac de pommes de terre à Bafoussam et le revend à Yaoundé, il réalise une opération commerciale à l'intérieur du pays. D'où la deuxième notion clé.

\begin{definition}[Le commerce intérieur]
Le \cle{commerce intérieur} est l'ensemble des opérations d'achat et de vente de biens réalisées à l'intérieur des frontières d'un pays. Il se distingue du commerce extérieur (importations et exportations), qui franchit les frontières nationales. Le commerce intérieur comprend le commerce de gros (achat en grande quantité aux producteurs), le commerce de demi-gros et le commerce de détail (vente aux consommateurs sur les marchés, dans les boutiques ou à la sauvette).
\end{definition}

\begin{attention}[Piège fréquent : confondre commerce intérieur et commerce extérieur]
Beaucoup d'élèves citent le port de Douala ou les exportations de pétrole et de cacao dans une leçon sur le commerce \emph{intérieur}. Erreur ! Le commerce intérieur s'arrête aux frontières nationales : le sac d'oignons de Maroua vendu à Yaoundé relève du commerce intérieur ; le même sac vendu à N'Djamena (Tchad) relève du commerce extérieur. Vérifie toujours si la marchandise franchit ou non une frontière.
\end{attention}

\courssub{Les échanges, ciment de l'intégration nationale}

Cette leçon s'inscrit dans la famille de situations \emph{l'intégration nationale}. Pourquoi ? Parce que les échanges inter-régionaux font bien plus que déplacer des marchandises : ils tissent des liens entre les régions et entre les Camerounais.

\begin{itemize}
\item \textbf{Sur le plan économique}, les échanges créent l'\cle{interdépendance} : l'Ouest a besoin d'écouler ses pommes de terre, Yaoundé a besoin de se nourrir, l'Extrême-Nord a besoin de vendre ses oignons. Chaque région devient utile aux autres.
\item \textbf{Sur le plan social et culturel}, les bayam selam de l'Ouest dorment à Maroua, les éleveurs du Nord descendent vendre leurs bœufs à Yaoundé, les pêcheurs du Sud alimentent le Centre. Les langues, les cuisines et les habitudes circulent avec les marchandises : on apprend à se connaître.
\item \textbf{Sur le plan politique}, un pays dont les régions commercent entre elles est un pays plus uni : le sentiment d'appartenir à une même nation se renforce quand le quotidien de chacun dépend du travail des autres.
\end{itemize}

\begin{aretenir}[L'essentiel du lien avec l'intégration nationale]
Les échanges inter-régionaux sont un instrument d'\cle{intégration nationale} : en rendant les régions économiquement interdépendantes et en brassant les populations, ils consolident l'unité du pays. Un réseau commercial intérieur dynamique est donc un enjeu de développement autant que de cohésion nationale.
\end{aretenir}

\courssub{Analyse de la situation-problème : le ravitaillement du marché Mokolo}

Revenons au marché Mokolo et analysons-le avec la méthode du géographe. Le document de départ est une simple observation de terrain : l'origine des produits vendus sur les étals.

\begin{methode}[Analyser une situation-problème en géographie]
\begin{enumerate}
\item \textbf{Identifier l'espace concerné} : ici, le marché Mokolo à Yaoundé (région du Centre), relié aux régions de production (Ouest, Sud, Nord, Extrême-Nord).
\item \textbf{Identifier les acteurs} : producteurs (agriculteurs de l'Ouest, pêcheurs du Sud, éleveurs du Nord), commerçants (bayam selam, grossistes, détaillants), transporteurs (camionneurs, conducteurs de bus et de pick-up), consommateurs (ménagères, restaurateurs).
\item \textbf{Repérer les flux} : sens des marchandises (des campagnes vers la ville), sens de l'argent (de la ville vers les campagnes), itinéraires empruntés (routes nationales N1, N4, N6, chemin de fer).
\item \textbf{Dégager les enjeux} : approvisionnement alimentaire des villes, revenus des producteurs et des commerçants, coût du transport répercuté sur les prix, dépendance aux infrastructures.
\end{enumerate}
\end{methode}

Appliquée à Mokolo, cette méthode révèle un système complet. La pomme de terre, par exemple, suit un circuit bien rôdé : récoltée sur les hautes terres de l'Ouest (Bafoussam, Dschang, Mbouda, Santa dans le Nord-Ouest), elle est achetée au champ ou au marché rural par des bayam selam, chargée dans des camions qui empruntent la nationale N4, déchargée au marché Mokolo où des grossistes et demi-grossistes la répartissent, puis vendue au détail à la ménagère. Chaque intermédiaire ajoute sa marge : c'est pourquoi le kilo de pommes de terre coûte nettement plus cher à Yaoundé qu'à Bafoussam.

\begin{popfigure}
\begin{center}
\begin{tikzpicture}[font=\small, >=Stealth, node distance=0.9cm]
\tikzstyle{etape}=[draw=popblue, fill=popblueL, rounded corners=3pt, minimum height=0.9cm, inner sep=5pt, align=center]
\tikzstyle{lieu}=[draw=popdark, fill=popcream, rounded corners=3pt, minimum height=0.9cm, inner sep=5pt, align=center]
\node[lieu, align=center] (prod){\textbf{Zone de production}\\Hautes terres de l'Ouest\\(Bafoussam, Dschang, Mbouda)};
\node[etape, right=1.1cm of prod, align=center] (achat){\textbf{Achat en gros}\\Bayam selam\\au champ ou au marché rural};
\node[etape, right=1.1cm of achat, align=center] (transp){\textbf{Transport}\\Camions sur la N4\\(environ 300 km)};
\node[etape, right=1.1cm of transp, align=center] (gros){\textbf{Marché Mokolo}\\Grossistes et\\demi-grossistes};
\node[lieu, right=1.1cm of gros, align=center] (conso){\textbf{Zone de consommation}\\Détaillants, ménagères\\(Yaoundé)};
\draw[->, very thick, poporange] (prod) -- (achat);
\draw[->, very thick, poporange] (achat) -- (transp);
\draw[->, very thick, poporange] (transp) -- (gros);
\draw[->, very thick, poporange] (gros) -- (conso);
\draw[->, very thick, popgreen, dashed] (conso.south) to[bend right=18] node[below, align=center]{\small Flux de capitaux (paiements)} (prod.south);
\end{tikzpicture}
\end{center}
\legende{Schéma 1 : le circuit de la pomme de terre, de la zone de production (Ouest) à la zone de consommation (Yaoundé). Les flèches orange représentent le flux des marchandises ; la flèche verte en pointillés représente le flux inverse des capitaux. À chaque étape, un intermédiaire prélève une marge, ce qui renchérit le prix final.}
\end{popfigure}

\begin{savaistu}[Des vivres qui voyagent très loin]
Une grande partie des vivres consommées à Douala et Yaoundé provient d'autres régions, parfois situées à plus de \SI{500}{km} : les oignons et le niébé descendent de l'Extrême-Nord (plus de \SI{1000}{km} par la route depuis Maroua), le bétail arrive de l'Adamaoua et du Nord, le macabo et le plantain remontent du Littoral et du Sud-Ouest. Les deux métropoles ne produisent qu'une faible part de ce qu'elles mangent : sans les échanges inter-régionaux, elles ne pourraient tout simplement pas se nourrir.
\end{savaistu}

\courssub{De la situation-problème à la problématique}

L'observation du marché Mokolo fait apparaître une contradiction qui est le cœur de la leçon : les échanges inter-régionaux sont \textbf{intenses} (les villes sont ravitaillées chaque jour, les bayam selam ne comptent plus les kilomètres), mais ils se déploient sur des \textbf{infrastructures fragiles} (routes dégradées, chemin de fer vétuste, ports mal connectés à l'intérieur). D'où la problématique déjà formulée : comment s'organisent les échanges inter-régionaux au Cameroun, quels facteurs les favorisent et quelles difficultés les freinent ?

\begin{exoresolu}[Identifier les éléments d'un flux commercial]
\textbf{Énoncé.} Un bayam selam achète à Maga (Extrême-Nord) 50 sacs d'oignons à \num{12000} FCFA le sac. Il paie \num{150000} FCFA de transport par camion jusqu'au marché Mokolo (Yaoundé), puis revend chaque sac \num{22000} FCFA. (1) Identifie la zone de production, la zone de consommation, les acteurs et le flux. (2) Calcule le bénéfice brut du bayam selam.
\tcblower
\textbf{Solution détaillée.}
(1) La zone de production est l'Extrême-Nord (Maga, plaine du Logone) ; la zone de consommation est Yaoundé (marché Mokolo). Les acteurs sont le producteur, le bayam selam (commerçant), le transporteur (camionneur) et les consommateurs. Le flux de marchandises va du Nord vers le Centre ; le flux de capitaux va en sens inverse.

(2) Coût total d'achat : $50 \times \num{12000} = \num{600000}$ FCFA. En ajoutant le transport :
\[ \text{Coût total} = \num{600000} + \num{150000} = \num{750000} \text{ FCFA}. \]
Recette de la vente : $50 \times \num{22000} = \num{1100000}$ FCFA. Bénéfice brut :
\[ \text{Bénéfice} = \num{1100000} - \num{750000} = \num{350000} \text{ FCFA}. \]
Ce calcul montre que le commerce intérieur peut être rentable, ce qui explique son dynamisme ; mais ce bénéfice fond vite si le camion tombe en panne ou si la route est coupée, d'où la fragilité du système.
\end{exoresolu}

\courssub{Annonce du plan}

Pour répondre à la problématique, la leçon suivra trois étapes logiques, auxquelles s'ajoute un travail pratique de cartographie :

\begin{enumerate}
\item \textbf{Les facteurs de développement des échanges} : le rôle des infrastructures (routes, chemin de fer, aéroports, ports) et l'interdépendance entre les régions, qui rend les échanges nécessaires.
\item \textbf{L'organisation du commerce intérieur} : les marchés périodiques, les bayam selam, les vendeurs à la sauvette et les centres commerciaux, qui structurent la distribution des marchandises.
\item \textbf{Les difficultés du commerce intérieur} : les obstacles qui freinent les flux et renchérissent les prix.
\item \textbf{TP 4} : la construction d'une carte des flux commerciaux inter-régionaux au Cameroun, pour passer de la description à la représentation géographique.
\end{enumerate}

\begin{exercice}[Pour vérifier ta compréhension]
1. Définis en deux phrases les échanges inter-régionaux et le commerce intérieur, en soulignant leur différence. 2. Pourquoi dit-on que les échanges inter-régionaux sont un facteur d'intégration nationale ? Donne deux arguments. 3. En t'inspirant du schéma 1, décris le circuit probable du poisson fumé du Sud vers Yaoundé (zone de production, acteurs, flux).
\end{exercice}

\coursec{Les facteurs de développement des échanges : les infrastructures}

Dans l'introduction de cette leçon, nous avons suivi le parcours d'une caisse de tomates partie du marché de Foumbot (Ouest) pour aboutir sur les étals du marché Mokolo à Yaoundé. Nous avons constaté que ce voyage de quelques centaines de kilomètres conditionne le prix final du produit, sa fraîcheur et parfois sa simple disponibilité. La question se pose alors : qu'est-ce qui permet, ou empêche, ces échanges entre les régions du Cameroun ? La réponse tient en grande partie à un mot : les \cle{infrastructures de transport}. C'est à leur étude que cette section est consacrée.

\courssub{Qu'est-ce qu'une infrastructure de transport ?}

\begin{definition}[Infrastructure de transport]
Une \cle{infrastructure de transport} est l'ensemble des équipements fixes — routes, voies ferrées, ponts, ports, aéroports, gares — qui permettent la circulation des personnes et des marchandises entre les lieux de production et les lieux de consommation. On distingue l'infrastructure (la route elle-même) du moyen de transport (le camion, le train, le bateau) qui l'utilise.
\end{definition}

L'intuition est simple : deux régions complémentaires ne peuvent échanger que si elles sont \emph{reliées}. Le Nord produit du coton et des oignons, le Littoral abrite les usines et le port, l'Ouest fournit les vivres : sans routes praticables, cette complémentarité reste théorique. Les infrastructures sont donc la condition matérielle de l'\cle{intégration nationale}. Mais attention : ce n'est pas seulement leur \emph{existence} qui compte, c'est aussi leur \emph{état}, leur \emph{entretien} et leur \emph{interconnexion}.

\courssub{Le réseau routier : en progression, mais inégalement entretenu}

La route est de loin le premier mode de transport au Cameroun : elle assure l'essentiel du trafic des voyageurs et des marchandises. Le réseau routier total est estimé à environ \num{120000} km (ordre de grandeur), toutes catégories confondues, dont seulement une fraction — de l'ordre de 10 \% — est bitumée.

\begin{exemplebox}[Les grands axes bitumés]
L'axe structurant du pays est la route \cle{Douala--Yaoundé--Ngaoundéré}. La \textbf{Nationale 1} relie Yaoundé à Ngaoundéré (environ 850 km) via Bafia, Bafoussam, Foumban, Nkambé et Meiganga, et se prolonge vers Garoua, Maroua et Kousseri à la frontière tchadienne. La \textbf{Nationale 3} relie Douala à Yaoundé (environ 250 km) via Edéa. La \textbf{Ring Road} (route circulaire du Nord-Ouest, environ 360 km, reliant Bamenda, Ndop, Nkambe, Kumbo et Fundong) est en construction et en réhabilitation depuis des décennies. D'autres axes majeurs ont été bitumés récemment : Yaoundé--Bertoua (N10), Bertoua--Garoua-Boulaï (N15), et l'autoroute Yaoundé--Nsimalen.
\end{exemplebox}

Le constat géographique est double :

\begin{itemize}
\item \textbf{Des progrès réels} : le kilométrage de routes bitumées augmente régulièrement (autoroute Yaoundé--Nsimalen, contournement de Douala, bitumage d'axes vers l'Est et l'Extrême-Nord), ce qui réduit les temps de parcours et les coûts de transport.
\item \textbf{Des limites fortes} : l'entretien est irrégulier, si bien que de nombreux tronçons bitumés se dégradent vite (nids-de-poule, chaussées effondrées). Surtout, la majorité du réseau est constituée de \cle{routes en terre} et de \cle{pistes rurales} impraticables en saison des pluies. Des bassins de production entiers (vallées de l'Ouest, départements de l'Est et de l'Adamaoua) se retrouvent alors coupés des marchés : les tomates pourrissent au champ pendant que les prix flambent en ville.
\end{itemize}

\begin{attention}[Erreur fréquente au Bac]
N'écris jamais que « le Cameroun n'a pas d'infrastructures » ou « manque totalement de routes ». C'est faux : le pays dispose d'un réseau routier, ferroviaire, portuaire et aéroportuaire réel et ancien. Le problème est ailleurs : la \textbf{vétusté} des équipements, l'\textbf{entretien irrégulier} et la \textbf{mauvaise interconnexion} entre les modes (le port n'est pas bien relié à l'intérieur, le rail ne dessert pas tout le pays). Une copie qui nuance ainsi montre une vraie compétence d'analyse.
\end{attention}

\courssub{Le chemin de fer : un réseau vétuste mais stratégique}

\begin{exemplebox}[Le Transcamerounais en chiffres]
Le réseau ferré camerounais compte environ \num{1000} km de voies à écartement métrique. Son axe principal, surnommé le \cle{Transcamerounais}, relie \textbf{Douala à Ngaoundéré} sur environ 620 km (via Yaoundé, Eseka, Belabo). Une seconde ligne, plus courte, relie Douala à Mbanga et Kumba (ligne dite « Transcamerounais Ouest », héritée de l'époque coloniale allemande). L'exploitation est assurée par \textbf{Camrail} (concession depuis 1999).
\end{exemplebox}

Le chemin de fer camerounais souffre de trois maux :

\begin{enumerate}
\item \textbf{La vétusté} : voies et matériel roulant datent pour beaucoup de plusieurs décennies ; la vitesse commerciale est faible (souvent moins de 40 km/h en moyenne).
\item \textbf{L'insécurité} : les déraillements sont fréquents — la catastrophe d'Eséka (octobre 2016, environ 80 morts) a tragiquement illustré l'état du réseau.
\item \textbf{La lenteur} : le trajet Yaoundé--Ngaoundéré peut durer plus de 15 heures, ce qui décourage les voyageurs.
\end{enumerate}

Pourtant, le rail garde un \textbf{rôle clé pour les marchandises lourdes} : le bois de l'Est et du Sud, le coton du Nord, le ciment, les produits pétroliers et les conteneurs à destination de Ngaoundéré (relais vers le Tchad) empruntent massivement le train, car le coût au tonne-kilomètre y est bien inférieur à celui du camion. Le rail soulage aussi la route nationale 1, saturée de poids lourds.

\courssub{Le transport aérien : des aéroports, mais une compagnie nationale fragilisée}

Le Cameroun possède plusieurs aéroports, dont trois internationaux : \textbf{Douala} (le plus fréquenté), \textbf{Yaoundé--Nsimalen} et \textbf{Garoua} ; \textbf{Maroua}, Bamenda, Bafoussam ou Bertoua disposent de pistes d'importance variable. L'avion est précieux pour les liaisons rapides entre les métropoles du Sud et les villes du Grand Nord, éloignées de plus de \num{1000} km par la route.

Mais la compagnie nationale, \textbf{Camair-Co} (créée en 2006 après la faillite de Cameroon Airlines), connaît des \cle{difficultés financières et opérationnelles} chroniques : flotte réduite, dettes, dessertes intérieures irrégulières. Résultat : les liaisons aériennes intérieures restent chères et peu fiables, et l'avion joue un rôle marginal dans le commerce intérieur des marchandises, sauf pour les produits périssables à forte valeur.

\courssub{Les ports : des façades maritimes mal connectées à l'intérieur}

Le Cameroun dispose de trois ports principaux :

\begin{itemize}
\item Le \textbf{port de Douala} (estuaire du Wouri) : premier port du pays, il traite l'essentiel du commerce extérieur (conteneurs, hydrocarbures, produits divers) et alimente le commerce intérieur en produits importés redistribués vers toutes les régions.
\item Le \textbf{port de Kribi} : port en \cle{eau profonde} inauguré en 2018, spécialisé dans les conteneurs et les grands navires ; il est aussi le débouché du pipeline tchado-camerounais.
\item Le \textbf{port de Limbe} : port secondaire, lié à la SONARA (raffinerie) et aux exportations de la zone côtière.
\end{itemize}

\begin{savaistu}[Douala, port du Cameroun... et de deux pays enclavés]
Le port de Douala ne sert pas que le Cameroun : il traite aussi l'essentiel du commerce extérieur du \textbf{Tchad} et de la \textbf{Centrafrique}, deux États sans littoral dont les marchandises transitent par le corridor Douala--Ngaoundéré (rail + route) puis par la route vers N'Djamena et Bangui. Paradoxe : ce port international reste \emph{mal connecté à son propre arrière-pays} — l'Adamaoua, le Nord et l'Est camerounais souffrent d'un enclavement qui renchérit tous les produits qui y arrivent.
\end{savaistu}

Le problème majeur est donc la \cle{connexion port--intérieur} : une seule ligne ferroviaire, une route nationale 1 saturée, et l'Est (Bertoua, Batouri, Yokadouma) qui ne dispose d'aucune voie ferrée. Les coûts de transport explosent avec la distance : un conteneur débarqué à Douala coûte beaucoup plus cher rendu à Maroua qu'à Yaoundé, ce qui creuse les inégalités de prix entre régions.

\courssub{Cartographier le réseau de transport du Cameroun}

\begin{methode}[Lire une carte de réseau de transport]
\begin{enumerate}
\item \textbf{Identifier les nœuds} : repère les villes principales (Douala, Yaoundé, Ngaoundéré, Garoua, Maroua, Bamenda, Bafoussam, Bertoua) — ce sont les points de concentration et de redistribution des flux.
\item \textbf{Identifier les axes} : distingue les routes bitumées (traits pleins), les voies ferrées (traits particuliers), les ports et aéroports (symboles).
\item \textbf{Repérer les zones bien desservies} : le triangle Douala--Yaoundé--Bafoussam concentre routes bitumées, rail et aéroports.
\item \textbf{Repérer les zones enclavées} : l'Est, l'Adamaoua et le Nord n'ont qu'un ou deux axes ; le Sud-Est en est presque dépourvu.
\item \textbf{Conclure} : relie la densité du réseau à l'intensité des échanges — là où les axes manquent, les produits sont rares et chers.
\end{enumerate}
\end{methode}

\begin{popfigure}
\begin{center}
\includegraphics[width=\linewidth,height=9.5cm,keepaspectratio]{N01/cmr_map.jpg}
\end{center}
\legende{Carte du réseau de transport du Cameroun (relief ombré) : routes principales (en rouge), voie ferrée Transcamerounais (Douala--Yaoundé--Ngaoundéré), capitales régionales. On observe la concentration des axes dans le triangle Douala--Yaoundé--Bafoussam, et l'éloignement des grands axes dans l'Est, l'Adamaoua et le Grand Nord. Ports (Douala, Kribi, Limbé) et aéroports sont à situer par l'élève sur ce fond. \\ Source : Wikimedia Commons, CIA, domaine public.}
\end{popfigure}

\courssub{Étude de documents : lire les infrastructures sur le terrain}

\begin{experience}[Analyse de deux photographies (démarche guidée)]
\textbf{Document 1 : photographie d'une route rurale dégradée en saison des pluies (Ouest Cameroun).} On observe une piste en terre battue transformée en bourbier ; un camion chargé de bananes plantain est embourbé ; des riverains déchargent la marchandise à dos d'homme pour alléger le véhicule.
\begin{itemize}
\item \emph{Observation} : la piste existe mais n'est pas praticable toute l'année.
\item \emph{Interprétation} : l'absence de bitumage et d'entretien interrompt les échanges ; les produits périssables risquent de pourrir, les coûts de transport augmentent (temps, main-d'œuvre, avaries).
\item \emph{Conséquence géographique} : les prix montent en ville, les revenus des producteurs baissent — l'enclavement pénalise les deux bouts de la chaîne.
\end{itemize}
\textbf{Document 2 : photographie du port de Douala.} On observe des grues à conteneurs, des navires à quai, des files de camions attendant à la sortie, et en arrière-plan l'embouteillage permanent du pont sur le Wouri.
\begin{itemize}
\item \emph{Observation} : le port est moderne et actif, mais congestionné.
\item \emph{Interprétation} : la faiblesse n'est pas le port lui-même mais sa \emph{connexion} à l'arrière-pays : une seule voie ferrée, des routes saturées.
\item \emph{Conséquence géographique} : le délai et le coût d'acheminement vers Ngaoundéré, Garoua ou Bertoua renchérissent tous les produits importés dans ces régions.
\end{itemize}
\end{experience}

\courssub{Des chiffres pour mesurer le retard routier}

Comparer les longueurs de routes bitumées et non bitumées permet de visualiser l'ampleur du défi. Les ordres de grandeur ci-dessous (réseau total d'environ \num{120000} km, dont environ 10 \% bitumés) sont indicatifs : les sources officielles varient selon les années et la prise en compte ou non des pistes rurales.

\begin{popfigure}
\begin{center}
\begin{tikzpicture}
\begin{axis}[
  ybar, bar width=1.6cm,
  width=0.85\linewidth, height=6.5cm,
  ymin=0, ymax=125000,
  ylabel={Longueur (km)},
  symbolic x coords={Routes bitumées, Routes non bitumées},
  xtick=data,
  nodes near coords,
  every node near coord/.append style={font=\small\bfseries},
  ymajorgrids=true, grid style={popline!40},
  axis line style={popdark},
]
\addplot[fill=popblue, draw=popdark] coordinates {(Routes bitumées,10000) (Routes non bitumées,110000)};
\end{axis}
\end{tikzpicture}
\end{center}
\legende{Longueurs comparées des routes bitumées et non bitumées au Cameroun (ordres de grandeur, réseau total environ \num{120000} km). Moins d'un kilomètre sur dix est bitumé : c'est toute la difficulté du commerce intérieur, surtout en saison des pluies.}
\end{popfigure}

\courssub{Compléments utiles au Bac : pistes rurales et nouvelles technologies}

\textbf{Les pistes rurales, maillon invisible mais décisif.} Entre le champ et le marché hebdomadaire, ce ne sont ni le rail ni le goudron qui circulent, mais la piste : motos-taxis (\emph{bendskins}), vélos, portage à dos. Les programmes de \cle{désenclavement des bassins de production} (réhabilitation de pistes, ponts ruraux) ont un effet direct sur les échanges : là où une piste est rendue praticable, les \cle{bayam selam} reviennent, les prix payés aux producteurs montent, et l'approvisionnement des villes s'améliore.

\textbf{La téléphonie mobile et le mobile money, infrastructures immatérielles.} Les échanges ne dépendent pas que des routes : ils dépendent aussi de l'\emph{information} et du \emph{paiement}. Avec un taux de pénétration du téléphone mobile très élevé (plusieurs dizaines de millions d'abonnements pour environ 28 millions d'habitants, ordre de grandeur), un commerçant de Maroua peut connaître le prix des oignons à Douala avant d'expédier, et un acheteur peut payer par \cle{mobile money} (MTN Mobile Money, Orange Money) sans transporter d'espèces sur des routes peu sûres. Ces infrastructures immatérielles réduisent les coûts de transaction et sécurisent le commerce intérieur.

\begin{exoresolu}[Expliquer le rôle des infrastructures dans les échanges inter-régionaux]
\textbf{Énoncé.} À partir de la carte du réseau de transport et de tes connaissances, montre en une quinzaine de lignes comment les infrastructures de transport favorisent ou freinent les échanges inter-régionaux au Cameroun.
\tcblower
\textbf{Solution détaillée.} Les infrastructures de transport sont la condition matérielle des échanges inter-régionaux : elles relient les zones de production aux zones de consommation. Au Cameroun, leur rôle est ambivalent. D'un côté, elles favorisent les échanges : le bitumage croissant des routes (axe Douala--Yaoundé--Ngaoundéré, Ring Road en construction) réduit les temps de parcours ; le Transcamerounais (environ 620 km entre Douala et Ngaoundéré) transporte à faible coût les marchandises lourdes (bois, coton, ciment) ; les ports de Douala et de Kribi alimentent tout le pays en produits importés ; les aéroports relient rapidement le Grand Nord aux métropoles du Sud. De l'autre côté, elles freinent les échanges : l'entretien routier est irrégulier et la plupart des pistes rurales sont impraticables en saison des pluies, ce qui enclave les bassins de production ; le chemin de fer est vétuste, lent et sujet aux déraillements ; la compagnie Camair-Co est en difficultés ; enfin, les ports sont mal connectés à l'intérieur, si bien que l'Adamaoua, le Nord et l'Est paient leurs produits beaucoup plus cher. En définitive, ce n'est pas l'absence d'infrastructures qui pose problème, mais leur vétusté, leur entretien irrégulier et leur mauvaise interconnexion : moderniser et interconnecter le réseau est donc un enjeu majeur de l'intégration nationale.
\end{exoresolu}

\begin{aretenir}[L'essentiel de la section]
\begin{itemize}
\item Les infrastructures de transport (routes, rail, ports, aéroports) sont la condition matérielle des échanges inter-régionaux.
\item Le réseau routier s'étend et se bitume, mais l'entretien est irrégulier et les pistes rurales sont impraticables en saison des pluies.
\item Le Transcamerounais (environ 620 km, Douala--Ngaoundéré) est vétuste et lent, mais reste stratégique pour les marchandises lourdes.
\item Les aéroports sont nombreux, mais Camair-Co est fragilisée ; les ports (Douala, Kribi, Limbe) sont mal connectés à l'intérieur.
\item Le vrai problème n'est pas l'absence d'infrastructures, mais leur vétusté, leur entretien et leur interconnexion — auxquels s'ajoutent désormais les infrastructures immatérielles (téléphonie mobile, mobile money).
\end{itemize}
\end{aretenir}

Ces infrastructures, qu'elles soient efficaces ou défaillantes, dessinent la géographie des échanges. Mais pourquoi les régions camerounaises ont-elles tant besoin d'échanger entre elles ? C'est ce que nous verrons dans la section suivante, consacrée à l'interdépendance entre les régions.

\coursec{L'interdépendance entre les régions}

Dans la section précédente, nous avons vu que les infrastructures de transport — routes bitumées, chemin de fer, ports, aéroports — constituent le \cle{support matériel} des échanges. Mais une route ne justifie son existence que si des marchandises ont une raison de circuler sur elle. Pourquoi un camion quitte-t-il Bafoussam à 3 heures du matin chargé de pommes de terre vers Douala ? Pourquoi des bœufs de l'Adamaoua descendent-ils vers Yaoundé ? La réponse tient en un mot : l'\cle{interdépendance}. C'est le moteur profond du commerce intérieur camerounais, et c'est elle que nous étudions maintenant.

\courssub{Le principe : aucune région n'est autosuffisante}

\begin{definition}[Interdépendance régionale]
L'\cle{interdépendance régionale} est la situation dans laquelle les régions d'un pays dépendent les unes des autres pour leur \textbf{approvisionnement} (elles achètent ce qu'elles ne produisent pas) et pour l'\textbf{écoulement} de leurs produits (elles vendent leur surplus aux autres régions). Elle résulte de la \textbf{complémentarité} des milieux naturels, des spécialisations de production et de la répartition inégale de la population et des activités.
\end{definition}

L'intuition est simple, et tu peux la vérifier sur n'importe quel étal de marché. Le Cameroun s'étire sur plus de \SI{1200}{km} du sud au nord, de la forêt équatoriale humide du Golfe de Guinée jusqu'à la savane sèche du bassin du lac Tchad. Ce gradient climatique exceptionnel — qui fait dire qu'il est « l'Afrique en miniature » — signifie que \textbf{chaque région produit ce que son milieu permet} : le cacao exige la forêt humide du Sud, l'oignon et le niébé prospèrent dans la savane sèche de l'Extrême-Nord, la pomme de terre réclame les hautes terres fraîches de l'Ouest, le zébu paît dans les pâturages de l'Adamaoua. Inversement, aucune région ne peut tout produire : Douala ne cultive pas de cacao dans ses rues, Maroua ne pêche pas de capitaine dans la Mayo Kani, et l'Extrême-Nord importe l'huile de palme qu'il ne peut pas faire pousser.

\begin{aretenir}[La logique de l'interdépendance]
\begin{itemize}
\item Les \textbf{milieux naturels} diffèrent (climat, sols, relief) $\Rightarrow$ les \textbf{productions} diffèrent.
\item La \textbf{population} et les \textbf{activités industrielles} sont concentrées dans quelques villes $\Rightarrow$ la \textbf{consommation} est concentrée ailleurs que la production.
\item Résultat : chaque région \textbf{vend son surplus} et \textbf{achète son déficit}. Les échanges inter-régionaux naissent de cette double nécessité.
\end{itemize}
\end{aretenir}

\begin{attention}[Ne pas confondre interdépendance et dépendance simple]
Erreur fréquente dans les copies : écrire que « le Nord dépend du Sud » ou que « les campagnes dépendent des villes ». C'est faux ou, du moins, incomplet. La \textbf{dépendance simple} est un flux à sens unique (une région reçoit sans rien donner). L'\textbf{interdépendance} est un échange \textbf{mutuel} : le Nord fournit bœuf, oignons et coton au Sud, et reçoit en retour huile de palme, manioc, poisson fumé et produits manufacturés. Au Bac, montre toujours les \textbf{deux sens} des flux : c'est précisément ce que le mot « inter- » signifie.
\end{attention}

\courssub{Étude de documents : les grandes spécialisations régionales}

Observons d'abord les faits. Le tableau ci-dessous récapitule les principales productions commerciales par grande région (ordres de grandeur et spécialisations dominantes ; les chiffres exacts varient selon les campagnes agricoles).

\begin{center}
\begin{tabularx}{\linewidth}{|l|X|X|}
\hline
\rowcolor{popblueL} \textbf{Grande zone (régions)} & \textbf{Principales productions vendues aux autres régions} & \textbf{Principaux achats aux autres régions} \\
\hline
Extrême-Nord & Oignons, niébé, sorgho, riz (Logone, SEMRY), coton, poisson séché & Huile de palme, manioc, bananes, produits manufacturés \\
\hline
Nord & Bœuf, coton, sorgho, maïs & Produits manufacturés, denrées forestières \\
\hline
Adamaoua & Bœuf (premier pâturage national), maïs, ignames & Produits manufacturés, sel, poisson \\
\hline
Ouest -- Nord-Ouest & Pommes de terre, légumes (carottes, choux, poireaux), maïs, haricots, café arabica & Poisson, huile de palme, produits importés \\
\hline
Littoral & Bananes, plantain, huile de palme, caoutchouc, produits industriels et importés (port de Douala) & Vivres frais, bœuf, oignons \\
\hline
Centre -- Sud -- Est & Cacao, manioc, plantain, poisson fumé, bois & Bœuf, oignons, pommes de terre, produits manufacturés \\
\hline
Sud-Ouest & Bananes, cacao, huile de palme, poisson de mer & Vivres, bœuf, produits manufacturés \\
\hline
\end{tabularx}
\end{center}

La lecture de ce tableau est instructive : \textbf{aucune ligne n'est vide dans la colonne des ventes, aucune dans celle des achats}. C'est la preuve documentaire de l'interdépendance. Complétons par l'observation des photographies de marchés d'approvisionnement : au marché de Mokolo à Yaoundé ou au marché Mboppi à Douala, on voit côte à côte des oignons de Garoua, des pommes de terre de Santa (Nord-Ouest), du gombo de Sangmélima, du poisson fumé de Kribi et du bœuf de Ngaoundéré. Chaque étal est une leçon de géographie à lui tout seul.

\courssub{Les trois grands axes de complémentarité}

\textbf{1. Les complémentarités Nord--Sud.} C'est l'axe le plus spectaculaire. Le Nord (Extrême-Nord, Nord, Adamaoua) fournit au Sud :
\begin{itemize}
\item les \textbf{oignons} de l'Extrême-Nord (vallée de la Logone, environs de Maroua et Garoua), qui approvisionnent tout le pays ;
\item le \textbf{bœuf} de l'Adamaoua et du Nord, premier cheptel national (plusieurs millions de têtes ; l'Adamaoua et le Nord détiennent l'essentiel du troupeau camerounais, de l'ordre de 6 à 7 millions de bovins au niveau national) ;
\item le \textbf{coton} (SODECOTON), qui descend en partie vers les filatures et l'exportation par Douala ;
\item le \textbf{riz} de la SEMRY (Yagoua) et des périmètres irrigués de la Logone.
\end{itemize}
En sens inverse, le Sud envoie vers le Nord les \textbf{produits manufacturés} (savons, tissus, quincaillerie, boissons, matériaux), le \textbf{sel}, l'\textbf{huile de palme}, le \textbf{manioc} et les produits importés par le port de Douala.

\textbf{2. Les complémentarités Ouest/Nord-Ouest -- Métropoles.} Les hautes terres de l'Ouest et du Nord-Ouest, fraîches et fertiles, sont le « grenier maraîcher » du pays. Pommes de terre, légumes (carottes, choux, poireaux, tomates), maïs, haricots et café arabica descendent chaque nuit par camions vers Douala et Yaoundé.

\begin{exemplebox}[L'Ouest et le Nord-Ouest, fournisseurs des métropoles]
L'Ouest et le Nord-Ouest fournissent la \textbf{majeure partie} des pommes de terre et des légumes frais consommés à Douala et Yaoundé : on estime couramment que les hautes terres de l'Ouest (Dschang, Bafoussam, Mbouda) et du Nord-Ouest (Santa, Bambui, Wum) assurent plus des \textbf{deux tiers} de l'approvisionnement national en pommes de terre (ordre de grandeur : la production nationale dépasse 300 000 tonnes par an, dont l'essentiel vient de ces deux régions). Un sac de pommes de terre acheté au marché Mboppi a donc très probablement été arraché la veille sur les pentes des monts Bamboutos.
\end{exemplebox}

\textbf{3. Les complémentarités forêt--savane.} Le Sud forestier (Centre, Sud, Est, Littoral, Sud-Ouest) envoie vers les grandes villes et vers le Nord le \textbf{cacao}, le \textbf{café robusta}, les \textbf{bananes} et le \textbf{plantain}, le \textbf{manioc} (bâtons, gari), l'\textbf{huile de palme} et le \textbf{poisson fumé} des fleuves et des côtes. Ces produits remontent le long de l'axe Yaoundé--Ngaoundéré (route et chemin de fer Transcam) jusqu'aux marchés de Garoua et Maroua, où ils sont revendus aux consommateurs de la savane.

\begin{popfigure}
\begin{center}
\begin{tikzpicture}[scale=1.0]
% Régions (blocs)
\fill[popgoldL] (-3.2,2.2) rectangle (-0.4,4.2);
\node[font=\small\bfseries] at (-1.8,3.6) {NORD};
\node[font=\scriptsize, align=center] at (-1.8,2.9) {bœuf, oignons,\\coton, riz, niébé};
\fill[popgreenL] (-4.6,-0.6) rectangle (-1.6,1.4);
\node[font=\small\bfseries] at (-3.1,0.9) {OUEST / N.-OUEST};
\node[font=\scriptsize, align=center] at (-3.1,0.1) {pommes de terre,\\légumes, maïs, café};
\fill[popblueL] (0.6,-0.6) rectangle (3.8,1.4);
\node[font=\small\bfseries] at (2.2,0.9) {SUD FORESTIER};
\node[font=\scriptsize, align=center] at (2.2,0.1) {cacao, bananes, manioc,\\poisson fumé, huile de palme};
\fill[poppinkL] (-0.8,-3.2) rectangle (2.0,-1.4);
\node[font=\small\bfseries] at (0.6,-1.9) {DOUALA -- YAOUNDÉ};
\node[font=\scriptsize, align=center] at (0.6,-2.7) {produits manufacturés,\\importations, redistribution};
% Flèches proportionnelles
\draw[-{Stealth[length=4mm]}, line width=3.2pt, poporange] (-1.6,2.1) -- (0.2,-1.3);
\draw[-{Stealth[length=4mm]}, line width=2.4pt, poppurple] (0.9,-1.3) -- (-1.5,2.1);
\draw[-{Stealth[length=4mm]}, line width=3.6pt, popgreen] (-2.6,-0.7) -- (-0.2,-1.5);
\draw[-{Stealth[length=4mm]}, line width=1.8pt, popblue] (0.9,-1.5) -- (-2.2,-0.7);
\draw[-{Stealth[length=4mm]}, line width=2.8pt, popblue] (2.0,-0.7) -- (1.1,-1.4);
\draw[-{Stealth[length=4mm]}, line width=2.0pt, poporange] (1.4,-1.4) -- (2.6,-0.7);
\draw[-{Stealth[length=4mm]}, line width=2.2pt, popblue] (2.6,1.5) -- (-0.5,2.6);
\draw[-{Stealth[length=4mm]}, line width=1.6pt, popgold] (-0.3,2.5) -- (2.4,1.5);
% Légende des flèches
\node[font=\scriptsize, anchor=west] at (4.3,3.4) {\textbf{Flux (épaisseur $\propto$ volume)}};
\draw[line width=3.4pt, popgreen] (4.3,2.9) -- (5.3,2.9); \node[font=\scriptsize, anchor=west] at (5.4,2.9) {produits agricoles};
\draw[line width=3.4pt, poporange] (4.3,2.4) -- (5.3,2.4); \node[font=\scriptsize, anchor=west] at (5.4,2.4) {élevage (bœuf)};
\draw[line width=3.4pt, popblue] (4.3,1.9) -- (5.3,1.9); \node[font=\scriptsize, anchor=west] at (5.4,1.9) {pêche, denrées forestières};
\draw[line width=3.4pt, poppurple] (4.3,1.4) -- (5.3,1.4); \node[font=\scriptsize, anchor=west] at (5.4,1.4) {produits manufacturés};
\draw[line width=3.4pt, popgold] (4.3,0.9) -- (5.3,0.9); \node[font=\scriptsize, anchor=west] at (5.4,0.9) {coton, riz, oignons};
\end{tikzpicture}
\end{center}
\legende{Diagramme de flux des principaux échanges inter-régionaux au Cameroun. L'épaisseur des flèches est proportionnelle à l'importance estimée des échanges ; les couleurs distinguent les grandes catégories de produits. Schéma simplifié, sans échelle.}
\end{popfigure}

\courssub{Le rôle pivot des grandes métropoles}

Au cœur de ce système, \textbf{Douala} et \textbf{Yaoundé} jouent un triple rôle :

\begin{enumerate}
\item \textbf{Zones de commercialisation} : leurs grands marchés (Mboppi, Mokolo, Etoudi, Sandaga) sont les plaques tournantes où convergent les produits de toutes les régions et d'où repartent les flux de redistribution.
\item \textbf{Zones de consommation} : avec à elles deux plus de 6 millions d'habitants (ordre de grandeur), elles concentrent le pouvoir d'achat, les administrations, les industries et les services : c'est là que la demande est la plus forte.
\item \textbf{Zones de redistribution} : Douala, premier port du pays (par où transite l'essentiel du commerce extérieur), redistribue les produits importés (riz, farine de blé, matériel, carburant) vers tout l'hinterland, y compris vers le Tchad et la Centrafrique. Yaoundé redistribue vers le Centre, le Sud et l'Est.
\end{enumerate}

\begin{popfigure}
\begin{center}
\includegraphics[width=\linewidth,height=9.5cm,keepaspectratio]{N01/cmr_map.jpg}
\end{center}
\legende{Fond de carte pour la carte de synthèse des échanges inter-régionaux : on y repère les grandes villes de commercialisation et de consommation (Douala, Yaoundé, Bafoussam, Ngaoundéré, Garoua, Maroua), les routes et la voie ferrée qui les relient. La carte de synthèse y ajoute les zones de production (noms de produits en couleur), les points pour les villes et les flèches pour les principaux flux. \\ Source : Wikimedia Commons, CIA, domaine public.}
\end{popfigure}

\begin{methode}[Construire un diagramme de flux]
\begin{enumerate}
\item \textbf{Réunir les données} : liste des produits échangés, origines, destinations, volumes (ou ordres de grandeur).
\item \textbf{Placer les pôles} : régions ou villes sous forme de points ou de blocs, en respectant leur position relative.
\item \textbf{Tracer les flèches} dans le sens de circulation des marchandises (origine $\rightarrow$ destination).
\item \textbf{Proportionner l'épaisseur} de chaque flèche au volume échangé : un flux majeur = trait épais, un flux secondaire = trait fin. Choisis une échelle simple (par exemple 1 mm d'épaisseur pour 100 000 tonnes).
\item \textbf{Soigner la légende} : couleurs par catégorie de produits, échelle d'épaisseur, titre, orientation. Sans légende, un diagramme de flux ne vaut rien au Bac.
\end{enumerate}
\end{methode}

\begin{exoresolu}[Lire et interpréter un diagramme de flux]
\textbf{Énoncé.} À partir du diagramme de flux ci-dessus : (1) identifie les deux régions dont les échanges avec les métropoles sont les plus intenses ; (2) montre que ces échanges traduisent une interdépendance et non une dépendance simple ; (3) explique le rôle de Douala dans ce système.
\tcblower
\textbf{Solution détaillée.}
(1) Les flèches les plus épaisses relient d'une part l'\textbf{Ouest/Nord-Ouest} (pommes de terre, légumes, maïs) et d'autre part le \textbf{Nord/Adamaoua} (bœuf, oignons) aux métropoles Douala et Yaoundé : ce sont les flux vivriers les plus massifs du pays.
(2) Les flux ne sont pas à sens unique : pendant que le Nord envoie bœuf et oignons vers le Sud, il reçoit en retour produits manufacturés, huile de palme, manioc et produits importés (flèches Sud $\rightarrow$ Nord sur le diagramme). Chaque région est à la fois \textbf{fournisseuse} et \textbf{cliente} des autres : c'est la définition même de l'interdépendance.
(3) Douala est à la fois \textbf{porte d'entrée} (le port y décharge les importations), \textbf{marché de consommation} (plus grande agglomération du pays) et \textbf{centre de redistribution} (les produits importés et manufacturés y sont rechargés vers l'intérieur, jusqu'à Maroua et au-delà vers N'Djamena). Elle est le nœud principal du réseau d'échanges national.
\end{exoresolu}

\begin{savaistu}[Le bœuf descend du Nord... sur pied ou par camion]
Le bœuf de l'Adamaoua et du Nord est traditionnellement transporté \textbf{sur pied} : des troupeaux conduits par des bergers peuls parcourent des centaines de kilomètres en plusieurs semaines, de Ngaoundéré ou Garoua jusqu'aux abattoirs et marchés de Yaoundé et de Douala — c'est la \textbf{transhumance commerciale}. Aujourd'hui, les camions-bétaillères et les wagons du Transcam (Transcamerounais, jusqu'à Ngaoundéré) prennent le relais, mais la « descente » des bœufs à pied reste pratiquée. Un zébu acheté 250 000 FCFA à Ngaoundéré peut valoir le double à Douala : cette marge rémunère toute une chaîne d'acteurs (éleveurs, commerçants, transporteurs, bayam selam, bouchers) qui fait vivre l'interdépendance au quotidien.
\end{savaistu}

\courssub{Conséquence : les échanges, ciment de l'intégration nationale}

L'interdépendance n'est pas seulement un fait économique : elle a une portée \textbf{politique et sociale} majeure, au cœur de la famille de situations « l'intégration nationale ».

\begin{aretenir}[De l'interdépendance économique à l'intégration nationale]
\begin{itemize}
\item Les échanges créent une \textbf{solidarité économique concrète} : le consommateur de Yaoundé a besoin de l'éleveur de l'Adamaoua, qui a besoin du port de Douala, qui a besoin du cacao de l'Est.
\item Ils font \textbf{circuler les hommes} autant que les marchandises : commerçants bamilékés installés à Maroua, bouchers haoussa à Douala, bayam selam de toutes origines sur tous les marchés — le marché est un lieu de \textbf{brassage des populations}.
\item Ils justifient et financent les \textbf{infrastructures nationales} (routes, Transcam), qui relient physiquement les régions.
\item Ils diffusent une \textbf{monnaie, des prix et des habitudes de consommation communs} : on mange des oignons du Nord à Ebolowa et du manioc du Sud à Garoua.
\end{itemize}
Ainsi, le commerce intérieur tisse chaque jour la \textbf{nation camerounaise} : faire circuler les produits, c'est faire exister le pays comme un tout solidaire.
\end{aretenir}

\begin{exercice}[Pour t'entraîner]
À partir du tableau des productions régionales et de la carte de synthèse : (1) construis un tableau à double entrée montrant, pour trois régions de ton choix, ce qu'elles vendent et ce qu'elles achètent ; (2) rédige un paragraphe de 10 lignes montrant que l'approvisionnement du marché Mokolo (Yaoundé) illustre l'interdépendance régionale ; (3) explique pourquoi une rupture de la route Douala--Bafoussam (éboulement, pont coupé) se traduit immédiatement par une hausse des prix des pommes de terre à Douala.
\end{exercice}

\begin{corrige}[Exercice]
(1) Exemple avec l'Extrême-Nord : vend oignons, riz, niébé, bœuf ; achète huile de palme, manioc, produits manufacturés. Ouest : vend pommes de terre, légumes, maïs, café ; achète poisson, huile de palme, produits importés. Centre : vend cacao, manioc, plantain ; achète bœuf, oignons, pommes de terre.
(2) Le marché Mokolo reçoit chaque jour des produits venus de presque toutes les régions : pommes de terre de l'Ouest et du Nord-Ouest, oignons de l'Extrême-Nord, bœuf de l'Adamaoua, poisson fumé du Sud et des côtes, bananes et manioc du Centre et du Littoral. Inversement, les recettes de ces ventes permettent aux commerçants d'acheter des produits manufacturés et importés redistribués depuis Yaoundé et Douala. Chaque région y est donc à la fois vendeuse et acheteuse : Mokolo matérialise l'interdépendance.
(3) La pomme de terre consommée à Douala vient pour l'essentiel des hautes terres de l'Ouest et du Nord-Ouest, par la route Douala--Bafoussam. Si cet axe est coupé, l'offre s'effondre brutalement à Douala alors que la demande reste forte : la rareté fait monter les prix. Cette sensibilité des prix à une route prouve à quel point la métropole \textbf{dépend} de ses régions d'approvisionnement — et donc la force de l'interdépendance.
\end{corrige}

Retenons l'essentiel : les infrastructures étudiées dans la section précédente ne sont que des instruments ; c'est l'interdépendance — née de la diversité des milieux et de la concentration urbaine de la consommation — qui met les marchandises en mouvement. Reste à voir comment ce commerce intérieur est concrètement \textbf{organisé} : marchés périodiques, bayam selam, vendeurs à la sauvette et centres commerciaux. C'est l'objet de la section suivante.

\coursec{L'organisation du commerce intérieur}

Nous venons de voir que les régions du Cameroun sont profondément interdépendantes : le maïs de l'Ouest, les oignons du Nord, le cacao du Centre, le poisson du Littoral circulent en permanence d'une région à l'autre. Mais une question demeure : \emph{comment ces produits passent-ils concrètement des mains du producteur à celles du consommateur ?} Qui les collecte, qui les transporte, qui les vend, et où ? Autrement dit, comment le commerce intérieur est-il \cle{organisé} ?

Pour répondre à cette question, partons d'une observation simple. À Yaoundé, une ménagère peut acheter ses tomates de quatre façons très différentes : au marché Mokolo auprès d'une revendeuse, à un vendeur installé sur le trottoir au carrefour, dans une boutique de quartier, ou dans le supermarché d'un centre commercial. Ces quatre lieux de vente coexistent, mais ils ne fonctionnent pas du tout de la même manière et ne s'adressent pas aux mêmes clients. C'est toute la richesse — et la complexité — du commerce intérieur camerounais, que nous allons maintenant étudier en quatre temps : les marchés périodiques, les revendeurs de vivres, le commerce informel de rue et le commerce moderne.

\courssub{Les marchés périodiques : le cœur du commerce rural}

\textbf{Intuition et observation.} Si tu te rends un vendredi à Bafou (Ouest) ou un jeudi à un marché de la plaine des Mbo, tu verras des milliers de personnes converger vers un même lieu : des paysannes portant des cuvettes de maïs, de haricots ou de bananes sur la tête, des commerçants venus de la ville avec des camionnettes, des acheteurs qui marchandent. Le lendemain, le même lieu est presque désert. Pourquoi ce rythme ?

\begin{definition}[Marché périodique]
Un \cle{marché périodique} est un marché qui se tient à \cle{jours fixes} — le plus souvent une fois par semaine (marché hebdomadaire) — dans une localité donnée, et où convergent producteurs et acheteurs venus des villages environnants. Il est dit « périodique » parce qu'il ne fonctionne pas tous les jours, contrairement au marché permanent.
\end{definition}

\textbf{Mécanisme et rôle économique.} Le marché périodique joue un double rôle fondamental dans le commerce intérieur :

\begin{itemize}
\item un rôle de \cle{collecte} : les petits producteurs des villages environnants y apportent leurs surplus (vivres, mais aussi bétail, volailles, artisanat). Les commerçants urbains ou les revendeurs y achètent en gros ;
\item un rôle de \cle{redistribution} : les villageois y achètent en retour ce qu'ils ne produisent pas (huile, savon, tissus, poisson fumé, sel, outils, produits manufacturés importés).
\end{itemize}

Le marché périodique est donc un \emph{nœud} du réseau commercial : il concentre une offre dispersée et diffuse les marchandises venues d'ailleurs. Il a aussi une fonction sociale (rencontres, informations, mariages) qui explique sa vitalité.

\begin{exemplebox}[Les grands marchés périodiques du Cameroun]
Dans la région de l'Ouest, les marchés hebdomadaires des Bamboutos et des Hauts-Plateaux (marchés de Mbouda, Bafou, Dschang et des villages de la plaine des Mbo) collectent le maïs, le haricot, les pommes de terre et les légumes destinés à Douala et Yaoundé. Dans le Nord, les grands marchés hebdomadaires (par exemple autour de Garoua, Maroua et des lamidats) sont des lieux d'échange de mil, d'oignons, de coton et surtout de \cle{bétail} : les marchés à bœufs y approvisionnent tout le sud du pays. Chaque localité a son « jour de marché », décalé par rapport aux localités voisines, ce qui permet aux commerçants de tourner d'un marché à l'autre au cours de la semaine.
\end{exemplebox}

\begin{attention}[Ne pas confondre marché périodique et marché permanent]
Le marché périodique se tient à jours fixes (une fois par semaine en général), surtout en milieu rural. Le marché permanent fonctionne tous les jours, surtout en ville (marché Mokolo à Yaoundé, marché central et marché Mboppi à Douala). Dans une copie, précise toujours de quel type de marché tu parles.
\end{attention}

\courssub{Les revendeurs de vivres : les « bayam selam », maillons indispensables}

\textbf{Observation.} Au marché Mokolo (Yaoundé), dès 5 heures du matin, des femmes déchargent des sacs de macabo, de manioc, de plantain arrivés pendant la nuit depuis les villages du Centre, du Sud et de l'Ouest. Ces commerçantes ne cultivent pas : elles achètent pour revendre. On les appelle les \emph{bayam selam}.

\begin{definition}[Bayam selam]
Un(e) \cle{bayam selam} est un revendeur ou une revendeuse de vivres qui achète les produits agricoles dans les zones de production (villages, marchés périodiques de collecte) pour les revendre dans les marchés urbains. Le terme vient de l'anglais pidgin \emph{« buyam sellam »} (contraction de \emph{buy them, sell them}), littéralement « achète-les, vends-les ». Ce métier est majoritairement exercé par des femmes.
\end{definition}

\textbf{Mécanisme : le circuit bayam selam.} Le bayam selam est un \cle{intermédiaire} entre le producteur et le consommateur urbain. Son circuit suit des étapes bien identifiables : achat au village ou au marché de collecte, groupage des marchandises, transport (camions, taxis-brousse, pick-up), revente en gros ou en demi-gros au marché urbain, puis vente au détail par les détaillantes aux consommateurs. Chaque étape ajoute un coût (transport, taxes de marché, pertes) et une marge, ce qui explique l'écart de prix entre le village et la ville.

\begin{popfigure}
\begin{center}
\begin{tikzpicture}[font=\small, node distance=0.55cm,
  etape/.style={draw=popblue, fill=popblueL, rounded corners=2pt, minimum height=0.85cm, minimum width=3.1cm, align=center},
  fleche/.style={-{Stealth[length=2.5mm]}, thick, poporange}]
\node[etape, align=center] (v){Village de production\\ (paysans)};
\node[etape, below=of v, align=center] (mc){Marché périodique\\ de collecte};
\node[etape, below=of mc, fill=popgoldL, draw=popgold, align=center] (b){\textbf{Bayam selam}\\ achat, groupage, transport};
\node[etape, below=of b, align=center] (mu){Marché urbain de gros\\ (Mokolo, Mboppi)};
\node[etape, below=of mu, align=center] (d){Détaillants\\ (étals, boutiques)};
\node[etape, below=of d, fill=popgreenL, draw=popgreen] (c) {Consommateur urbain};
\draw[fleche] (v) -- (mc);
\draw[fleche] (mc) -- (b);
\draw[fleche] (b) -- node[right, popdark, font=\scriptsize]{camion, taxi-brousse} (mu);
\draw[fleche] (mu) -- (d);
\draw[fleche] (d) -- (c);
\end{tikzpicture}
\end{center}
\legende{Le circuit bayam selam : du champ du producteur à l'assiette du consommateur urbain. À chaque étape, le prix augmente (coûts de transport, marges, taxes de marché).}
\end{popfigure}

\textbf{Pourquoi sont-ils indispensables ?} Sans les bayam selam, le paysan n'aurait pas les moyens d'écouler ses vivres jusqu'en ville (pas de véhicule, pas de réseau de clients), et le citadin n'aurait pas accès aux produits frais. Ils assurent donc la \cle{jonction spatiale} entre zones de production et zones de consommation, et financent souvent la campagne agricole en avançant de l'argent aux producteurs. En contrepartie, leur position d'intermédiaire leur donne un pouvoir sur les prix, parfois critiqué.

\begin{savaistu}[D'où vient le mot « bayam selam » ?]
L'expression vient de l'anglais pidgin parlé dans les régions du Nord-Ouest et du Sud-Ouest : \emph{buyam sellam}, contraction de \emph{buy them, sell them} (« achète-les, vends-les »). Elle désigne à l'origine les commerçantes qui achetaient les vivres dans les zones de production anglophones pour les revendre à Douala. Le mot est aujourd'hui passé dans le langage courant de tout le Cameroun, et le métier reste l'un des principaux secteurs d'activité des femmes dans l'économie urbaine.
\end{savaistu}

\courssub{Les vendeurs à la sauvette : le commerce de rue}

\textbf{Observation.} Aux carrefours de Douala (Akwa, Deïdo) et de Yaoundé (Poste centrale, rond-point Express), des centaines de vendeurs étalent leurs marchandises sur des bâches, des cagettes ou simplement à même le sol : vêtements, chaussures, téléphones, fruits, beignets. Dès l'arrivée des agents municipaux, tout disparaît en quelques secondes... pour réapparaître une heure plus tard.

\begin{definition}[Vendeur à la sauvette]
Un \cle{vendeur à la sauvette} est un commerçant informel qui vend ses marchandises dans la rue, sur les trottoirs ou aux carrefours, sans local commercial ni autorisation administrative. Il ne paie généralement ni loyer ni impôt, et peut déménager son activité à tout moment — d'où l'expression « à la sauvette ».
\end{definition}

\textbf{Importance et limites.} Ce commerce de rue joue un rôle social majeur : il offre un revenu à des milliers de jeunes et de chômeurs dans les grandes villes, et propose aux consommateurs pauvres des marchandises à petits prix, vendues au détail le plus fin (une cigarette, un sachet d'eau, un fruit). Mais il pose de sérieux problèmes :

\begin{itemize}
\item \cle{occupation de l'espace public} : trottoirs et chaussées encombrés, gêne pour les piétons et la circulation, insalubrité ;
\item \cle{fiscalité} : ces vendeurs échappent à l'impôt et aux taxes de marché, ce qui prive l'État et les communes de ressources et crée une concurrence jugée déloyale par les commerçants installés ;
\item précarité : pas de protection sociale, marchandises parfois saisies, revenus irréguliers.
\end{itemize}

\courssub{Le commerce moderne : boutiques, supermarchés et centres commerciaux}

À l'autre extrémité du système se trouve le \cle{commerce moderne} : boutiques fixes avec registres et licences, supermarchés (à Douala et Yaoundé notamment) et \cle{centres commerciaux} regroupant magasins, cinémas et restaurants sous un même toit, comme le \textbf{Douala Grand Mall}, le plus grand centre commercial du pays, inauguré en 2020 à Douala (quartier Bonanjo, près de l'aéroport), ou les supermarchés de Yaoundé (Mahima, Casino, Supermarché du Centre, quartiers du centre-ville et de la périphérie résidentielle comme Bastos, Nlongkak).

Ce commerce formel offre des produits standardisés, étiquetés, souvent importés, dans un cadre climatisé et sécurisé. Mais il cible une \cle{clientèle urbaine aux revenus élevés} : cadres, expatriés, classes moyennes. Il reste donc marginal en volume dans la distribution des vivres, même s'il est très visible et en croissance.

\begin{exemplebox}[Le poids écrasant du secteur informel]
Au Cameroun, le secteur informel (bayam selam, vendeurs à la sauvette, petits commerçants, marchés traditionnels) assure l'essentiel de la distribution des vivres dans les villes : la très grande majorité des ménages urbains s'approvisionne en produits frais sur les marchés et auprès des revendeuses, et non dans les supermarchés. Le secteur informel fournirait, selon les estimations de l'INS et du BIT (ordres de grandeur), environ \textbf{90 \% de l'emploi total hors agriculture} et une part prépondérante du PIB non pétrolier. Le commerce moderne, bien que dynamique, ne concerne qu'une minorité de la consommation alimentaire quotidienne.
\end{exemplebox}

\courssub{La hiérarchie des lieux de vente}

L'ensemble de ces lieux de vente forme une \cle{hiérarchie commerciale} : à la base, une multitude de petits points de vente fréquentés quotidiennement ; au sommet, quelques grands équipements rares, fréquentés occasionnellement, qui offrent des services plus élaborés.

\begin{popfigure}
\begin{center}
\begin{tikzpicture}[font=\small]
\fill[popgreenL, draw=popgreen] (-4.5,0) -- (4.5,0) -- (3.2,1.3) -- (-3.2,1.3) -- cycle;
\fill[popblueL, draw=popblue] (-3.2,1.3) -- (3.2,1.3) -- (2.1,2.6) -- (-2.1,2.6) -- cycle;
\fill[popgoldL, draw=popgold] (-2.1,2.6) -- (2.1,2.6) -- (1.1,3.9) -- (-1.1,3.9) -- cycle;
\fill[poppinkL, draw=poppink] (-1.1,3.9) -- (1.1,3.9) -- (0,5.2) -- cycle;
\node at (0,0.65) {\textbf{Marchés périodiques ruraux} (très nombreux, petits, hebdomadaires)};
\node at (0,1.95) {\textbf{Marchés urbains permanents} (Mokolo, Central, Mboppi)};
\node at (0,3.25) {\textbf{Boutiques et supermarchés} (villes)};
\node at (0,4.55) {\textbf{Centres commerciaux} (rares)};
\node[rotate=90, popdark, font=\scriptsize] at (-5.3,2.6) {prix et services croissants};
\draw[-{Stealth}, thick, popdark] (-5.0,0.2) -- (-5.0,5.0);
\node[rotate=90, popdark, font=\scriptsize] at (5.3,2.6) {nombre de lieux décroissant};
\draw[-{Stealth}, thick, popdark] (5.0,5.0) -- (5.0,0.2);
\end{tikzpicture}
\end{center}
\legende{Pyramide de la hiérarchie commerciale au Cameroun : à la base, des lieux de vente nombreux, proches et bon marché ; au sommet, des équipements rares, chers et spécialisés. Les vendeurs à la sauvette et les bayam selam s'insèrent à tous les niveaux inférieurs et intermédiaires.}
\end{popfigure}

On lit cette pyramide ainsi : le consommateur se rend chaque jour au marché ou au vendeur de rue (biens courants, faible rayon d'attraction), mais seulement de temps en temps au supermarché ou au centre commercial (biens durables, vêtements de marque, loisirs), ce qui suppose un revenu suffisant et souvent un moyen de transport.

\begin{attention}[Erreur fréquente à éviter]
Ne réduis jamais le commerce intérieur camerounais au secteur moderne (supermarchés, centres commerciaux). C'est le piège classique du candidat qui ne décrit que ce qu'il voit dans les quartiers riches. En réalité, le commerce \cle{traditionnel et informel} — marchés périodiques, marchés urbains, bayam selam, vendeurs à la sauvette — domine largement, tant en volume de marchandises distribuées qu'en nombre d'emplois. Une bonne copie présente toujours les deux secteurs et précise leur poids respectif.
\end{attention}

\begin{methode}[Comment décrire et analyser une photographie de géographie commerciale (Bac)]
\begin{enumerate}
\item \textbf{Identifier le document} : nature (photo), lieu approximatif (indices : architecture, végétation, pancartes), date si connue.
\item \textbf{Décrire ce qui est visible} : disposition de l'espace (étals, bâches, rayonnages, bitume/terre), acteurs (hommes/femmes, âge, tenues), marchandises (type, conditionnement, quantité), moyens de transport présents.
\item \textbf{Classer le type de commerce} : marché périodique / permanent, commerce informel (sauvette, bayam selam), commerce moderne (boutique, supermarché, centre commercial). Justifier par les indices vus.
\item \textbf{Analyser les fonctions} : collecte, redistribution, vente au détail, services. Liens avec l'aménagement (voirie, électricité, stationnement).
\item \textbf{Conclure} : insérer le lieu dans la hiérarchie commerciale et l'interdépendance des régions.
\end{enumerate}
\end{methode}

\courssub{Étude de documents : quatre photographies du commerce intérieur}

\begin{exoresolu}[Décrire et comparer quatre lieux de vente]
\textbf{Énoncé.} On dispose de quatre photographies : (A) un marché périodique de l'Ouest un jour de marché ; (B) des bayam selam déchargeant des sacs de vivres au marché Mokolo à l'aube ; (C) des vendeurs à la sauvette sur un trottoir d'Akwa (Douala) ; (D) l'intérieur climatisé du Douala Grand Mall. Pour chaque document, décris ce que tu vois, identifie le type de commerce, puis dégage un point commun et deux différences entre ces lieux de vente.

\tcblower
\textbf{Solution détaillée.}
\begin{itemize}
\item \textbf{Photo A} : une foule dense en plein air, sur sol souvent non bitumé, des cuvettes de maïs et de haricots posées au sol ou sur des tables sommaires, des camionnettes bâchées chargées en arrière-plan. Il s'agit d'un \cle{marché périodique} rural : il assure la collecte des produits agricoles et la redistribution des marchandises vers les villages.
\item \textbf{Photo B} : des femmes (majoritairement) marchandant de gros sacs de plantain, de macabo, de manioc à côté d'un camion stationné sur une voirie bitumée, sous des hangars en tôle. Ce sont des \cle{bayam selam}, intermédiaires qui achètent les vivres dans les zones de production pour les revendre dans les marchés urbains de gros.
\item \textbf{Photo C} : des marchandises (chaussures, téléphones portables, accessoires, fruits) étalées sur une bâche ou des cartons sur le trottoir bitumé, à proximité immédiate de la chaussée, pas de local en dur. C'est le commerce informel des \cle{vendeurs à la sauvette}, sans autorisation ni fiscalité, à la merci des opérations de déguerpissement.
\item \textbf{Photo D} : des rayonnages métalliques éclairés par néons, des produits emballés et étiquetés (codes-barres), des chariots de course, une clientèle peu nombreuse, climatisation apparente. C'est le \cle{commerce moderne} (centre commercial), qui cible une clientèle urbaine aux revenus élevés.
\end{itemize}
\textbf{Point commun} : les quatre lieux assurent la même fonction — mettre des marchandises à la disposition des consommateurs — et participent tous au commerce intérieur national. \textbf{Différences} : (1) le statut juridique et l'aménagement : A, B et C relèvent du commerce traditionnel ou informel (peu ou pas d'infrastructures lourdes), D du commerce formel et moderne (infrastructures lourdes, normes) ; (2) la clientèle et la fixation des prix : A, B et C servent l'ensemble de la population, y compris les plus pauvres, avec des prix négociés à l'unité, tandis que D s'adresse à une clientèle aisée avec des prix fixes, affichés et payables souvent par carte. \textbf{Conclusion} : le commerce intérieur camerounais est un système à deux vitesses où le secteur informel et traditionnel domine en volume et en emploi, tandis que le secteur moderne, minoritaire, se développe dans les grandes métropoles.
\end{exoresolu}

\begin{exercice}[Pour t'entraîner]
\begin{enumerate}
\item Définis les termes suivants : marché périodique, bayam selam, vendeur à la sauvette.
\item Reproduis et complète le schéma du circuit bayam selam en indiquant, à chaque étape, un facteur qui fait augmenter le prix du produit.
\item Explique en cinq à huit lignes pourquoi le commerce informel domine la distribution des vivres dans les villes camerounaises, malgré le développement des supermarchés.
\end{enumerate}
\end{exercice}

\begin{corrige}[Exercice]
\begin{enumerate}
\item \textbf{Marché périodique} : marché se tenant à jours fixes (généralement hebdomadaire) en milieu rural, rassemblant producteurs locaux et acheteurs pour la collecte et la redistribution. \textbf{Bayam selam} : revendeur(se) de vivres assurant l'intermédiation entre zones de production rurales et marchés urbains (achat, groupage, transport, revente). \textbf{Vendeur à la sauvette} : commerçant informel occupant l'espace public (trottoirs, carrefours) sans local ni autorisation administrative, vendant au détail fin.
\item \textbf{Facteurs de hausse du prix par étape} : 
\begin{itemize}
\item Village $\to$ Marché de collecte : coût de transport primaire (moto, portage), taxe de collecte locale.
\item Marché de collecte $\to$ Bayam selam : marge d'achat du bayam selam, coût de groupage.
\item Bayam selam $\to$ Marché urbain (transport) : carburant, péages, entretien véhicule (routes dégradées), pertes/perte de qualité (pourriture), taxes de barreau/entrée de ville.
\item Marché urbain (gros) $\to$ Détaillant : taxe de place de marché, marge du grossiste/demi-gros, manutention.
\item Détaillant $\to$ Consommateur : marge de la détaillante, conditionnement (sachets, bols), localisation du point de vente (loyer si boutique).
\end{itemize}
\item Le commerce informel domine la distribution des vivres urbains pour trois raisons principales. (1) \textbf{Adaptation au pouvoir d'achat} : il pratique la vente à l'unité ou au tas (détail fin), prix négociables, et parfois le crédit de confiance, accessibles aux ménages à faibles revenus majoritaires. (2) \textbf{Proximité et accessibilité} : marchés de quartier, vendeurs de rue au coin de la rue, horaires élargis (très tôt le matin, tard le soir), accessibles à pied ou par transport en commun. (3) \textbf{Faibles barrières à l'entrée} : pas de capital important, pas de formalités administratives, main-d'œuvre familiale, ce qui en fait le premier pourvoyeur d'emplois urbains (environ 90 \% selon l'INS/BIT). Les supermarchés, eux, imposent des prix fixes, des conditionnements standards, une localisation périphérique exigeant un véhicule, et ciblent les classes moyennes et supérieures (moins de 10 \% des ménages).
\end{enumerate}
\end{corrige}

\begin{aretenir}[L'essentiel de la section]
Le commerce intérieur camerounais repose sur quatre formes complémentaires et hiérarchisées : les \cle{marchés périodiques} (collecte et redistribution en milieu rural, hebdomadaires), les \cle{bayam selam} (intermédiaires indispensables entre campagnes et villes, majoritairement femmes), les \cle{vendeurs à la sauvette} (commerce informel de rue, source massive d'emplois mais problèmes d'espace public et de fiscalité) et le \cle{commerce moderne} (supermarchés et centres commerciaux comme Douala Grand Mall, réservé à une clientèle aisée). Le secteur traditionnel et informel domine largement la distribution des vivres (volume, emplois) : c'est la réalité structurelle qu'il faut toujours mettre en avant dans une copie du Baccalauréat.
\end{aretenir}

\coursec{Les difficultés du commerce intérieur}

Nous venons de voir que le commerce intérieur camerounais repose sur un tissu vivant de marchés périodiques, de \cle{bayam-sellam}, de vendeurs à la sauvette et de centres commerciaux modernes. Mais ce système, si dynamique soit-il, fonctionne dans la douleur. Qui n'a jamais entendu un commerçant de Mokolo à Yaoundé se plaindre : « Les tomates arrivent chères et pourries depuis le Nord-Ouest » ? Qui n'a jamais vu un camion de vivres immobilisé trois jours sur une piste défoncée entre Batouri et Bertoua ? Ces scènes quotidiennes révèlent un paradoxe : le Cameroun produit abondamment, mais ses produits circulent mal et coûtent cher. Comprendre les \cle{difficultés du commerce intérieur}, c'est comprendre pourquoi le panier de la ménagère s'alourdit alors que les champs regorgent de denrées.

\courssub{Un problème d'ensemble : présentation et démarche}

Les difficultés du commerce intérieur ne sont pas un simple inventaire de maux isolés : elles forment un \cle{système de contraintes} qui s'enchaînent. Une route dégradée allonge le temps de transport ; le temps de transport pourrit les denrées ; les pertes obligent le commerçant à répercuter ses frais sur le prix ; le prix élevé décourage le consommateur et appauvrit le producteur qui n'a pas été payé à sa juste valeur. Pour analyser ce système, les géographes utilisent un outil puissant : l'\cle{arbre des causes}.

\begin{methode}[Construire un arbre des causes]
Pour analyser un problème géographique complexe (ici : « le commerce intérieur camerounais fonctionne mal »), procède en quatre étapes :
\begin{enumerate}
\item \textbf{Formuler le problème central} en une phrase claire, placée au sommet ou au centre du schéma.
\item \textbf{Classer les causes par catégories} (infrastructures, coûts, acteurs, financement, sécurité...) : chaque catégorie forme une branche principale.
\item \textbf{Ramifier} : sous chaque catégorie, précise les causes concrètes (ex. branche « infrastructures » $\rightarrow$ « routes dégradées », « rail vétuste »...).
\item \textbf{Associer une conséquence et une solution à chaque cause} : un bon arbre des causes débouche toujours sur des pistes d'action, ce qui est exactement la démarche attendue au Baccalauréat dans un commentaire de document ou une étude de cas.
\end{enumerate}
\end{methode}

\begin{popfigure}
\begin{tikzpicture}[
  font=\small,
  problem/.style={rectangle, draw=popdark, fill=popink!20, rounded corners, align=center, text width=4.4cm, font=\bfseries},
  branche/.style={rectangle, draw=popblue, fill=popblueL, rounded corners, align=center, text width=2.9cm, font=\bfseries\small},
  cause/.style={rectangle, draw=popmuted, fill=popcream, rounded corners, align=center, text width=3.2cm, font=\footnotesize},
  lien/.style={-{Stealth[length=2mm]}, thick, popdark}
]
\node[problem] (P) at (0,0) {Difficultés du commerce intérieur au Cameroun};

\node[branche] (B1) at (-6.0,-2.2) {Infrastructures déficientes};
\node[branche] (B2) at (-3.0,-2.2) {Coût élevé du transport};
\node[branche] (B3) at (0,-2.2) {Multiplicité des intermédiaires};
\node[branche] (B4) at (3.0,-2.2) {Tracasseries et insécurité};
\node[branche] (B5) at (6.0,-2.2) {Financement et conservation faibles};

\node[cause] (C1) at (-6.0,-4.5) {Routes dégradées, rail vétuste, ports mal connectés, aérien en difficulté};
\node[cause] (C2) at (-3.0,-4.5) {Carburant, usure véhicules, pertes post-récolte};
\node[cause] (C3) at (0,-4.5) {Circuits allongés, prix gonflés au consommateur final};
\node[cause] (C4) at (3.0,-4.5) {Contrôles multiples, racket, grand banditisme, crises régionales};
\node[cause] (C5) at (6.0,-4.5) {Manque de crédit, absence de chaîne du froid et d'entrepôts};

\draw[lien] (P) -- (B1); \draw[lien] (P) -- (B2); \draw[lien] (P) -- (B3);
\draw[lien] (P) -- (B4); \draw[lien] (P) -- (B5);
\draw[lien] (B1) -- (C1); \draw[lien] (B2) -- (C2); \draw[lien] (B3) -- (C3);
\draw[lien] (B4) -- (C4); \draw[lien] (B5) -- (C5);
\end{tikzpicture}
\legende{Arbre des causes des difficultés du commerce intérieur au Cameroun. Lecture : le problème central (en haut) se ramifie en cinq catégories de causes, elles-mêmes déclinées en causes concrètes. Chaque cause appelle une conséquence (hausse des prix, pertes, enclavement) et une solution ciblée.}
\end{popfigure}

\courssub{Les difficultés liées aux infrastructures de transport}

\begin{definition}[Enclavement]
L'\cle{enclavement} est la situation d'une région mal desservie par les voies de communication (routes, voies ferrées, voies navigables), ce qui isole sa production des marchés de consommation. Une région enclavée produit mais ne vend pas, ou vend à perte.
\end{definition}

L'intuition est simple : on ne peut pas commercer ce qu'on ne peut pas transporter. Or l'observation du réseau camerounais montre quatre faiblesses majeures.

\textbf{Des routes dégradées et des bassins de production enclavés.} Le réseau routier bitumé s'allonge (environ \num{6500} km de routes revêtues sur un réseau classé d'environ \num{50000} km, ordres de grandeur 2023), mais l'entretien est irrégulier. En saison des pluies, des axes stratégiques comme la route Batouri--Bertoua (Est), l'axe Foumban--Mamfé (Ouest--Nord-Ouest) ou certaines pistes de l'Adamaoua deviennent quasi impraticables. Conséquence directe : les bassins de production vivrière (Ouest, Nord-Ouest, Adamaoua) et les zones d'élevage (Nord, Adamaoua) sont périodiquement coupés des grandes villes consommatrices que sont Douala et Yaoundé.

\textbf{Un chemin de fer vétuste.} La Transcamerounaise (Douala--Yaoundé--Ngaoundéré, environ \num{886} km) souffre d'un matériel roulant vieillissant, de déraillements répétés (celui d'Eséka en octobre 2016 a marqué les esprits) et de lenteurs. Le fret ferroviaire, pourtant le moins cher pour les marchandises lourdes (bois, coton, ciment, hydrocarbures), ne joue plus pleinement son rôle de désenclavement du Nord.

\textbf{Des ports mal connectés à l'intérieur.} Le port de Douala concentre l'essentiel du trafic maritime, mais sa desserte vers l'hinterland repose presque exclusivement sur la route et un rail fatigué. Le port en eau profonde de Kribi, mis en service en 2014 (première phase), attend encore le plein achèvement de ses connexions ferroviaires et routières vers l'arrière-pays.

\textbf{Une compagnie aérienne en difficultés.} Le fret aérien, marginal pour les vivres, reste handicapé par les difficultés récurrentes de la compagnie nationale Camair-Co, ce qui pénalise les produits à haute valeur et périssables (fleurs coupées, fruits d'exportation, poissons frais).

\courssub{Le coût élevé du transport et les pertes post-récolte}

Le transport est le premier poste de coût dans le prix d'une denrée vendue en ville. Carburant, péages, usure accélérée des camions sur pistes défoncées, retards : tout se paie, et c'est le consommateur qui paie au bout de la chaîne.

\begin{exemplebox}[Le voyage d'un cageot de tomates]
Un cageot de tomates acheté \num{5000} FCFA au producteur à Santa (Nord-Ouest) peut se vendre \num{12000} à \num{15000} FCFA au détail au marché Mokolo à Yaoundé, soit un doublement, voire un triplement du prix. La différence ne profite pas au producteur : elle absorbe le transport, les commissions des intermédiaires, les tracasseries et les pertes dues au pourrissement en route (ordres de grandeur variables selon les saisons).
\end{exemplebox}

\begin{popfigure}
\begin{center}
\begin{tikzpicture}
\begin{axis}[
  width=0.9\linewidth, height=7cm,
  xlabel={Étape du circuit de commercialisation},
  ylabel={Prix du cageot de tomates (FCFA)},
  symbolic x coords={Santa (champ), Collecteur rural, Transporteur, Grossiste Yaoundé, Détaillant Mokolo},
  xtick=data, x tick label style={rotate=20, anchor=east, font=\footnotesize},
  ymin=0, ymax=17000,
  ymajorgrids=true, grid style={popline!40},
  legend style={at={(0.02,0.98)}, anchor=north west, font=\small}
]
\addplot[ybar, fill=popblue, bar width=14pt] coordinates {
  (Santa (champ),5000) (Collecteur rural,6500) (Transporteur,8500) (Grossiste Yaoundé,11000) (Détaillant Mokolo,14000)
};
\addlegendentry{Prix cumulé du cageot}
\addplot[popcurve, poporange, mark=*] coordinates {
  (Santa (champ),5000) (Collecteur rural,6500) (Transporteur,8500) (Grossiste Yaoundé,11000) (Détaillant Mokolo,14000)
};
\addlegendentry{Tendance : quasi-triplement du prix}
\end{axis}
\end{tikzpicture}
\end{center}
\legende{Évolution indicative du prix d'un cageot de tomates du champ au consommateur urbain (valeurs en ordres de grandeur). Chaque intermédiaire et chaque étape de transport ajoute sa marge : le prix final peut atteindre près de trois fois le prix au producteur.}
\end{popfigure}

À ce coût s'ajoutent les \cle{pertes post-récolte} : faute de camions frigorifiques, de hangars ventilés et d'entrepôts, tomates, bananes, mangues, choux et poissons pourrissent ou s'avariennent pendant le transport et le stockage.

\begin{savaistu}[Le gâchis invisible]
Faute de \cle{chaîne du froid} et d'infrastructures de stockage, une part importante des fruits et légumes produits au Cameroun — les estimations évoquent couramment 30 à 40 \% pour les denrées les plus périssables (ordre de grandeur) — se perd avant même d'atteindre le consommateur. C'est un double gâchis : le producteur perd son revenu, et la ville perd des vivres qui auraient fait baisser les prix.
\end{savaistu}

\courssub{La multiplicité des intermédiaires}

Entre le producteur et le consommateur, la marchandise passe souvent par cinq ou six mains : collecteur au village, revendeuse au marché périodique, bayam-sellam qui finance et achemine, grossiste urbain, demi-grossiste, détaillant. Chaque passage de main ajoute une marge légitime (chacun doit vivre), mais l'addition finale est lourde.

\begin{propriete}[Effet de l'allongement des circuits sur les prix]
Plus le circuit de commercialisation compte d'intermédiaires, plus l'écart se creuse entre le \textbf{prix au producteur} (faible) et le \textbf{prix au consommateur} (élevé). Ce double niveau de prix appauvrit à la fois le paysan, mal rémunéré, et le citadin, qui paie cher : c'est le phénomène de « ciseaux » des prix agricoles.
\end{propriete}

\begin{attention}[Erreur fréquente au Baccalauréat]
Ne te contente jamais de citer une difficulté (« les routes sont mauvaises ») sans la relier à sa \textbf{conséquence concrète}. Le correcteur attend un raisonnement en chaîne : route dégradée $\rightarrow$ temps de transport allongé $\rightarrow$ pourrissement des denrées et coût du fret majoré $\rightarrow$ hausse des prix en ville et enclavement du bassin de production. Une difficulté sans conséquence, c'est une copie à moitié réussie.
\end{attention}

\courssub{Tracasseries routières, insécurité et concurrence extérieure}

\textbf{Les tracasseries routières.} Sur les grands axes (Douala--Yaoundé, Yaoundé--Bafoussam, corridor Ngaoundéré--Garoua), les transporteurs dénoncent la multiplication des contrôles (police, gendarmerie, douanes, contrôle vétérinaire, poids et mesures) et surtout les \cle{péages informels} : sommes indues exigées à chaque barrage, véritable racket qui renchérit mécaniquement le coût des marchandises. Un camion de vivres peut s'arrêter des dizaines de fois sur 300 km ; chaque arrêt coûte du temps et de l'argent, répercutés sur le prix final.

\textbf{L'insécurité.} Le grand banditisme (coupeurs de route) sur certains axes du Nord et de l'Est, ainsi que les crises sécuritaires (régions du Nord-Ouest et du Sud-Ouest depuis 2017, Extrême-Nord face à Boko Haram) perturbent gravement les flux : marchés périodiques fermés ou désertés, itinéraires détournés et allongés, primes de risque exigées par les transporteurs.

\textbf{La concurrence des produits importés.} Riz importé d'Asie, huile, poulets congelés, oignons, tissus : ces produits, souvent moins chers à l'arrivée que les productions locales pénalisées par leurs coûts de transport internes, concurrencent déloyalement le paysan camerounais... sur son propre marché. Paradoxe amer : il arrive que du riz venu de l'autre bout du monde coûte moins cher à Douala que le riz de Yagoua, faute de voies pour descendre celui-ci.

\courssub{La faiblesse du financement, de la conservation et la concurrence de la grande distribution}

Le petit commerçant camerounais travaille sans filet. Les banques classiques prêtent peu aux acteurs de l'informel, jugés non solvables ; le crédit, quand il existe, passe par les \cle{tontines} ou par des prêteurs à taux élevés. Résultat : la bayam-sellam ne peut pas acheter en gros, ni louer un camion seule, ni investir dans un entrepôt. L'absence quasi totale de chaîne du froid et de magasins de stockage dans les zones de production achève de fragiliser le système : sans conservation, pas de possibilité d'attendre le bon prix ; on vend vite, et on vend mal.

\textbf{Les centres commerciaux modernes face aux mêmes maux.} Les supermarchés et centres commerciaux (ex. Mahima, Spar, Casino, Carrefour Market à Douala et Yaoundé) ne sont pas épargnés : ils subissent la rupture de stock due aux routes, la flambée du prix du carburant pour leurs groupes électrogènes (coupures d'électricité fréquentes), et la difficulté à s'approvisionner localement en produits frais de qualité constante, ce qui les pousse à importer davantage, accentuant la concurrence déloyale vis-à-vis des filières locales.

\courssub{Les pistes de solutions}

À chaque cause identifiée dans notre arbre doit répondre une solution concrète :
\begin{itemize}
\item \textbf{Réhabiliter et surtout entretenir les routes} : un entretien régulier coûte moins cher qu'une reconstruction et désenclave durablement les bassins de production (Programme d'Urgence Triennal, contrats d'entretien routier).
\item \textbf{Moderniser le rail} : rénovation de la Transcamerounaise et extension envisagée vers Kribi et le Nord pour un fret lourd et bon marché.
\item \textbf{Connecter les ports à l'intérieur} : achever les liaisons rail-route du port de Kribi vers son arrière-pays.
\item \textbf{Appuyer le secteur informel} : microcrédit adapté aux bayam-sellam et petits commerçants, coopératives d'achat groupé, réglementation allégée mais protectrice.
\item \textbf{Construire des infrastructures de stockage et de conservation} : entrepôts, marchés de gros modernes, unités de transformation (séchage, conserves) et chaîne du froid pour réduire les pertes post-récolte.
\item \textbf{Lutter contre les tracasseries} : réduction et mutualisation des contrôles, sanctions contre les péages informels.
\end{itemize}

\begin{exoresolu}[Étude de cas : pourquoi le chou coûte-t-il plus cher à Douala qu'à Dschang ?]
\textbf{Énoncé.} Un sac de choux est acheté \num{8000} FCFA au producteur à Dschang (Ouest). Le transport jusqu'à Douala coûte \num{2000} FCFA par sac, les frais de contrôles et péages informels sont estimés à \num{500} FCFA, et le grossiste puis le détaillant ajoutent chacun une marge de \num{20} \% sur leur prix d'achat. 1) Calcule le prix final payé par le consommateur à Douala. 2) Calcule le taux de hausse entre Dschang et Douala. 3) Identifie deux difficultés du commerce intérieur illustrées par cet exemple.
\tcblower
\textbf{Solution détaillée.}
1) Prix après transport et tracasseries :
\[ 8000 + 2000 + 500 = \num{10500} \text{ FCFA} \]
Marge du grossiste (\num{20} \% de \num{10500}) :
\[ \num{10500} \times \num{1,20} = \num{12600} \text{ FCFA} \]
Marge du détaillant (\num{20} \% de \num{12600}) :
\[ \num{12600} \times \num{1,20} = \num{15120} \text{ FCFA} \]
Le consommateur paie donc environ \num{15120} FCFA le sac.
2) Taux de hausse :
\[ \frac{\num{15120} - 8000}{8000} \times 100 = \num{89} \% \]
Le prix a presque doublé entre le champ et l'étal.
3) Deux difficultés illustrées : le \textbf{coût élevé du transport} (routes longues et dégradées entre l'Ouest et le Littoral) et la \textbf{multiplicité des intermédiaires} (grossiste puis détaillant, chacun prélevant sa marge), auxquelles s'ajoutent les \textbf{tracasseries routières} (\num{500} FCFA de frais informels).
\end{exoresolu}

\begin{exercice}[À toi de jouer]
Un sac d'oignons vaut \num{10000} FCFA à Maroua. Le transport jusqu'à Yaoundé coûte \num{3000} FCFA, les tracasseries \num{700} FCFA ; le grossiste ajoute \num{15} \% et le détaillant \num{25} \% sur leurs prix d'achat respectifs. Calcule le prix final et le taux de hausse, puis propose deux solutions qui permettraient de réduire cet écart.
\end{exercice}

\begin{corrige}[Exercice]
Prix après transport et tracasseries : $10000 + 3000 + 700 = \num{13700}$ FCFA. Après la marge du grossiste : $\num{13700} \times \num{1,15} = \num{15755}$ FCFA. Après la marge du détaillant : $\num{15755} \times \num{1,25} \approx \num{19694}$ FCFA. Taux de hausse : $\dfrac{\num{19694} - 10000}{10000} \times 100 \approx \num{97} \%$, soit un quasi-doublement. Solutions possibles : moderniser le rail Douala--Ngaoundéré et le prolonger vers le Nord pour un fret moins cher ; construire des entrepôts et marchés de gros pour réduire le nombre d'intermédiaires ; supprimer les péages informels.
\end{corrige}

\begin{aretenir}[L'essentiel de la section]
Les difficultés du commerce intérieur camerounais forment un système : infrastructures déficientes (routes dégradées, rail vétuste, ports mal connectés, aérien en crise), coût élevé du transport et pertes post-récolte, multiplicité des intermédiaires qui gonfle les prix, tracasseries routières, faiblesse du financement et de la conservation, insécurité et concurrence des importations. Conséquence majeure : le prix d'une denrée peut doubler ou tripler entre le champ et la ville, ce qui appauvrit le producteur et pèse sur le consommateur. Les solutions passent par l'entretien des routes, la modernisation du rail, l'appui au secteur informel et la construction d'infrastructures de stockage.
\end{aretenir}

\coursec{TP 4 : Carte des flux commerciaux inter-régionaux au Cameroun}

Dans la section précédente, nous avons dressé le bilan des difficultés qui freinent le commerce intérieur camerounais : routes dégradées, chemin de fer vétuste, enclavement de certaines régions, intermédiaires nombreux. Ces difficultés ne sont pas abstraites : elles se \emph{voient} sur une carte. Une route coupée, c'est un flux qui ralentit ; une région enclavée, c'est une flèche qui manque. Ce travail pratique te propose donc de passer de l'analyse à la \cle{représentation cartographique} : construire, pas à pas, une carte des flux commerciaux inter-régionaux du Cameroun. C'est l'une des compétences les plus attendues au Baccalauréat, et l'une des plus formatrices : celui qui sait cartographier les échanges comprend d'un seul regard la géographie économique de son pays.

\courssub{Objectif du TP et démarche générale}

\begin{definition}[Objectif du TP 4]
Représenter sur un fond de carte du Cameroun les \cle{produits} échangés entre les régions, leurs \cle{zones de production}, les \cle{zones de commercialisation et de consommation}, ainsi que les \cle{flux} commerciaux qui relient ces espaces, puis interpréter la carte obtenue pour dégager les grands axes et les déséquilibres spatiaux du commerce intérieur.
\end{definition}

L'intuition est simple : une carte de flux est une \emph{photographie synthétique} de la circulation des marchandises. Là où les flèches sont épaisses et nombreuses, l'intégration nationale fonctionne ; là où elles sont fines ou absentes, elle fait défaut. La démarche se déroule en cinq étapes successives, que nous allons détailler : inventorier et classer les produits, localiser les espaces, tracer les flux, rédiger la légende, interpréter la carte.

\begin{methode}[Matériel et documents nécessaires]
\begin{enumerate}
\item Un \textbf{fond de carte vierge du Cameroun} (limites extérieures et, si possible, limites régionales), à l'échelle d'une page.
\item Les \textbf{documents fournis} : tableaux statistiques de production par région, carte des axes de transport, photographies de marchés, liste des produits échangés.
\item Du \textbf{matériel de cartographie} : crayons de couleur, feutres fins de plusieurs épaisseurs (ou règle pour graduer les flèches), gomme, règle.
\item Une \textbf{feuille de brouillon} pour tester les figurés avant de les reporter sur la carte définitive.
\end{enumerate}
\end{methode}

\courssub{Étape 1 : inventorier et classer les produits échangés}

Avant de dessiner quoi que ce soit, il faut savoir \emph{ce qui circule}. À partir des documents fournis (tableaux de production, descriptions de marchés), tu dresses l'inventaire des produits échangés entre les régions, puis tu les \textbf{classes par grandes catégories}. Ce classement est indispensable : il déterminera le nombre de figurés de ta légende.

\begin{exemplebox}[Inventaire et classement des produits]
\begin{center}
\begin{tabularx}{\linewidth}{|l|X|X|}
\hline
\rowcolor{popblueL} \textbf{Catégorie} & \textbf{Produits} & \textbf{Principales zones de production} \\
\hline
Produits vivriers & Maïs, mil, sorgho, manioc, plantain, igname, haricot, pomme de terre & Ouest et Nord-Ouest (maïs, pomme de terre), Extrême-Nord (mil, sorgho), Centre et Sud (manioc, plantain) \\
\hline
Élevage et pêche & Bœufs, moutons, chèvres, volailles, poissons (frais, fumés, séchés) & Adamaoua et Nord (bovins), côte et fleuves (pêche maritime), lac Tchad et Lagdo (pêche continentale) \\
\hline
Cultures de rente & Cacao, café, coton, banane, huile de palme & Centre, Sud et Littoral (cacao), Ouest (café arabica), Nord (coton) \\
\hline
Produits industriels & Boissons, ciment, savons, textiles, produits pétroliers raffinés & Douala et Yaoundé (grandes zones industrielles), Edéa, Garoua \\
\hline
\end{tabularx}
\end{center}
\end{exemplebox}

\begin{attention}[Piège : trop de figurés tue la carte]
Ne représente pas chaque produit individuellement : vingt symboles différents rendent la carte illisible. Regroupe par catégorie (4 à 6 figurés maximum) et choisis des symboles \emph{évocateurs} : un épi pour les vivres, une tête de bœuf pour l'élevage, un poisson pour la pêche, une roue dentée ou une usine pour l'industrie.
\end{attention}

\courssub{Étape 2 : localiser les zones de production et de consommation}

Deuxième étape : placer sur le fond de carte les \textbf{espaces} entre lesquels circulent les produits.

\begin{itemize}
\item Les \cle{zones de production} sont figurées par des symboles disposés dans les régions concernées (par exemple trois épis dans l'Ouest pour le maïs, une tête de bœuf sur l'Adamaoua). Un même symbole peut être répété plusieurs fois si la production est dispersée.
\item Les \cle{zones de commercialisation et de consommation} sont les grandes villes : \textbf{Douala} (plus de 3 millions d'habitants, premier marché du pays), \textbf{Yaoundé} (capitale, environ 4 millions d'habitants), \textbf{Garoua}, \textbf{Maroua}, \textbf{Bamenda}, \textbf{Bafoussam}. On les représente par des points nommés, dont la taille peut traduire l'importance de la population.
\end{itemize}

La règle d'or : chaque point, chaque symbole posé sur la carte devra être \emph{expliqué dans la légende}. Note au brouillon, au fur et à mesure, la liste de tout ce que tu figures.

\courssub{Étape 3 : tracer les flux}

C'est le cœur du TP. Un \cle{flux} commercial se représente par une \textbf{flèche orientée} (du lieu de production vers le lieu de consommation) dont \textbf{l'épaisseur est proportionnelle à l'importance de l'échange}.

\begin{methode}[Représenter des flux lisibles]
\begin{enumerate}
\item \textbf{Orienter} chaque flèche : la pointe indique la destination (la ville consommatrice). Une flèche sans pointe est une erreur éliminatoire.
\item \textbf{Graduer les épaisseurs} : choisis deux ou trois épaisseurs (par exemple fine, moyenne, épaisse) correspondant à des flux faibles, moyens et forts, et annonce cette échelle dans la légende.
\item \textbf{Suivre les axes réels} : fais passer les flèches le long des routes et voies ferrées effectivement empruntées (axe Douala--Yaoundé, chemin de fer Transcamerounais vers Ngaoundéré, route nationale n\textsuperscript{o}1 vers Garoua et Maroua).
\item \textbf{Éviter les croisements illisibles} : si deux flux se croisent, décale légèrement l'un des tracés ou utilise des couleurs différentes par type de produit.
\end{enumerate}
\end{methode}

\begin{exemplebox}[Les axes majeurs du commerce intérieur]
Les flux les plus intenses du pays se concentrent sur deux grands axes : l'axe \textbf{Douala--Yaoundé} (environ 230 km, parcouru par des centaines de camions et de cars chaque jour, qui alimente la capitale en vivres, produits importés et manufacturés) et l'axe \textbf{Yaoundé--Ngaoundéré} (environ 800 km, dont la partie ferroviaire, le Transcamerounais, achemine vers le Sud bois, coton, bétail sur pied et produits du Grand Nord). À l'inverse, les flux entre l'Extrême-Nord et le Nord-Ouest, ou entre l'Est et le reste du pays, restent ténus : ce sont les régions enclavées.
\end{exemplebox}

\courssub{Étape 4 : rédiger une légende complète et ordonnée}

\begin{methode}[Légender une carte : la règle de correspondance]
\begin{enumerate}
\item \textbf{Correspondance stricte} : tout élément figuré sur la carte doit apparaître dans la légende, et \emph{réciproquement} tout élément de la légende doit figurer sur la carte. Aucun symbole « orphelin », aucune ligne de légende inutile.
\item \textbf{Titre précis} : il répond à trois questions — \emph{quoi} (les flux commerciaux inter-régionaux), \emph{où} (au Cameroun), \emph{quand} (au début des années 2020). Exemple : « Les principaux flux commerciaux inter-régionaux au Cameroun ».
\item \textbf{Ordre logique} : on présente d'abord les zones de production (figurés par produit), puis les villes (points), puis les flux (flèches graduées), enfin les éléments d'orientation (nord, échelle).
\item \textbf{Éléments techniques} : flèche du nord et échelle indicative (par exemple 1 cm pour 100 km) complètent la légende.
\end{enumerate}
\end{methode}

\begin{attention}[Erreur fréquente au Bac]
Dessiner des flèches \textbf{sans orientation} (simples traits), \textbf{sans proportionnalité} (toutes de même épaisseur alors que les flux sont inégaux) et \textbf{sans mention dans la légende} : voilà les trois fautes qui coûtent le plus de points dans un croquis de flux. Relis ta carte en te posant la question : un correcteur qui n'a jamais vu le Cameroun peut-il comprendre ma carte avec ma seule légende ?
\end{attention}

\courssub{Modèle de carte à reproduire}

La figure ci-dessous présente un modèle simplifié de ce que tu dois produire. Le fond de carte est celui d'une carte générale ; l'essentiel est la cohérence entre les figurés, les flèches et la légende.

\begin{popfigure}
\begin{center}
\includegraphics[width=\linewidth,height=9.5cm,keepaspectratio]{N01/cmr_map.jpg}
\end{center}
\legende{Fond de carte du Cameroun sur lequel construire la carte des flux (régions, villes, routes et voie ferrée sont déjà dessinées). La carte modèle situe : 1. les zones de production par des figurés (vivres à l'Ouest, bétail au Grand Nord, cacao au Centre--Sud, pêche sur le littoral) ; 2. les grandes villes de commercialisation et de consommation par des points (Douala, Yaoundé) ; 3. les flux par des flèches orange d'épaisseur proportionnelle à leur intensité : flux le plus épais Douala--Yaoundé, flux moyens Yaoundé--Ngaoundéré et Ouest vers Douala/Yaoundé, flux fin vers le Nord ; 4. une légende ordonnée. Ces éléments sont à reporter sur le fond par l'élève. \\ Source : Wikimedia Commons, CIA, domaine public.}
\end{popfigure}

\begin{popfigure}
\begin{tikzpicture}[every node/.style={font=\small}]
\draw[thick, popdark, fill=popcream] (0,0) rectangle (8.6,6.4);
\node[font=\bfseries, anchor=north] at (4.3,6.3) {Légende : les flux commerciaux inter-régionaux au Cameroun};
\draw[popline] (0.3,5.7) -- (8.3,5.7);
% Zones de production
\node[anchor=west, font=\bfseries\footnotesize] at (0.4,5.3) {Zones de production};
\node[popgreen, anchor=west, font=\footnotesize] at (0.6,4.9) {\textbf{Vivres} : maïs, plantain, manioc, mil};
\node[poppurple, anchor=west, font=\footnotesize] at (0.6,4.5) {\textbf{Bétail} : bovins, ovins, caprins};
\node[popblue, anchor=west, font=\footnotesize] at (0.6,4.1) {\textbf{Cacao} et cultures de rente};
\node[poporange, anchor=west, font=\footnotesize] at (0.6,3.7) {\textbf{Pêche} : maritime et continentale};
% Villes
\node[anchor=west, font=\bfseries\footnotesize] at (0.4,3.2) {Villes de commercialisation et de consommation};
\fill[popink] (0.8,2.75) circle (2.5pt); \node[anchor=west, font=\footnotesize] at (1.1,2.75) {Grande ville (Douala, Yaoundé)};
\fill[popink] (0.8,2.35) circle (1.8pt); \node[anchor=west, font=\footnotesize] at (1.1,2.35) {Ville régionale (Garoua, Maroua, Bamenda, Bafoussam)};
% Flux
\node[anchor=west, font=\bfseries\footnotesize] at (0.4,1.85) {Flux commerciaux};
\draw[line width=4pt, ->, poporange] (0.7,1.4) -- (1.9,1.4); \node[anchor=west, font=\footnotesize] at (2.1,1.4) {Flux intense (axe Douala--Yaoundé)};
\draw[line width=2.5pt, ->, poporange] (0.7,1.0) -- (1.9,1.0); \node[anchor=west, font=\footnotesize] at (2.1,1.0) {Flux moyen (vivres de l'Ouest, Transcamerounais)};
\draw[line width=1.2pt, ->, poporange] (0.7,0.6) -- (1.9,0.6); \node[anchor=west, font=\footnotesize] at (2.1,0.6) {Flux faible (régions enclavées)};
\end{tikzpicture}
\legende{Exemple de légende normalisée : titre précis (quoi, où), cadre, rubriques ordonnées (production, villes, flux), symboles identiques à ceux de la carte. L'orientation (flèche du nord) et l'échelle figurent sur la carte elle-même.}
\end{popfigure}

\courssub{Étape 5 : interpréter la carte et conclure}

Une carte n'est pas une fin en soi : elle sert à \emph{raisonner}. L'interprétation s'organise en trois questions.

\begin{enumerate}
\item \textbf{Quels sont les axes majeurs ?} La carte fait apparaître la dominante écrasante de l'axe Douala--Yaoundé, prolongée vers Ngaoundéré par le Transcamerounais : c'est la « colonne vertébrale » du commerce intérieur, qui draine vivres, bétail, bois et produits manufacturés.
\item \textbf{Quelles régions dominent ?} Le triangle Douala--Yaoundé--Bafoussam concentre la production vivrière commerciale, l'industrie et la consommation : c'est le cœur économique du pays. Le Grand Nord fournit bétail et coton mais reste dépendant du Sud pour les produits manufacturés.
\item \textbf{Quelles régions sont enclavées ?} L'Est et une partie de l'Adamaoua n'envoient que des flux ténus : mauvais état des pistes, éloignement, faible densité. La carte visualise ainsi directement les difficultés étudiées dans la section précédente.
\end{enumerate}

\begin{exoresolu}[Rédiger la conclusion du TP]
Énoncé : à partir de ta carte des flux commerciaux inter-régionaux, rédige une conclusion de huit à douze lignes dégageant l'essentiel de l'organisation du commerce intérieur camerounais.
\tcblower
Solution détaillée. La carte des flux commerciaux inter-régionaux révèle un commerce intérieur \textbf{fortement polarisé}. L'essentiel des échanges s'organise autour de deux métropoles, Douala (port et capitale économique) et Yaoundé (capitale politique et premier marché de consommation de l'intérieur), reliées par l'axe routier le plus chargé du pays. Autour de ce duo, trois grands flux structurent l'espace national : les flux vivriers descendant de l'Ouest et du Nord-Ouest vers les deux métropoles ; les flux de bétail et de coton venant du Grand Nord par la route nationale n\textsuperscript{o}1 et le Transcamerounais jusqu'à Ngaoundéré puis Yaoundé ; et, en sens inverse, les flux de produits manufacturés et importés remontant de Douala vers toutes les régions. Cette organisation traduit une réelle \textbf{interdépendance} entre régions complémentaires, mais elle révèle aussi de profonds \textbf{déséquilibres} : l'Est et l'Extrême-Nord restent à l'écart des grands flux, faute d'infrastructures entretenues. Améliorer les routes et le chemin de fer apparaît ainsi comme la condition première d'une meilleure intégration nationale par les échanges.
\end{exoresolu}

\courssub{Variante : croquis simplifié ou diagramme de flux}

Si le temps manque (une heure au lieu de deux), deux solutions allégées permettent d'atteindre l'objectif.

\begin{exemplebox}[Deux variantes rapides]
\begin{itemize}
\item \textbf{Le croquis simplifié} : on ne représente que trois catégories de produits (vivres, bétail, produits manufacturés), quatre villes (Douala, Yaoundé, Garoua, Bamenda) et trois flèches graduées. La légende tient en cinq lignes, mais la règle de correspondance carte--légende reste intégralement respectée.
\item \textbf{Le diagramme de flux} : on abandonne le fond de carte pour un schéma où les régions sont des rectangles disposés approximativement selon leur position réelle (Nord en haut, Sud en bas), reliés par des flèches épaisses ou fines portant le nom du produit. On perd la précision spatiale, mais on conserve l'essentiel : l'origine, la destination et l'intensité des échanges.
\end{itemize}
\end{exemplebox}

\begin{savaistu}[La carte, un outil de diagnostic]
Une bonne carte de flux permet de visualiser d'un seul coup d'œil les déséquilibres spatiaux du commerce intérieur camerounais : les planificateurs de l'État s'en servent pour décider où bitumer une route ou réhabiliter une voie ferrée. Quand tu traces une flèche fine vers l'Est, tu montres exactement le même problème que les ingénieurs du ministère des Travaux publics : l'enclavement. Cartographier, c'est déjà diagnostiquer — et diagnostiquer, c'est préparer la décision.
\end{savaistu}

\begin{exercice}[À toi de jouer]
Sur un fond de carte vierge du Cameroun, construis la carte complète des flux commerciaux inter-régionaux en suivant les cinq étapes du TP. Contraintes : au moins quatre figurés de production, six villes nommées, trois épaisseurs de flèches, une légende ordonnée avec titre, nord et échelle. Rédige ensuite un commentaire de dix lignes intitulé « Les flux commerciaux révèlent-ils un Cameroun intégré ? ».
\end{exercice}

\begin{corrige}[Exercice n\textsuperscript{o}1]
Éléments de correction. La carte doit présenter : figurés vivriers dans l'Ouest, le Nord-Ouest et le Centre-Sud ; figuré bétail sur l'Adamaoua et le Nord ; figuré pêche sur la côte et le lac Tchad ; figuré industriel à Douala et Yaoundé. Les flèches les plus épaisses relient Douala à Yaoundé et l'Ouest à ces deux villes ; des flèches moyennes suivent le Transcamerounais vers Ngaoundéré ; des flèches fines relient Garoua, Maroua et l'Est. La légende respecte la correspondance stricte et l'ordre production--villes--flux. Le commentaire doit nuancer : oui, les régions sont interdépendantes (vivres, bétail, produits manufacturés circulent dans tout le pays), mais l'intégration reste inachevée car les flux sont polarisés par Douala et Yaoundé et parce que l'Est et le Grand Nord demeurent mal desservis. La mention des infrastructures (routes bitumées en progrès mais entretien irrégulier, chemin de fer vétuste) est attendue pour expliquer ces contrastes.
\end{corrige}

\newpage

\coursec{Activité d'intégration}

\coursec{Les échanges inter-régionaux au Cameroun}

\courssub{Introduction : notions clés et situation-problème}

\begin{aretenir}[Objectifs de la leçon]
À l'issue de cette leçon, tu seras capable de :
\begin{itemize}
\item \cle{observer}, \cle{localiser} et \cle{décrire} les flux commerciaux inter-régionaux à partir de documents variés (tableaux statistiques, cartes, photographies) ;
\item \cle{inventorier} et \cle{classer} les produits selon leurs régions de production ;
\item \cle{construire} et \cle{légender} une carte des flux commerciaux (flèches proportionnelles aux tonnages) ;
\item \cle{interpréter} l'interdépendance entre les régions du Cameroun et expliquer les difficultés du commerce intérieur ;
\item formuler des \cle{solutions} argumentées pour améliorer le ravitaillement des villes et l'intégration nationale.
\end{itemize}
\end{aretenir}

\textbf{Situation-problème : Le ravitaillement de Yaoundé, enquête sur les flux vivriers.}\par
Le marché Mokolo, situé au cœur de Yaoundé, est l'un des plus grands marchés vivriers de la capitale politique du Cameroun (plus de 4 millions d'habitants, ordre de grandeur 2024). Chaque semaine, des dizaines de camions y déchargent des produits agricoles venus de plusieurs régions du pays. Ces marchandises sont achetées dans les zones de production par des revendeurs de vivres appelés \cle{bayam selam} (de l'anglais \emph{buy them sell them} : « acheter pour revendre »), qui louent des camions, chargent les produits dans les marchés périodiques de gros (Bafoussam, Foumbot, Bafia, Abong-Mbang\ldots) et les revendent aux détaillantes de Mokolo.

Une enquête fictive mais réaliste, menée auprès des grossistes du marché Mokolo pendant une semaine du mois de mars, a permis d'établir le tableau suivant (les tonnages sont des ordres de grandeur plausibles, donnés à titre pédagogique).

\begin{center}
\begin{tabularx}{\linewidth}{|l|X|c|c|}
\hline
\rowcolor{popblueL} \textbf{Produit} & \textbf{Région d'origine (zone de production)} & \textbf{Tonnage hebdomadaire (t)} & \textbf{Axe principal emprunté} \\
\hline
Pommes de terre & Ouest (Bafoussam, Dschang) & 120 & Route N4 (Bafoussam--Yaoundé) \\
\hline
Carottes, choux, poireaux & Ouest (hauts plateaux) & 60 & Route N4 \\
\hline
Maïs, haricots & Ouest (Foumbot, Mbouda) & 80 & Route N4 \\
\hline
Plantains, macabos & Centre (Bafia, Obala) et Sud & 150 & Routes N1 et N2 \\
\hline
Manioc, bâtons de manioc & Centre et Sud (Mbalmayo, Sangmélima) & 90 & Routes N1 et N2 \\
\hline
Tomates, oignons & Nord (Garoua, Maroua) et Adamaoua & 70 & Chemin de fer (Transcam) puis route \\
\hline
B\oe ufs sur pied (équivalent carcasse) & Adamaoua (Ngaoundéré) et Nord & 40 & Transcam (Ngaoundéré--Yaoundé) \\
\hline
Poissons fumés & Littoral (Youpwé, Limbé) et Sud & 30 & Routes N3 et N2 \\
\hline
Huile de palme, noix de palme & Sud-Ouest (Muyuka, Kumba) et Littoral & 50 & Route N3 puis N1 \\
\hline
\end{tabularx}
\end{center}

\textbf{Problème posé :} Comment Yaoundé, ville de plus de 4 millions d'habitants, parvient-elle à se ravitailler chaque semaine, et que révèle ce ravitaillement de l'intégration nationale ?

\courssub{Les facteurs de développement des échanges : les infrastructures}

\begin{definition}[Infrastructures de transport]
Ensemble des équipements fixes (routes, voies ferrées, ports, aéroports) et mobiles (camions, wagons, navires, avions) qui permettent la circulation des biens et des personnes. Au Cameroun, l'essor des échanges repose sur la densité et la qualité de ce réseau.
\end{definition}

\begin{propriete}[État du réseau camerounais (diagnostic partagé)]
\begin{itemize}
\item \textbf{Routes bitumées :} En augmentation (programmes d'urgence, routes à 2x2 voies Yaoundé--Douala, Yaoundé--Nsimalen), mais l'entretien est irrégulier (nids-de-poule, éboulements, chaussées déformées), surtout en saison des pluies.
\item \textbf{Chemin de fer (Transcam) :} Ligne unique Yaoundé--Ngaoundéré (environ 622 km), vétuste, vitesse moyenne faible (30--40 km/h), pannes fréquentes, capacité limitée.
\item \textbf{Ports :} Douala (premier port d'Afrique centrale, 95\% du trafic maritime) et Kribi (port en eau profonde, phase 1 opérationnelle). Problème : connexion ferroviaire/routière insuffisante avec l'intérieur (corridor Douala--Ngaoundéré saturé).
\item \textbf{Aéroports :} Douala et Yaoundé-Nsimalen (internationaux), aéroports secondaires (Garoua, Maroua, Bafoussam). Compagnie nationale (Camair-Co) en difficultés financières structurelles, fret aérien marginal pour les vivriers.
\end{itemize}
\end{propriete}

\begin{experience}[Analyse de photographies : l'état des infrastructures]
\textbf{Document 3 (fourni dans le TP) :} Photographie montrant des bayam selam chargeant des sacs de pommes de terre sur un camion au marché de gros de Bafoussam, sous la pluie, sur une chaussée dégradée (ornières, boue, absence de quai de chargement).
\textbf{Interprétation :} L'absence d'infrastructures logistiques adaptées (aires de stationnement, hangars de stockage, quais) aux marchés de gros allonge les temps de chargement, abîme les produits et les véhicules, et augmente les coûts.
\end{experience}

\courssub{L'interdépendance entre les régions}

\begin{definition}[Interdépendance régionale]
Relation de dépendance mutuelle entre des espaces géographiques qui s'échangent des biens, des services, des capitaux ou de la main-d'œuvre. Aucune région ne peut subvenir seule à ses besoins ; la spécialisation productive (climat, sol, histoire) crée la complémentarité.
\end{definition}

\begin{propriete}[Spécialisation et complémentarité au Cameroun]
\begin{itemize}
\item \textbf{Région de l'Ouest (Hauts plateaux) :} Climat tempéré d'altitude $\rightarrow$ cultures maraîchères (pommes de terre, carottes, choux, poireaux), céréales (maïs, haricots). \emph{Rôle : grenier à légumes et céréales.}
\item \textbf{Régions du Centre et du Sud :} Climat équatorial humide $\rightarrow$ cultures de base (plantain, macabo, manioc, igname), produits forestiers. \emph{Rôle : bassin vivrier de base.}
\item \textbf{Grand Nord (Adamaoua, Nord, Extrême-Nord) :} Climat soudanien/sahélien $\rightarrow$ élevage bovin (bœufs sur pied), cultures de rente (coton, oignon, tomate de saison sèche). \emph{Rôle : bassin de protéines animales et légumes de contre-saison.}
\item \textbf{Régions du Littoral et du Sud-Ouest :} Climat équatorial littoral $\rightarrow$ cultures industrielles (palmier à huile, bananier, hévéa), pêche artisanale et industrielle. \emph{Rôle : huiles, poissons, fruits d'exportation et consommation locale.}
\end{itemize}
\end{propriete}

\begin{savaistu}[L'interdépendance en chiffres : le cas de Mokolo]
D'après le tableau de l'enquête (Document 1), le tonnage total hebdomadaire est de \textbf{690 t}. La répartition par grande région d'origine montre une dépendance quasi équilibrée :
\begin{itemize}
\item Ouest : 260 t (37,7 \%) -- Légumes, pommes de terre, céréales.
\item Centre--Sud : 255 t (37,0 \%) -- Plantain, manioc, poisson (part Sud).
\item Grand Nord : 110 t (15,9 \%) -- Viande, oignons, tomates.
\item Littoral--Sud-Ouest : 65 t (9,4 \%) -- Huile de palme, poisson (part Littoral).
\end{itemize}
Yaoundé ne produit presque rien de ce qu'elle consomme ; les régions productrices ont besoin du pouvoir d'achat yaoundéen. C'est l'interdépendance \emph{mesurée}.
\end{savaistu}

\courssub{L'organisation du commerce intérieur}

\begin{definition}[Marché périodique]
Marché qui se tient à jours fixes (ex : tous les 8 jours -- « jour de marché » -- ou hebdomadaire) dans une localité rurale ou urbaine secondaire. C'est le lieu de rencontre entre producteurs locaux (paysans) et collecteurs (bayam selam). Exemples : Foumbot (maïs), Bafia (plantain), Magba (bétail).
\end{definition}

\begin{definition}[Bayam selam (collecteurs-revendeurs)]
Acteurs charnières du commerce intérieur. Ils assurent la \cle{collecte} à la ferme ou au marché périodique, le \cle{transport} (location de camions), le \cle{stockage} temporaire et la \cle{vente en gros} aux détaillants des grands marchés urbains (Mokolo, Nkolndongo, Mfoundi à Yaoundé ; Congo, Bépanda à Douala). Ils supportent les risques (pertes, vol, variation des prix).
\end{definition}

\begin{definition}[Vendeurs à la sauvette]
Commerçants ambulants non sédentarisés, installés sur les trottoirs, abords de marchés ou carrefours. Ils écoulent de petits volumes (détail de proximité), souvent des produits frais ou manufacturés importés. Ils témoignent de la vitalité de l'informel et de l'insuffisance des infrastructures marchandes formelles.
\end{definition}

\begin{definition}[Centre commercial]
Grande surface moderne, climatisée, offrant une large gamme de produits (alimentaires, non alimentaires) avec prix fixes, hygiène contrôlée, parking. Exemples : \emph{Super U} (Douala, Yaoundé), \emph{Casino} (Yaoundé), \emph{Mahima} (Douala). Ils captent une clientèle aisée mais restent marginaux dans le volume global du commerce vivrier (< 5 \% du ravitaillement vivrier).
\end{definition}

\begin{methode}[Analyser l'organisation du commerce intérieur]
\begin{enumerate}
\item Identifier les \textbf{lieux} : marchés de gros ruraux (collecte), marchés de gros urbains (répartition), marchés de détail (consommation).
\item Identifier les \textbf{acteurs} : producteurs $\rightarrow$ bayam selam (collecteurs/grossistes) $\rightarrow$ détaillants (boutiquiers, vendeuses assises) $\rightarrow$ consommateurs.
\item Identifier les \textbf{flux} : produits, quantités, axes de transport, coûts, délais.
\item Repérer les \textbf{dysfonctionnements} : pertes, racket, absence de chaîne du froid, informel dominant.
\end{enumerate}
\end{methode}

\courssub{Les difficultés du commerce intérieur}

\begin{aretenir}[Principales difficultés (à connaître pour le Bac)]
\begin{enumerate}
\item \textbf{Mauvais état des infrastructures routières :} Dégradation rapide (climat, surcharge des camions), entretien curatif et non préventif, coupures saisonnières (ponts effondrés, éboulements sur N4, N6, N10).
\item \textbf{Coût et lenteur des transports :} Prix du carburant, péages, usure des véhicules, vitesse moyenne 40--50 km/h sur route, 30 km/h sur rail. Le transport représente 30 à 50 \% du prix final à Mokolo.
\item \textbf{Racket et contrôles multiples :} Postes de police, gendarmerie, douane, conseil régional, délégués de groupements \ldots jusqu'à 10--15 barrières sur Yaoundé--Bafoussam (270 km). Temps perdu et argent prélevé.
\item \textbf{Pertes post-récolte et absence de chaîne du froid :} 30 à 40 \% des productions maraîchères (tomates, carottes) et halieutiques perdues entre le champ et l'étal (Document 3 : chargement sous la pluie, sacs non palettisés).
\item \textbf{Insécurité :} Coupeurs de route (Grand Nord, Est), vols à la tire dans les marchés.
\item \textbf{Faible modernisation des circuits :} Dominance de l'informel (bayam selam, vendeurs à la sauvette), absence de marchés de gros modernes avec chambres froides, standardisation des conditionnements (cagettes vs sacs de 100 kg).
\end{enumerate}
\end{aretenir}

\courssub{TP 4 : Carte des flux commerciaux inter-régionaux au Cameroun}

\courssub{Consignes du Travail Pratique}

\textbf{Durée :} 2 heures (1 h 30 de travail de groupe / individuel, 30 min de restitution et correction).\par
\textbf{Matériel :} Fond de carte du Cameroun (contour, limites régionales, villes principales, axes N1, N2, N3, N4, voie ferrée), crayons de couleur, règle, calculatrice, documents 1 à 4.

\begin{enumerate}
\item \textbf{Exploitation du Document 1 (Tableau statistique) :}
\begin{enumerate}
\item Inventorie les produits vendus à Mokolo et classe-les par région d'origine (Ouest ; Centre--Sud ; Grand Nord : Adamaoua, Nord, Extrême-Nord ; Littoral--Sud-Ouest).
\item Calcule le tonnage total hebdomadaire fourni par chaque grande région.
\item Calcule la part en pourcentage de chaque région dans le total. Arrondis au dixième de pourcent.
\end{enumerate}
\item \textbf{Construction de la carte des flux (Cœur du TP) :}
Sur le fond de carte fourni :
\begin{enumerate}
\item \textbf{Zones de production :} Colorie chaque grande région d'origine avec une couleur distincte (ex : Vert = Ouest ; Jaune = Centre--Sud ; Orange = Grand Nord ; Bleu = Littoral--Sud-Ouest). Légende ces couleurs.
\item \textbf{Zone de consommation :} Repère Yaoundé par un symbole fort (étoile rouge, cercle plein noir) et nomme-le « Yaoundé (Marché Mokolo) ».
\item \textbf{Flux :} Trace \textbf{quatre flèches} convergentes vers Yaoundé, partant du centre approximatif de chaque zone de production colorée.
\begin{itemize}
\item L'\textbf{épaisseur} de chaque flèche doit être \textbf{proportionnelle au tonnage} total de la région (échelle suggérée : 1 mm d'épaisseur pour 50 t).
\item L'\textbf{origine} de la flèche est la zone de production, sa \textbf{pointe} touche Yaoundé.
\item Pour le Grand Nord, indique le parcours combiné (rail + route) par une flèche en pointillés ou une annotation.
\end{itemize}
\item \textbf{Éléments obligatoires de la carte :} Titre précis, Légende organisée (Zones de production, Zone de consommation, Flux + échelle des largeurs), Orientation (flèche Nord), Échelle indicative (ex : 1 cm $\approx$ 100 km), Auteur/Date.
\end{enumerate}
\item \textbf{Analyse des Documents 2, 3 et 4 :}
\begin{enumerate}
\item À l'aide du Document 2 (carte réseau), nomme les axes routiers (N1, N2, N3, N4) et la ligne ferroviaire (Transcam) empruntés par les flux identifiés en question 1.
\item À l'aide du Document 3 (photo chargement Bafoussam), décris l'état de la voirie et les conditions de chargement. Quelles en sont les conséquences sur les coûts et les pertes ?
\item À l'aide du Document 4 (photo Mokolo), décris l'organisation de l'espace marchand (étals formels vs vendeurs à la sauvette sur trottoirs). Quel problème d'urbanisme et de commerce cela pose-t-il ?
\end{enumerate}
\item \textbf{Synthèse rédactionnelle (15--20 lignes) :}
Rédige un texte intitulé : \textbf{« Mokolo, miroir de l'interdépendance régionale »}.
Structure attendue :
\begin{itemize}
\item (a) Comment l'interdépendance entre les régions se traduit concrètement dans le ravitaillement de Yaoundé (données chiffrées à l'appui).
\item (b) Les difficultés du commerce intérieur mises en évidence par les documents (infrastructures, acteurs, pertes, urbanisme).
\item (c) \textbf{Deux solutions concrètes et argumentées} pour améliorer ces échanges (ex : réhabilitation axes + marchés de gros frigorifiques ; organisation bayam selam en coopératives + suppression barrières illégales).
\end{itemize}
\item \textbf{Restitution orale :} Présente ta carte en 5 minutes (titre, légende, flux principaux, difficultés, solutions). Participe à la correction collective.
\end{enumerate}

\courssub{Ressources à mobiliser (Rappels de cours)}

\begin{aretenir}[Savoirs indispensables pour le TP]
\begin{itemize}
\item \textbf{Facteurs de développement des échanges :} Infrastructures (routes bitumées en augmentation mais entretien irrégulier ; chemin de fer vétuste ; aéroports ; ports mal connectés à l'intérieur) et Interdépendance régionale (spécialisation climatique $\rightarrow$ complémentarité).
\item \textbf{Organisation du commerce intérieur :} Marchés périodiques (ruraux, collecte), Bayam selam (collecteurs-grossistes itinérants), Vendeurs à la sauvette (détail informel urbain), Centres commerciaux (formel moderne, minoritaire).
\item \textbf{Difficultés du commerce intérieur :} État routes/rail, coût/lenteur transport, racket/barrières, pertes post-récolte (absence chaîne du froid), insécurité, informel dominant.
\item \textbf{Savoir-faire cartographique (exigible au Bac) :} Une carte \textbf{doit} comporter un \textbf{TITRE}, une \textbf{LEGENDE} organisée (hiérarchisée), une \textbf{ORIENTATION} (flèche Nord), une \textbf{ECHELLE}. Les flux = flèches d'épaisseur proportionnelle aux quantités (largeur $\propto$ tonnage).
\item \textbf{Calculs :} Part $\% = \dfrac{\text{Partie}}{\text{Total}} \times 100$. Largeur flèche (mm) $= \dfrac{\text{Tonnage (t)}}{50}$ (si échelle 1 mm = 50 t).
\end{itemize}
\end{aretenir}

\courssub{Démarche et corrigé commenté du TP}

\begin{exoresolu}[Le ravitaillement du marché Mokolo : corrigé guidé complet]
Énoncé : Classer les produits par région, cartographier les flux vers Yaoundé, analyser l'interdépendance, les difficultés et proposer deux solutions.
\tcblower

\textbf{Étape 1 : Classement des produits par région et calcul des parts (Question 1).}\par
On regroupe les lignes du Document 1 par grande région d'origine. Pour les produits partagés (poissons fumés), on répartit le tonnage au prorata des zones citées (estimation pédagogique : 15 t Sud / 15 t Littoral).

\begin{center}
\begin{tabularx}{\linewidth}{|l|X|c|c|}
\hline
\rowcolor{popblueL} \textbf{Grande région} & \textbf{Produits concernés (Document 1)} & \textbf{Tonnage (t)} & \textbf{Part (\%)} \\
\hline
\textbf{Ouest} & Pommes de terre (120), Légumes hauts plateaux (60), Maïs/Haricots (80) & $120+60+80 = \mathbf{260}$ & $\frac{260}{690}\times100 \approx \mathbf{37,7}$ \\
\hline
\textbf{Centre--Sud} & Plantains/Macabos (150), Manioc (90), Poissons fumés Sud (15 est.) & $150+90+15 = \mathbf{255}$ & $\frac{255}{690}\times100 \approx \mathbf{37,0}$ \\
\hline
\textbf{Grand Nord} & Tomates/Oignons (70), Bœufs (40) & $70+40 = \mathbf{110}$ & $\frac{110}{690}\times100 \approx \mathbf{15,9}$ \\
\hline
\textbf{Littoral--Sud-Ouest} & Huile/Palme (50), Poissons fumés Littoral (15 est.) & $50+15 = \mathbf{65}$ & $\frac{65}{690}\times100 \approx \mathbf{9,4}$ \\
\hline
\textbf{TOTAL} & & \textbf{690} & \textbf{100,0} \\
\hline
\end{tabularx}
\end{center}

\textbf{Interprétation :} Yaoundé dépend de \textbf{quatre grands bassins} pour 690 t/semaine. L'Ouest et le Centre-Sud assurent les 3/4 de l'approvisionnement (produits de base + maraîchage). Le Grand Nord apporte les protéines animales et les légumes de contre-saison. Le Littoral/Sud-Ouest apporte l'huile et le poisson. L'interdépendance est \textbf{structurelle, massive et réciproque} (débouchés pour les producteurs, alimentation pour la capitale).

\textbf{Étape 2 : Construction de la carte des flux (Question 2).}\par
\emph{Choix cartographiques :}
\begin{itemize}
\item \textbf{Échelle des largeurs :} 1 mm = 50 t. $\rightarrow$ Ouest : 5,2 mm ; Centre-Sud : 5,1 mm ; Grand Nord : 2,2 mm ; Littoral-SO : 1,3 mm.
\item \textbf{Couleurs :} Vert (Ouest), Jaune (Centre-Sud), Orange (Grand Nord), Bleu (Littoral-SO).
\item \textbf{Flèche Grand Nord :} Trait mixte (pointillés pour le rail Ngaoundéré--Yaoundé, plein pour la route finale) ou flèche unique avec pictogramme train+camion dans la légende.
\item \textbf{Légende :} Organisée en 3 items : 1. Zones de production (carrés de couleur), 2. Zone de consommation (étoile), 3. Flux hebdomadaires (échantillons de flèches avec largeurs réelles + valeurs en tonnes).
\end{itemize}

\begin{popfigure}
\begin{center}
\includegraphics[width=\linewidth,height=11cm,keepaspectratio]{N09/cmr_admin.png}
\end{center}
\legende{Fond de carte des régions du Cameroun (chiffres : 1 Extrême-Nord, 2 Nord, 3 Adamaoua, 4 Nord-Ouest, 5 Sud-Ouest, 6 Ouest, 7 Centre, 8 Est, 9 Littoral, 10 Sud) pour la carte des flux vivriers vers le marché Mokolo (Yaoundé, région 7). La carte attendue colorie les zones de production : Ouest (260 t, région 6), Centre--Sud (255 t, régions 7 et 10), Grand Nord (110 t, régions 1 à 3) et Littoral--Sud-Ouest (65 t, régions 9 et 5) ; elle place une étoile rouge à Yaoundé et trace des flèches d'épaisseur proportionnelle aux tonnages (1 mm = 50 t), celle du Grand Nord combinant rail (Transcam) et route. \\ Source : Wikimedia Commons, TUBS, CC BY-SA 3.0. Tonnages : enquête fictive du Document 1.}
\end{popfigure}

\textbf{Étape 3 : Infrastructures et difficultés (Question 3).}\par
\begin{itemize}
\item \textbf{Axes utilisés (Doc 2) :} N4 (Bafoussam--Yaoundé) pour l'Ouest ; N1 (Douala--Yaoundé) et N2 (Yaoundé--Sangmélima) pour Centre-Sud et Littoral-SO ; Transcam (Ngaoundéré--Yaoundé) pour Grand Nord (rail + route finale).
\item \textbf{État et difficultés (Doc 2 \& 3) :} Routes bitumées mais dégradées (nids-de-poule, bas-côtés effondrés), voirie des marchés de gros inexistante (chargement dans la boue, Doc 3), chemin de fer vétuste (vitesse 30 km/h, retards fréquents, wagons insuffisants).
\item \textbf{Conséquences pour les bayam selam :} Temps de trajet doublé (Bafoussam--Yaoundé : 4--5 h au lieu de 3 h), casses mécaniques fréquentes, pertes de produits (écrasement, pourrissement sous la pluie), coûts de location camion majorés, racket aux barrières (police, gendarmerie, conseils régionaux) $\rightarrow$ renchérissement du prix à Mokolo.
\end{itemize}

\textbf{Étape 4 : Synthèse rédactionnelle (Question 4).}\par
\textbf{Mokolo, miroir de l'interdépendance régionale}\par
(a) Le marché Mokolo illustre de manière spectaculaire l'interdépendance nationale. Chaque semaine, 690 tonnes de vivriers convergent vers Yaoundé depuis les quatre coins du pays : l'Ouest fournit 37,7 \% du tonnage (pommes de terre, légumes, maïs), le Centre-Sud 37,0 \% (plantain, manioc), le Grand Nord 15,9 \% (viande, oignons, tomates) et le Littoral-Sud-Ouest 9,4 \% (huile, poisson). Yaoundé, ville administrative à faible production agricole, ne survit que grâce à cette solidarité territoriale ; inversement, les régions productrices trouvent dans la capitale leur principal débouché monétaire. (b) Ce flux vital est entravé par des difficultés structurelles : le réseau routier (N4, N1, N2), bien que bitumé, est mal entretenu (Doc 3 : chargement dans la boue à Bafoussam) ; le Transcam, seul lien ferré avec le Nord, est vétuste et lent ; les barrières de rackets multiplient les coûts et les délais ; l'absence de chaîne du froid provoque 30--40 \% de pertes sur les produits périssables ; enfin, l'étalement des vendeurs à la sauvette sur les trottoirs (Doc 4) révèle l'insuffisance des infrastructures marchandes urbaines. (c) Deux solutions prioritaires : (1) La \textbf{réhabilitation lourde et l'entretien programmé des axes structurants} (N4, Transcam) couplée à la construction de \textbf{marchés de gros modernes équipés de chambres froides} à Yaoundé et dans les pôles de production (Bafoussam, Ngaoundéré) pour réduire les pertes et fluidifier la logistique. (2) L'\textbf{organisation professionnelle des bayam selam en coopératives ou groupements d'intérêt économique} pour mutualiser le transport (camions pleins, assurance), négocier la suppression des barrières illégales avec l'État, et accéder au financement bancaire pour moderniser leurs outils de travail.

\begin{attention}[Pièges fréquents à éviter au Bac]
\begin{itemize}
\item \textbf{Confusion Production / Consommation :} La flèche part \textbf{toujours} de la zone de production (origine) vers la zone de consommation (destination). Ne trace pas de flèche « retour » (argent, biens manufacturés) sauf si le sujet le demande explicitement.
\item \textbf{Légende incomplète :} Une carte sans titre, sans légende, sans nord, sans échelle = note éliminatoire. La légende doit expliquer \textbf{chaque} symbole (couleur des zones, symbole ville, épaisseur des flèches avec l'échelle numérique).
\item \textbf{Proportionnalité approximative :} « Proportionnel » ne veut pas dire « gros, moyen, petit » au feeling. Il faut une \textbf{échelle numérique} (ex: 1 mm = 50 t) et mesurer les largeurs à la règle.
\item \textbf{Vocabulaire :} Ne dis pas « le Nord envoie des bœufs ». Dis : « La région de l'Adamaoua (bassin d'élevage) approvisionne Yaoundé en bœufs sur pied via le Transcam ». Précise les \textbf{noms de régions, villes, axes}.
\end{itemize}
\end{attention}

\courssub{Bilan de la leçon et du TP 4}

Cette leçon et le TP 4 ont permis de mobiliser l'ensemble des compétences du programme : lecture critique de tableau statistique, calculs de pourcentages, construction rigoureuse d'une carte de flux (savoir-faire cartographique majeur), analyse de photographies géographiques, rédaction argumentée.

Trois idées forces structurent l'intégration nationale par le commerce :
\begin{enumerate}
\item \textbf{L'interdépendance est une réalité mesurable :} Aucune région n'est autosuffisante ; la complémentarité des climats et des sols (maraîchage Ouest, vivriers Centre-Sud, élevage Nord, huile/poisson Littoral-SO) tisse un maillage de flux vitaux.
\item \textbf{Les infrastructures sont le goulot d'étranglement :} L'état des routes, la vétusté du rail, l'absence de logistique du froid transforment l'avantage comparatif en coût comparatif désavantageux (prix chers, pertes, insécurité alimentaire relative).
\item \textbf{L'organisation des acteurs est le levier humain :} Les bayam selam, malgré leur rôle central, restent vulnérables. Leur professionnalisation (coopératives), l'aménagement de marchés de gros modernes et la suppression des tracasseries administratives sont les conditions \emph{sine qua non} d'une intégration nationale efficace et équitable.
\end{enumerate}

\end{exoresolu}

\begin{savaistu}[Prolongement : Le corridor Douala--Ngaoundéré]
Le corridor Douala--Ngaoundéré (route N1/N4 + rail Transcam) est la colonne vertébrale de l'intégration Nord-Sud. Son aménagement (projet de route à 2x2 voies Yaoundé--Bafoussam, réhabilitation Transcam financé par la BEI/AFD, port sec de Ngaoundéré) est la priorité n°1 de la Stratégie Nationale de Développement (SND30) pour réduire les coûts de transport de 30 \% et booster les échanges inter-régionaux.
\end{savaistu}

\newpage

\coursec{Méthodes et exercices résolus}

\coursec{Les échanges inter-régionaux au Cameroun}

\courssub{Introduction : notions clés et situation-problème}

\begin{exemplebox}[Situation-problème : le marché Central de Douala, un carrefour national]
Chaque matin, au marché Central de Douala (région du Littoral), des camions bâchés déchargent des régimes de plantains venus de l'Ouest (Bafoussam), des sacs de maïs du Nord-Ouest (Bamenda), des carcasses de bœufs du Grand Nord (Garoua, Ngaoundéré), des caisses de poisson fumé de la côte (Kribi, Youpwe) et des sacs de cacao du Centre-Sud (Ebolowa, Sangmélima) destinés à l'exportation par le port autonome de Douala. En sens inverse, des camions repartent chargés de riz importé, d'huile, de tissus, de motos et de médicaments vers l'intérieur du pays.
\textbf{Problématique :} Comment les infrastructures de transport, l'organisation des acteurs du commerce (bayam selam, marchés périodiques, vendeurs à la sauvette) et la complémentarité des régions permettent-elles — ou freinent-elles — l'intégration nationale par les échanges ?
\end{exemplebox}

\begin{definition}[Vocabulaire clé de la leçon]
\begin{itemize}
\item \cle{Flux commercial} : déplacement d'un volume de marchandises entre une zone de production (ou d'entrée) et une zone de consommation (ou de sortie).
\item \cle{Interdépendance régionale} : relation de complémentarité où chaque région a besoin des productions des autres pour satisfaire ses besoins (ex. : le Nord fournit la viande, le Sud fournit le cacao et les produits manufacturés/importés).
\item \cle{Marché périodique} : marché qui se tient à jours fixes (hebdomadaire, bi-hebdomadaire) dans une localité rurale ou semi-urbaine, point de rencontre entre producteurs locaux et collecteurs (bayam selam).
\item \cle{Bayam selam} (ou \emph{buyam-sellam}) : commerçants itinérants, souvent des femmes, qui achètent les vivres en gros sur les marchés de production rurale pour les revendre en gros ou demi-gros sur les marchés urbains. Terme pidgin : « buy them, sell them ».
\item \cle{Vendeur à la sauvette} : détaillant informel urbain vendant sur le trottoir, aux carrefours, sans local fixe, souvent des produits vivriers ou manufacturés de petite valeur.
\item \cle{Centre commercial} : grande surface moderne (supermarché, hypermarché) à gestion centralisée, approvisionnée par des centrales d'achat, ciblant une clientèle urbaine aisée (ex. : Super U, Casino, Mahima, Spar au Cameroun).
\end{itemize}
\end{definition}

\courssub{Les facteurs de développement des échanges : les infrastructures}

\begin{propriete}[L'état des infrastructures de transport au Cameroun (ordre de grandeur 2023-2024)]
Le réseau de transport structure l'espace camerounais mais présente de forts contrastes :
\begin{itemize}
\item \textbf{Routes :} Réseau estimé à \num{50000} km dont environ \num{6000} à \num{6500} km bitumés (en augmentation grâce aux programmes d'urgence et aux chantiers routiers). Les axes structurants sont la N1 (Douala--Yaoundé--Bertoua--Ngaoundéré--Garoua--Maroua), la N3 (Yaoundé--Bafoussam--Bamenda), la N10 (Bafoussam--Banyo--Ngaoundéré). \textbf{Limite :} entretien irrégulier, dégradation rapide en saison des pluies, rackets aux barrages, pistes rurales souvent impraticables (enclavement des bassins de production).
\item \textbf{Chemin de fer :} \cle{Transcamerounais} (Douala--Ngaoundéré, \SI{886}{km}), concessionné à \textbf{Camrail} (groupe Bolloré/APM Terminals). Trafic : \num{1} à \num{1,5} million de tonnes/an (fret dominant : hydrocarbures, coton, bauxite, bois, conteneurs). \textbf{Limite :} voie unique, vétusté du matériel roulant, déraillements fréquents, vitesse moyenne faible (\SI{30}{km/h}), ne dessert pas l'Ouest ni le Nord-Ouest agricoles.
\item \textbf{Ports :} \textbf{Port autonome de Douala} (95 \% du trafic portuaire, \num{12} à \num{14} millions de tonnes/an), saturé, accès routier/ferroviaire congestionné. \textbf{Port en eau profonde de Kribi} (phase 1 opérationnelle depuis 2014, géré par \textbf{Kribi Conteneur Terminal} -- KCT, filiale de Bolloré/APM Terminals ; phase 2 en projet), spécialisé vrac solide (fer, bauxite), conteneurs, hydrocarbures (GNL). \textbf{Limite :} connexion ferroviaire Kribi--Mbalam (fer) en projet, route Kribi--Edea--Douala unique.
\item \textbf{Transport aérien :} Aéroports internationaux de Douala, Yaoundé-Nsimalen, Garoua, Maroua. Compagnie nationale \textbf{Camair-Co} en grande difficulté financière (flotte réduite, vols intérieurs irréguliers), concurrence de compagnies étrangères (Ethiopian, ASKY, Air France, Brussels Airlines, Turkish Airlines). Rôle marginal pour le fret intérieur (coût élevé).
\end{itemize}
\end{propriete}

\begin{savaistu}[Les corridors de développement]
Le Cameroun est le poumon logistique de la CEMAC. Trois corridors structurent les échanges internationaux et inter-régionaux :
\begin{enumerate}
\item \textbf{Corridor Douala--N'Djamena (Tchad) :} Route N1 + Rail Transcamerounais. Vital pour l'import/export tchadien (coton, bétail, pétrole).
\item \textbf{Corridor Douala--Bangui (RCA) :} Route N1 puis N4 (Bertoua--Garoua-Boulaï--Frontière).
\item \textbf{Corridor Kribi--Mbalam (projet) :} Liaison ferroviaire dédiée au minerai de fer (Sundance Resources / AustSino) vers le port de Kribi.
\end{enumerate}
La réhabilitation du Transcamerounais (projet financé par la BEI, l'AFD, l'État) est la priorité n°1 pour fluidifier les flux Nord-Sud.
\end{savaistu}

\begin{exoresolu}[Analyse d'une carte des infrastructures de transport]
\textbf{Énoncé :} À partir de la carte schématique ci-dessous, identifiez les deux principales villes portuaires, tracez l'axe ferroviaire, nommez trois villes reliées par le rail et expliquez pourquoi l'Ouest agricole n'est pas connecté au rail.
\tcblower
\textbf{Solution détaillée}
\begin{popfigure}
\begin{center}
\includegraphics[width=\linewidth,height=11cm,keepaspectratio]{N09/cmr_rail.png}
\end{center}
\legende{Carte du réseau ferré du Cameroun (trait vert) : la ligne relie Douala à Yaoundé, puis, par Bélabo, à Ngaoundéré ; des embranchements desservent Kumba, Nkongsamba et Mbalmayo. On y lit : 1. le port de Douala, à l'extrémité sud-ouest de la ligne (Kribi est à situer plus au sud, sur la côte) ; 2. les villes reliées par le Transcamerounais (Douala, Yaoundé, Bélabo, Ngaoundéré) ; 3. l'absence de voie ferrée à l'Ouest (Bafoussam) et au Nord-Ouest (Bamenda), qui n'apparaissent pas sur le tracé. \\ Source : Wikimedia Commons, Trizek, CC BY-SA 3.0.}
\end{popfigure}
\emph{Réponse :} Les ports sont Douala et Kribi. L'axe ferroviaire (Transcamerounais) relie Douala, Yaoundé, Ngaoundéré, Garoua (et Maroua par la route). L'Ouest (Bafoussam) et le Nord-Ouest (Bamenda) sont en dehors du tracé ferroviaire historique (coloniale : Douala--Yaoundé--Ngaoundéré pour l'évacuation du coton/bétail du Nord), ce qui les oblige à dépendre uniquement de la route (N3, N6) pour écouler leurs productions (pommes de terre, maïs, haricots, café arabica).
\end{exoresolu}

\courssub{L'interdépendance entre les régions}

\begin{definition}[Complémentarité des espaces de production et de consommation]
Le territoire camerounais se divise en grandes régions agricoles et écologiques spécialisées, créant des flux obligatoires vers les métropoles de consommation (Douala, Yaoundé) et le port d'exportation (Douala, Kribi).
\end{definition}

\begin{exemplebox}[Tableau de synthèse : Spécialisations régionales et flux principaux (estimations annuelles, ordres de grandeur)]
\begin{center}
\begin{tabularx}{\linewidth}{|l|X|X|X|}\hline
\rowcolor{popblueL} \textbf{Région / Zone} & \textbf{Productions phares (Export / Intérieur)} & \textbf{Flux sortants principaux} & \textbf{Besoins / Flux entrants} \\ \hline
\textbf{Grand Nord} (Adamaoua, Nord, Extrême-Nord) & Bétail (bœufs, moutons), Coton, Oignon, Mil/Sorgho, Arachide & Bétail sur pied $\rightarrow$ Yaoundé, Douala ; Coton $\rightarrow$ Usines (Garoua) / Export Douala/Kribi ; Oignon $\rightarrow$ Tout le pays & Riz, Huile, Poisson, Produits manufacturés, Carburant (venus du Sud) \\ \hline
\textbf{Ouest / Nord-Ouest} (Bafoussam, Bamenda, Dschang) & Pomme de terre, Maïs, Haricot, Chou, Carotte, Café arabica, Porc, Volaille & Vivres frais $\rightarrow$ Yaoundé, Douala (flux massifs quotidiens) ; Café $\rightarrow$ Export Douala & Riz, Huile, Poisson, Sel, Produits manufacturés, Engrais \\ \hline
\textbf{Centre / Sud / Est} (Yaoundé, Ebolowa, Sangmélima, Bertoua, Abong-Mbang) & Cacao, Café robusta, Bois, Manioc, Macabo, Plantain, Ananas & Cacao/Bois $\rightarrow$ Port Douala/Kribi (Export) ; Vivres $\rightarrow$ Yaoundé, Douala & Poisson, Viande (Nord), Riz, Huile, Ciment, Carburant \\ \hline
\textbf{Littoral / Sud-Ouest} (Douala, Edéa, Nkongsamba, Buea, Limbe, Kumba) & Banane (plantations CDC, PHP, Boh), Palmier à huile (Socapalm, Safacam), Hévéa (Hévécam), Pétrole (offshore), Pêche industrielle/artisanale & Banane/Huile/Caoutchouc $\rightarrow$ Export Douala/Kribi ; Poisson $\rightarrow$ Intérieur (Nord, Yaoundé) ; Pétrole $\rightarrow$ Raffinerie Sonara (Limbe) $\rightarrow$ Intérieur & Vivres (Ouest, Nord, Centre), Bétail (Nord), Main-d'œuvre, Matériel industriel \\ \hline
\end{tabularx}
\end{center}
\legende{Synthèse des interdépendances. Sources : MINADER, MINEPAT, INS, rapports Camrail, Port Autonome de Douala (données annuelles récentes arrondies).}
\end{exemplebox}

\begin{exoresolu}[Caractériser l'interdépendance Nord-Sud à partir de données]
\textbf{Énoncé :} Le Grand Nord fournit environ 80 \% de la viande bovine consommée à Douala et Yaoundé (estim. \num{150000} tonnes équivalent carcasse/an). Le prix du bœuf à Douala a augmenté de 40 \% en 5 ans. Expliquez ce lien par les infrastructures.
\tcblower
\textbf{Solution détaillée}
L'interdépendance est forte : le Nord produit (pâturages, eau), le Sud consomme (population dense, pouvoir d'achat). Le flux emprunte la N1 (bitumée) et le Transcamerounais. La hausse de 40 \% s'explique par : (1) la hausse du coût du transport (carburant, péages, rackets, entretien camions sur route dégradée par endroits) ; (2) les pertes au chargement/déchargement et la mortalité en route (vétusté wagons bétail, longs arrêts) ; (3) la demande urbaine croissante (démographie, urbanisation) non suivie d'une offre de transport efficace. Le rail, s'il était performant, diviserait le coût par 2 à 3.
\end{exoresolu}

\courssub{L'organisation du commerce intérieur}

\begin{definition}[Les trois piliers du commerce intérieur camerounais]
\begin{enumerate}
\item \textbf{Les marchés périodiques (ruraux et périurbains) :} Points de collecte hebdomadaires (ex. : Bafia, Mbouda, Magba, Guider, Mokolo-Yaoundé jour de marché). Rôle : fixation des prix à la production, rencontre producteurs/bayam selam, lien social. Fréquence : 1 jour/semaine (jour de marché) ou 2 jours (ex. : 8 jours).
\item \textbf{Les bayam selam (collecteurs / grossistes itinérants) :} Majoritairement des femmes. Circuit : Marché rural (achat sacs 50/100 kg) $\rightarrow$ Camion (location ou propriété) $\rightarrow$ Marché urbain de gros (ex. : Marché Central Douala, Marché Mfoundi Yaoundé, Marché Nkoulouloun) $\rightarrow$ Revendeurs détaillants. Elles assurent le financement (avances aux producteurs), le tri, le conditionnement, le transport, le risque commercial.
\item \textbf{Les vendeurs à la sauvette (détaillants informels urbains) :} Dernier maillon. Vente à l'unité ou petit conditionnement (tas de manioc, bols de maïs, morceaux de viande). Emplois de survie, faibles barrières à l'entrée, exposition aux intempéries et aux opérations de « nettoyage » municipal.
\end{enumerate}
\end{definition}

\begin{propriete}[L'émergence du commerce moderne : les centres commerciaux]
Depuis les années 2010, développement de grandes surfaces à Douala et Yaoundé : \textbf{Super U} (groupe CFAO/Carrefour), \textbf{Casino} (groupe CFAO), \textbf{Mahima} (groupe L'Anecdote), \textbf{Spar} (franchise), \textbf{City Market}, \textbf{Marche Central moderne} (réhabilité).
\begin{itemize}
\item \textbf{Part de marché :} estimée < 10 \% du commerce de détail alimentaire (majorité reste informelle).
\item \textbf{Approvisionnement :} Centrales d'achat importatrices (riz, huile, conserves, surgelés) + contrats avec quelques gros producteurs locaux (poulets, légumes calibrés, ananas). Marges sur les produits locaux souvent élevées.
\item \textbf{Clientèle :} Classes moyennes/aisées, expatriés. Paiement carte/espèces. Hygiène, climatisation, parking.
\item \textbf{Impact :} Concurrence pour les bayam selam sur les produits standards ; mais créent des débouchés pour quelques coopératives structurées (ex. : pommes de terre calibrées, poulets de chair).
\end{itemize}
\end{propriete}

\begin{exoresolu}[Analyse d'un organigramme de la filière pomme de terre Ouest $\rightarrow$ Yaoundé]
\textbf{Énoncé :} Reconstituez la chaîne des acteurs et les marges types (ordres de grandeur) pour 1 kg de pomme de terre variété \emph{Nicola} ou \emph{Spunta} produite à Dschang (Ouest) et consommée à Yaoundé.
\tcblower
\textbf{Solution détaillée}
\begin{center}
\begin{tabularx}{\linewidth}{|l|X|c|X|}\hline
\rowcolor{popblueL} \textbf{Étape / Acteur} & \textbf{Opération} & \textbf{Prix indicatif (FCFA/kg)} & \textbf{Marge / Coût} \\ \hline
Producteur (Dschang) & Culture, récolte, tri grossier & 150 -- 200 & Revenu producteur \\ \hline
Bayam selam (collecte) & Achat, sac 50kg, chargement camion & 220 -- 250 & Marge collecte + risque \\ \hline
Transport (Camion 10-20t) & Dschang $\rightarrow$ Yaoundé (\SI{300}{km}, 6-8h) & + 30 -- 50 & Carburant, péages, rackets, entretien \\ \hline
Bayam selam (gros Yaoundé) & Déchargement, tri, reconditionnement, vente demi-gros & 300 -- 350 & Marge gros + stockage \\ \hline
Revendeuse détaillée (Marché Mfoundi / sauvette) & Vente au tas / kg & 400 -- 500 & Marge détail + pertes (germes, pourriture) \\ \hline
\textbf{Consommateur final} & & \textbf{400 -- 500} & \\ \hline
\end{tabularx}
\end{center}
\emph{Enseignement :} Le producteur ne touche que 30 à 50 \% du prix final. Les coûts de transport et de commercialisation (informels) absorbent le reste. L'amélioration des routes et l'organisation en coopératives (vente groupée, stockage frigorifique) permettraient de réduire les intermédiaires et les pertes.
\end{exoresolu}

\courssub{Les difficultés du commerce intérieur}

\begin{methode}[Analyser les freins au commerce intérieur : la chaîne causale]
Pour toute difficulté identifiée, construire l'enchaînement :
\textbf{Cause structurelle} (ex. : route dégradée) $\rightarrow$ \textbf{Contrainte opérationnelle} (ex. : temps de trajet doublé, camion en panne) $\rightarrow$ \textbf{Coût économique} (ex. : surcoût transport + 30 \%, pertes 15 \% produit) $\rightarrow$ \textbf{Conséquence commerciale} (ex. : prix final haut, repli sur circuit court/informel, désincitation production) $\rightarrow$ \textbf{Conséquence spatiale} (ex. : enclavement région, baisse flux inter-régionaux).
\end{methode}

\begin{aretenir}[Les 4 grandes difficultés structurelles]
\begin{enumerate}
\item \textbf{Infrastructures de transport défaillantes :} Pistes rurales impraticables 3 à 4 mois/an (saison des pluies) $\rightarrow$ rupture d'approvisionnement villes, pourritures récoltes (manioc, tomate, plantain). Routes bitumées : nids-de-poule, poids-lourds surchargés (destruction chaussée), barrages multiples (police, gendarmerie, douane, conseil régional) $\rightarrow$ temps perdu, coûts illicites.
\item \textbf{Inadéquation modale :} Rail vétuste (part modale fret < 10 \%), absent de l'Ouest/Nord-Ouest. Transport fluvial (Wouri, Sanaga, Bénoué, Logone) quasi inexistant pour le fret commercial. Tout repose sur le camion (coût élevé, pollution, accidents).
\item \textbf{Informel dominant et non accompagné :} Bayam selam et vendeurs à la sauvette : pas d'accès au crédit bancaire (garanties), pas d'assurance, pas de stockage frigorifique, fiscalité harcelante (patente, taxes marchés, redevances journalières), insécurité (vols, rackets). Cela maintient des coûts de transaction élevés.
\item \textbf{Insécurité et instabilité :} Crise anglophone (Nord-Ouest, Sud-Ouest) depuis 2016 : routes coupées, marchés fermés, « ghost towns », déplacement populations $\rightarrow$ effondrement flux pommes de terre/maïs/haricots vers Douala/Yaoundé, hausse prix. Insécurité Grand Nord (Boko Haram) : perturbation marchés bétail (Maroua, Kousseri), transhumance difficile.
\end{enumerate}
\end{aretenir}

\begin{exoresolu}[Étude de cas : l'impact de la crise du Nord-Ouest/Sud-Ouest sur le marché de la pomme de terre à Douala]
\textbf{Énoncé :} Avant 2016, le Nord-Ouest (Bamenda, Kumbo, Wum) fournissait 60 \% de la pomme de terre consommée à Douala. Depuis la crise, l'approvisionnement repose quasi exclusivement sur l'Ouest (Dschang, Bafoussam, Santchou). Le prix moyen du kg a doublé (de 250 à 500 FCFA). Analysez cette situation.
\tcblower
\textbf{Solution détaillée}
\emph{1. Rupture de l'offre principale :} L'insécurité (blocages routes, incendies marchés, fuite agriculteurs) a stoppé la production et l'évacuation du Nord-Ouest. Perte de 60 \% de l'offre structurelle.
\emph{2. Report sur l'Ouest :} L'Ouest (capacité limitée, terres saturées, pression foncière) ne peut compenser pleinement. Hausse de la demande sur l'Ouest $\rightarrow$ renchérissement prix à la production (150 $\rightarrow$ 250 FCFA/kg).
\emph{3. Coûts de transport accrus :} Déviation des camions (évitement zones chaudes), allongement trajets, risques vols/rançons $\rightarrow$ surprime transport.
\emph{4. Conséquence :} Prix final doublé pour le consommateur doualais. Substitution partielle par le riz/pâtes (importés) ou manioc/local. Perte de revenus pour les producteurs restants du Nord-Ouest. Illustration parfaite : instabilité politique $\rightarrow$ rupture flux commercial $\rightarrow$ renchérissement vie $\rightarrow$ dépendance accrue aux importations.
\end{exoresolu}

\courssub{TP 4 : Carte des flux commerciaux inter-régionaux au Cameroun}

\begin{methode}[Conduire le TP 4 : de l'inventaire à la carte finale]
\begin{enumerate}
\item \textbf{Inventorier} : à partir des documents fournis (cartes de production agricole, tableaux de tonnages, photographies de marchés), dresser la liste des couples produit / zone de production / zone de consommation. Exemples : cacao (Centre, Sud, Est $\rightarrow$ port de Douala) ; bananes (Littoral, Sud-Ouest $\rightarrow$ Douala et exportation) ; bétail (Grand Nord $\rightarrow$ villes du Sud) ; pommes de terre (Ouest, Nord-Ouest $\rightarrow$ Douala, Yaoundé) ; poisson fumé (Littoral, fleuves $\rightarrow$ tout le pays).
\item \textbf{Classer} les flux : flux internes inter-régionaux (vivres, bétail) et flux d'exportation (cacao, café, bananes, coton vers le port de Douala/Kribi).
\item \textbf{Choisir les symboles} : couleurs de fond pour les zones de production, points pour les villes de commercialisation, flèches proportionnelles pour les flux, traits pour les axes (route, rail, port).
\item \textbf{Tracer} le fond de carte du Cameroun, placer les régions et les villes principales (Douala, Yaoundé, Garoua, Maroua, Ngaoundéré, Bafoussam, Bamenda, Buea, Bertoua, Ebolowa, Kribi).
\item \textbf{Légender} de façon exhaustive et hiérarchisée, ajouter titre, nord, échelle indicative.
\item \textbf{Rédiger un court commentaire} (5 lignes) de la carte obtenue : polarisation par Douala et Yaoundé, complémentarité Nord--Sud, rôle du corridor Douala--Ngaoundéré.
\end{enumerate}
\end{methode}

\begin{exoresolu}[TP 4 guidé : cartographier trois flux majeurs]
\textbf{Énoncé} : À partir de l'inventaire suivant — (a) cacao du Centre et du Sud vers le port de Douala ; (b) bétail du Grand Nord vers Douala et Yaoundé ; (c) pommes de terre et maïs de l'Ouest vers Yaoundé et Douala — construisez la carte des flux et rédigez un commentaire de 5 lignes.
\tcblower
\textbf{Solution détaillée}

\emph{Étape 1 — Choix des symboles.} Zones de production en teintes de fond distinctes ; villes de consommation en points bleus ; flux internes en flèches orange ; flux d'exportation en flèche verte vers le port ; port de Douala symbolisé par une ancre.

\begin{popfigure}
\begin{center}
\includegraphics[width=\linewidth,height=9.5cm,keepaspectratio]{N01/cmr_map.jpg}
\end{center}
\legende{Fond de carte du Cameroun (relief, régions, villes, routes, voie ferrée). Le croquis attendu situe : 1. la zone d'élevage (Grand Nord) ; 2. la zone cacaoyère (Centre--Sud) ; 3. la zone vivrière de l'Ouest et du Nord-Ouest ; 4. les villes de consommation et le port de Douala ; 5. le flux d'exportation du cacao vers Douala ; 6. les flux internes (bétail, vivres) ; 7. le Transcamerounais, visible sur le fond de carte. \\ Source : Wikimedia Commons, CIA, domaine public.}
\end{popfigure}

\emph{Étape 2 — Commentaire rédigé.} La carte montre que les flux commerciaux inter-régionaux convergent vers deux pôles, Douala et Yaoundé, et vers le port de Douala pour l'exportation (cacao). Ils traduisent une forte complémentarité : le Grand Nord fournit la viande, l'Ouest les vivres, le Centre-Sud le cacao. L'axe Douala--Ngaoundéré (rail et route) structure l'ensemble, mais sa vétusté et la mauvaise connexion des ports avec l'intérieur freinent ces échanges.

\emph{Conclusion.} Le TP confirme que le commerce intérieur camerounais repose sur l'interdépendance des régions, polarisée par les métropoles du Sud et dépendante d'infrastructures perfectibles.
\end{exoresolu}

\courssub{Savoir-faire 1 : Observer, localiser et décrire un phénomène géographique}

\begin{methode}[Décrire une carte ou une photographie de commerce intérieur]
La description est la première étape de tout commentaire de document au Bac. Elle doit être \textbf{objective} : on dit ce qu'on voit, pas ce qu'on pense.
\begin{enumerate}
\item \textbf{Identifier le document} : nature (carte, photographie, croquis, tableau), titre, source, date, échelle éventuelle.
\item \textbf{Localiser} : situer le lieu dans l'espace national (région, département, ville, axe de communication). Exemple : le marché de Mokolo se situe à Yaoundé, région du Centre, sur l'axe Yaoundé--Bafoussam.
\item \textbf{Décrire du général au particulier} : d'abord l'ensemble (un marché dense, un flux nord--sud), puis les détails (étals, produits visibles, moyens de transport, sens des flux).
\item \textbf{Nommer précisément} : utiliser le vocabulaire de la leçon (\cle{marché périodique}, \cle{bayam selam}, \cle{vendeur à la sauvette}, \cle{flux commercial}, \cle{zone de production}, \cle{zone de consommation}).
\item \textbf{Chiffrer si possible} : repérer les ordres de grandeur (épaisseur des flèches, tonnages, fréquence des marchés).
\item \textbf{Conclure la description} par une phrase bilan qui annonce l'interprétation à venir.
\end{enumerate}
\textbf{Réflexe} : interdiction absolue d'interpréter pendant la description (« ce marché montre que... » est à réserver à l'étape suivante).
\end{methode}

\begin{exoresolu}[Description d'une photographie de marché]
\textbf{Énoncé (type Bac)} : À partir de la photographie du marché hebdomadaire de Bafia (région du Centre), prise un jour de marché, décrivez ce que vous observez en 8 à 10 lignes, en localisant le lieu et en nommant les acteurs et les produits visibles. La photographie montre : des centaines de vendeuses accroupies derrière des cagettes de manioc, de macabo, d'arachides et de feuilles (nzolé) ; des camions bâchés garés en bordure ; des acheteurs venus en cars ; des revendeurs chargeant des sacs de 50 kg.
\tcblower
\textbf{Solution détaillée}

\emph{Étape 1 — Identification.} Il s'agit d'une photographie de paysage commercial, prise un jour de marché hebdomadaire à Bafia, dans le département du Mbam-et-Inoubou, région du Centre.

\emph{Étape 2 — Localisation.} Bafia se situe au nord de Yaoundé, sur l'axe routier Yaoundé--Bafoussam (N4 puis embranchements), ce qui en fait un point de collecte entre les zones de production rurales du Mbam et la métropole de Yaoundé.

\emph{Étape 3 — Description du général au particulier.} L'ensemble de la photographie montre un marché très fréquenté, organisé en plein air : c'est un \cle{marché périodique} (il se tient un jour fixe par semaine). Au premier plan, des vendeuses détaillantes proposent des vivres frais : manioc, macabo, arachides, légumes-feuilles. À l'arrière-plan, des camions bâchés et des sacs de 50 kg en cours de chargement indiquent une vente en gros et demi-gros. Les revendeurs qui achètent ces sacs pour les revendre dans les villes sont les \cle{bayam selam} (littéralement « ceux qui achètent et vendent »).

\emph{Étape 4 — Nommer et chiffrer.} Les produits visibles sont des vivres (manioc, macabo, arachides) typiques des zones de production de la région du Centre ; les sacs de 50 kg et les camions montrent que les quantités échangées dépassent la consommation locale.

\emph{Conclusion.} Cette photographie illustre un maillon essentiel du commerce intérieur camerounais : le marché périodique rural, point de rencontre entre producteurs, bayam selam et consommateurs, et point de départ des flux de vivres vers les grandes villes.
\end{exoresolu}

\courssub{Savoir-faire 2 : Construire et légender un croquis}

\begin{methode}[Construire un croquis de flux commerciaux inter-régionaux]
Le croquis est un exercice majeur au Bac de géographie. Il note la capacité à \textbf{sélectionner}, \textbf{hiérarchiser} et \textbf{légender}.
\begin{enumerate}
\item \textbf{Tracer le fond de carte} : contour très simplifié du Cameroun, orientation (flèche du nord), échelle indicative.
\item \textbf{Choisir un titre} précis : « Les flux de vivres et de produits tropicaux entre les régions du Cameroun ».
\item \textbf{Placer les zones} : zones de production (ex. Nord : bétail, coton ; Ouest et Nord-Ouest : pommes de terre, maïs ; Centre et Sud : manioc, cacao ; Littoral : bananes, huile de palme) et zones de consommation (Douala, Yaoundé, villes du Nord).
\item \textbf{Tracer les flux} : flèches dont l'\textbf{épaisseur est proportionnelle} à l'importance du flux ; indiquer le produit transporté le long de chaque flèche.
\item \textbf{Ajouter les axes} : routes bitumées (trait plein), chemin de fer (trait avec traverses ou pointillés spécifiques), port de Douala.
\item \textbf{Rédiger la légende} : 4 à 6 entrées maximum, classées (d'abord les zones, puis les flux, puis les axes). Chaque symbole du croquis doit figurer dans la légende, et rien dans la légende ne doit manquer sur le croquis.
\item \textbf{Vérifier} : titre, nord, échelle, légende, propreté.
\end{enumerate}
\textbf{Réflexe} : un croquis sans titre ou sans légende perd la moitié des points, même s'il est exact.
\end{methode}

\begin{exoresolu}[Croquis du flux bétail--viande Nord vers Sud]
\textbf{Énoncé (type Bac)} : Le Grand Nord (régions de l'Extrême-Nord, du Nord et de l'Adamaoua) fournit l'essentiel du bétail consommé dans les grandes villes du Sud (Douala, Yaoundé). Les animaux partent des marchés à bétail de Garoua, Ngaoundéré et Maroua et transitent par la route ou le train jusqu'à Ngaoundéré, puis par la route vers le Sud. Construisez un croquis simple montrant les zones de production, les zones de consommation et les flux de bétail.
\tcblower
\textbf{Solution détaillée}

\emph{Raisonnement.} On identifie d'abord les trois ensembles : (1) la zone de production = le Grand Nord ; (2) les zones de consommation = Douala et Yaoundé ; (3) le flux = flèches Nord $\rightarrow$ Sud, épaisses car le flux est massif (le Grand Nord fournit la grande majorité de la viande bovine du pays). On ajoute les villes-relais (Garoua, Ngaoundéré, Maroua) et l'axe ferroviaire Douala--Ngaoundéré (Transcamerounais), même si le bétail emprunte surtout la route.

\begin{popfigure}
\begin{center}
\includegraphics[width=\linewidth,height=11cm,keepaspectratio]{N09/cmr_rail.png}
\end{center}
\legende{Fond de carte : la voie ferrée du Cameroun (trait vert) de Douala à Ngaoundéré par Yaoundé et Bélabo, avec les embranchements vers Kumba, Nkongsamba et Mbalmayo. Le croquis attendu situe : 1. la zone de production (Grand Nord : élevage), au nord du terminus de Ngaoundéré ; 2. les villes-relais et marchés à bétail (Garoua, Maroua, Ngaoundéré) ; 3. les villes de consommation (Yaoundé, Douala) ; 4. le flux de bétail vers le sud, à tracer avec une épaisseur proportionnelle à son importance ; 5. le chemin de fer Transcamerounais, déjà visible sur le fond. \\ Source : Wikimedia Commons, Trizek, CC BY-SA 3.0.}
\end{popfigure}

\emph{Justification des choix.} Les flèches sont orientées du Nord vers le Sud, sens réel du flux ; leur épaisseur (2,5 à 3 pt) traduit un flux dominant. Les villes de consommation sont en bleu, la zone de production en orange, conformément à la légende.

\emph{Conclusion.} Ce croquis met en évidence l'\cle{interdépendance} entre les régions : le Grand Nord, zone d'élevage, approvisionne en viande les métropoles du Sud, tandis que les revenus de la vente du bétail irriguent l'économie nordiste.
\end{exoresolu}

\courssub{Savoir-faire 3 : Inventorier, classer et interpréter des documents statistiques}

\begin{methode}[Exploiter un tableau statistique et calculer des parts et des évolutions]
\begin{enumerate}
\item \textbf{Lire le tableau en entier} : titre, unités, dates, source. Ne jamais commenter un chiffre sans connaître son unité.
\item \textbf{Classer} : ranger les données (croissant ou décroissant), repérer le maximum et le minimum, calculer les écarts.
\item \textbf{Calculer une part en pourcentage} :
\[ \text{part} = \frac{\text{valeur de la partie}}{\text{valeur totale}} \times 100 \]
\item \textbf{Calculer un taux de variation} entre deux dates :
\[ t = \frac{V_{\text{finale}} - V_{\text{initiale}}}{V_{\text{initiale}}} \times 100 \]
Un résultat positif est une hausse, négatif une baisse. On peut aussi dire « la valeur a été multipliée par $V_f/V_i$ ».
\item \textbf{Interpréter} : relier chaque fait chiffré à une explication géographique (infrastructures, interdépendance, urbanisation).
\item \textbf{Rédiger} : une phrase par idée, avec le chiffre exact cité entre parenthèses.
\end{enumerate}
\textbf{Réflexe} : au Bac, tout chiffre calculé doit être arrondi proprement (une décimale suffit) et accompagné de son unité ou du symbole \%.
\end{methode}

\begin{exoresolu}[Calculs sur les flux de produits vers Douala et Yaoundé]
\textbf{Énoncé (type Bac)} : Le tableau ci-dessous donne une estimation des tonnages de produits alimentaires acheminés chaque année vers Douala et Yaoundé (données fictives mais d'ordre de grandeur réaliste).

\begin{center}
\begin{tabularx}{\linewidth}{|l|X|X|}\hline
\rowcolor{popblueL} \textbf{Produit} & \textbf{Tonnage (milliers de tonnes)} & \textbf{Origine principale} \\ \hline
Vivres (manioc, macabo, plantain) & 900 & Centre, Ouest, Littoral \\ \hline
Maïs et pommes de terre & 350 & Ouest, Nord-Ouest \\ \hline
Viande (bétail) & 150 & Grand Nord \\ \hline
Poissons et produits de mer & 100 & Littoral, importations \\ \hline
\textbf{Total} & \textbf{1500} & \\ \hline
\end{tabularx}
\end{center}

1. Calculez la part de chaque catégorie dans le total. 2. Le tonnage de viande était de 120 milliers de tonnes dix ans plus tôt : calculez le taux de variation. 3. Tirez deux enseignements géographiques.
\tcblower
\textbf{Solution détaillée}

\emph{1. Parts en pourcentage.} On applique $\text{part} = \dfrac{\text{partie}}{\text{total}} \times 100$ avec un total de \num{1500} milliers de tonnes :
\begin{align*}
\text{Vivres} &: \frac{900}{1500} \times 100 = 60\,\% \\
\text{Maïs et pommes de terre} &: \frac{350}{1500} \times 100 \approx 23{,}3\,\% \\
\text{Viande} &: \frac{150}{1500} \times 100 = 10\,\% \\
\text{Poissons} &: \frac{100}{1500} \times 100 \approx 6{,}7\,\%
\end{align*}
Vérification : $60 + 23{,}3 + 10 + 6{,}7 = 100\,\%$. Le calcul est cohérent.

\emph{2. Taux de variation de la viande.}
\[ t = \frac{150 - 120}{120} \times 100 = \frac{30}{120} \times 100 = 25\,\% \]
Le tonnage de viande a augmenté de 25 \% en dix ans (il a été multiplié par $150/120 = 1{,}25$).

\emph{3. Enseignements géographiques.}
\begin{itemize}
\item Les vivres dominent largement les flux (60 \%) : les métropoles dépendent des campagnes proches (Centre, Ouest, Littoral), ce qui illustre le rôle des \cle{bayam selam} et des marchés périodiques dans l'approvisionnement urbain.
\item La hausse de 25 \% des flux de viande traduit l'interdépendance croissante entre le Grand Nord (élevage) et les villes du Sud (consommation), rendue possible par les routes bitumées malgré leur entretien irrégulier.
\end{itemize}

\emph{Conclusion.} Les calculs confirment que le commerce intérieur camerounais repose d'abord sur les échanges de produits alimentaires entre régions complémentaires, flux en croissance avec l'urbanisation.
\end{exoresolu}

\courssub{Savoir-faire 4 : Commenter un document (méthode du commentaire composé de géographie)}

\begin{methode}[Rédiger un commentaire de document en trois étapes]
\begin{enumerate}
\item \textbf{Introduction} (3 à 4 lignes) : présenter le document (nature, titre, source, date), situer le thème dans la leçon, annoncer la problématique. Exemple de problématique : « Comment le commerce intérieur s'organise-t-il au Cameroun et quelles difficultés rencontre-t-il ? »
\item \textbf{Développement} en deux ou trois parties annoncées : I. Une description/analyse des faits (ce que montre le document) ; II. Les explications (facteurs : infrastructures, interdépendance) ; III. éventuellement les limites et difficultés (routes dégradées, chemin de fer vétuste, ports mal connectés à l'arrière-pays).
\item \textbf{Conclusion} : répondre à la problématique, ouvrir éventuellement (ex. : les projets de réhabilitation du Transcamerounais).
\end{enumerate}
\textbf{Réflexes de rédaction} : phrases courtes, vocabulaire de la leçon souligné mentalement, un chiffre du document cité par partie, aucune connaissance « hors document » non reliée au document.
\end{methode}

\begin{exoresolu}[Commentaire guidé d'un texte sur les bayam selam]
\textbf{Énoncé (type Bac)} : « Parties avant l'aube des marchés de l'Ouest, les bayam selam chargent chaque semaine des centaines de sacs de pommes de terre et de maïs dans des camions souvent surchargés. Après des heures de route sur des pistes dégradées en saison des pluies, elles revendent à Douala et Yaoundé, dans les marchés ou parfois sur les trottoirs, en concurrence avec les vendeurs à la sauvette. » Commentez ce texte en montrant le rôle des bayam selam dans le commerce intérieur et les difficultés qu'elles rencontrent (15 à 20 lignes).
\tcblower
\textbf{Solution détaillée (plan rédigé)}

\emph{Introduction.} Ce texte, extrait d'un reportage sur le commerce vivrier, décrit l'activité des \cle{bayam selam}, revendeuses et revendeurs de vivres entre les campagnes et les villes du Cameroun. On peut se demander quel rôle jouent ces commerçants dans le commerce intérieur et quelles difficultés ils affrontent.

\emph{I. Un rôle central : assurer l'interdépendance entre les régions.} Les bayam selam collectent les surplus des zones de production (pommes de terre et maïs de l'Ouest et du Nord-Ouest) et les acheminent vers les zones de consommation urbaines (Douala, Yaoundé). Ils relient ainsi les marchés périodiques ruraux aux marchés urbains permanents et garantissent l'approvisionnement alimentaire des métropoles : sans eux, l'interdépendance entre régions agricoles et villes resterait théorique. Ils jouent aussi un rôle de fixation des prix par l'offre et la demande.

\emph{II. Des difficultés liées aux infrastructures.} Le texte mentionne des « pistes dégradées en saison des pluies » et des « camions surchargés » : cela renvoie aux limites des infrastructures camerounaises — routes bitumées en augmentation mais à l'entretien irrégulier, chemin de fer vétuste qui ne dessert pas l'Ouest agricole. Ces contraintes allongent les délais, augmentent les coûts de transport et les pertes de produits périssables, et réduisent les marges des commerçants.

\emph{III. Une insertion urbaine précaire.} Arrivées en ville, les bayam selam revendent dans les marchés ou « sur les trottoirs », en concurrence avec les \cle{vendeurs à la sauvette} : le commerce intérieur reste largement informel, hors des circuits modernes des centres commerciaux, ce qui fragilise ces acteurs (pas de stockage, pas de crédit, expositions aux saisies municipales).

\emph{Conclusion.} Les bayam selam sont le maillon indispensable du commerce intérieur camerounais : ils matérialisent les flux entre régions complémentaires, mais pâtissent de l'état des infrastructures et de l'informalité. L'amélioration des routes rurales et l'organisation des marchés de gros renforceraient l'intégration nationale.
\end{exoresolu}

\courssub{Savoir-faire 5 : Réaliser le TP 4 — Carte des flux commerciaux inter-régionaux}

\begin{methode}[Conduire le TP 4 : de l'inventaire à la carte finale]
\begin{enumerate}
\item \textbf{Inventorier} : à partir des documents fournis (cartes de production agricole, tableaux de tonnages, photographies de marchés), dresser la liste des couples produit / zone de production / zone de consommation. Exemples : cacao (Centre, Sud, Est $\rightarrow$ port de Douala) ; bananes (Littoral, Sud-Ouest $\rightarrow$ Douala et exportation) ; bétail (Grand Nord $\rightarrow$ villes du Sud) ; pommes de terre (Ouest, Nord-Ouest $\rightarrow$ Douala, Yaoundé) ; poisson fumé (Littoral, fleuves $\rightarrow$ tout le pays).
\item \textbf{Classer} les flux : flux internes inter-régionaux (vivres, bétail) et flux d'exportation (cacao, café, bananes, coton vers le port de Douala).
\item \textbf{Choisir les symboles} : couleurs de fond pour les zones de production, points pour les villes de commercialisation, flèches proportionnelles pour les flux, traits pour les axes (route, rail, port).
\item \textbf{Tracer} le fond de carte du Cameroun, placer les régions et les villes principales (Douala, Yaoundé, Garoua, Maroua, Ngaoundéré, Bafoussam, Bamenda, Buea).
\item \textbf{Légender} de façon exhaustive et hiérarchisée, ajouter titre, nord, échelle indicative.
\item \textbf{Rédiger un court commentaire} (5 lignes) de la carte obtenue : polarisation par Douala et Yaoundé, complémentarité Nord--Sud, rôle du corridor Douala--Ngaoundéré.
\end{enumerate}
\end{methode}

\courssub{Savoir-faire 6 : Analyser les difficultés du commerce intérieur (étude de cas)}

\begin{methode}[Construire une étude de cas « infrastructures et commerce »]
\begin{enumerate}
\item \textbf{Dégager la thèse du document} : quelle difficulté est mise en avant (route dégradée, rail vétuste, compagnie aérienne en difficulté, port mal connecté) ?
\item \textbf{Relier la difficulté à ses conséquences commerciales} en chaîne logique : infrastructure défaillante $\rightarrow$ hausse des coûts et des délais $\rightarrow$ hausse des prix ou pertes de produits $\rightarrow$ repli sur le circuit informel (vendeurs à la sauvette) ou sur les importations.
\item \textbf{Nuancer} : rappeler les progrès (routes bitumées en augmentation, centres commerciaux modernes à Douala et Yaoundé) pour éviter le tableau uniquement négatif.
\item \textbf{Proposer une perspective} : entretien routier, réhabilitation du Transcamerounais, connexion port--arrière-pays (Douala, Kribi).
\item \textbf{Rédiger} en paragraphes liés par des connecteurs (d'abord, ensuite, en revanche, enfin).
\end{enumerate}
\end{methode}

\begin{exoresolu}[Étude de cas : la desserte de l'Ouest agricole]
\textbf{Énoncé (type Bac)} : « La région de l'Ouest, grenier vivrier du Cameroun, n'est desservie ni par le chemin de fer ni par un port ; ses pistes rurales sont difficilement praticables en saison des pluies. » Montrez, en une quinzaine de lignes, comment cette situation illustre les difficultés du commerce intérieur et quelles en sont les conséquences.
\tcblower
\textbf{Solution détaillée}

\emph{Thèse.} Le document souligne le paradoxe d'une région très productive mais mal desservie : c'est une illustration parfaite du rôle des infrastructures comme facteur — ou frein — des échanges.

\emph{Chaîne des conséquences.} L'absence de liaison ferroviaire (le Transcamerounais relie Douala à Yaoundé et Ngaoundéré, pas à l'Ouest) et la dégradation saisonnière des pistes rurales allongent les délais de collecte et augmentent les coûts de transport supportés par les bayam selam. Conséquences directes : pertes de produits périssables (pommes de terre, légumes), hausse des prix à Douala et Yaoundé, baisse des revenus des producteurs qui vendent au plus offrant sur place. Le commerce se replie alors sur les circuits courts et informels (marchés périodiques locaux, vendeurs à la sauvette en ville), au détriment de l'approvisionnement régulier des métropoles et des centres commerciaux.

\emph{Nuance.} Il faut cependant rappeler que le réseau de routes bitumées augmente (axes Douala--Bafoussam--Bamenda, Yaoundé--Bafoussam) et que l'Ouest reste le premier fournisseur vivrier des villes : la difficulté est relative, pas rédhibitoire.

\emph{Perspective.} L'entretien régulier des routes rurales et le désenclavement ferroviaire ou routier renforceraient l'interdépendance entre l'Ouest producteur et les métropoles consommatrices.

\emph{Conclusion.} Ce cas montre que le développement des échanges intérieurs dépend moins de la production que de la qualité des infrastructures de transport qui relient les régions.
\end{exoresolu}

\begin{attention}[Les 5 erreurs qui coûtent des points au Bac]
\begin{enumerate}
\item \textbf{Confondre description et interprétation.} Dans une question « décrivez », toute phrase explicative (« ce qui prouve que... ») est hors sujet. Décris d'abord, explique ensuite, uniquement si la consigne le demande.
\item \textbf{Oublier les éléments obligatoires du croquis ou de la carte} : titre, flèche du nord, échelle indicative, légende complète. Une carte sans légende, même exacte, est lourdement pénalisée. Vérifie aussi que chaque symbole de la légende figure bien sur la carte.
\item \textbf{Se tromper dans le sens des flux.} Le bétail va du Grand Nord vers les villes du Sud ; le cacao va des zones de production vers le port de Douala. Une flèche inversée détruit la crédibilité de toute la carte. Retiens : production $\rightarrow$ consommation, intérieur $\rightarrow$ port pour l'exportation.
\item \textbf{Faire des erreurs de calcul de pourcentage} : diviser par la partie au lieu du total, ou oublier de multiplier par 100 \%. Pour un taux de variation, le diviseur est toujours la valeur \emph{initiale} : $t = \frac{V_f - V_i}{V_i} \times 100$. Vérifie que la somme des parts fait 100 \%.
\item \textbf{Présenter le commerce informel comme uniquement négatif ou négligeable.} Bayam selam, marchés périodiques et vendeurs à la sauvette forment l'essentiel du commerce intérieur réel ; les centres commerciaux modernes ne concernent qu'une minorité urbaine. Un devoir qui oppose « moderne = bien » et « informel = mal » sans nuance perd des points d'analyse.
\end{enumerate}
\end{attention}

\newpage

\coursec{Exercices}

\courssub{Je vérifie mes connaissances}

\begin{exercice}[QCM : Infrastructures et acteurs du commerce]
Parmi les affirmations suivantes, cochez celle(s) qui sont exactes.
\begin{enumerate}
    \item Le réseau routier bitumé au Cameroun couvre plus de 10~000~km et son entretien est régulier sur l'ensemble du territoire.
    \item Le chemin de fer (Camrail) relie Douala à Yaoundé et Ngaoundéré ; il est performant pour le transport des hydrocarbures et du coton, mais son matériel roulant est vétuste.
    \item Les ports de Douala et de Kribi sont directement connectés au réseau ferré national pour l'évacuation des marchandises vers l'intérieur.
    \item Les « \cle{bayam sellam} » sont des revendeurs de vivres qui assurent le lien entre les marchés ruraux de production et les marchés urbains de consommation.
    \item Les vendeurs à la sauvette opèrent dans le secteur formel, paient des impôts et occupent des emplacements fixes dans les marchés modernes.
\end{enumerate}
\end{exercice}

\begin{exercice}[Vrai ou Faux justifié]
Indiquez si chaque affirmation est Vraie (V) ou Fausse (F) et justifiez brièvement votre réponse en vous appuyant sur les données du cours.
\begin{enumerate}
    \item L'interdépendance entre les régions du Cameroun repose uniquement sur l'échange de produits agricoles.
    \item L'essor des infrastructures routières a totalement résolu le problème de l'enclavement des zones de production du Nord et de l'Est.
    \item Les marchés périodiques (hebdomadaires) jouent un rôle central dans la redistribution des produits vivriers et l'approvisionnement des villes.
    \item La compagnie aérienne nationale Camair-Co assure actuellement une desserte intérieure dense et rentable entre toutes les capitales régionales.
\end{enumerate}
\end{exercice}

\begin{exercice}[Définitions et concepts clés]
Donnez une définition précise (2 à 3 lignes) pour chacun des termes suivants :
\begin{enumerate}
    \item Échanges inter-régionaux
    \item Commerce intérieur
    \item Interdépendance régionale
    \item Enclavement
    \item Bayam sellam
\end{enumerate}
\end{exercice}

\begin{exercice}[Calcul direct : Taux de variation du réseau routier]
Le réseau routier bitumé camerounais était estimé à environ 4~500~km en 2010. En 2023, il est évalué à environ 6~500~km (chiffres arrondis, sources : MINTP, Banque Mondiale).
\begin{enumerate}
    \item Calculez l'augmentation absolue du linéaire bitumé sur la période.
    \item Calculez le taux de variation (en \%) du réseau bitumé entre 2010 et 2023.
    \item En déduire l'augmentation moyenne annuelle (en km/an) sur cette période de 13 ans.
\end{enumerate}
\end{exercice}

\courssub{Je m'entraîne}

\begin{exercice}[Analyse de documents : L'interdépendance Nord-Sud]
\textbf{Document 1 :} Productions principales par région (extrait)
\begin{center}
\begin{tabularx}{\linewidth}{|l|X|}\hline
\textbf{Région} & \textbf{Productions dominantes} \\ \hline
Adamaoua, Nord, Extrême-Nord & Coton, bétail (viande, cuir), oignon, mil/sorgho, arachide \\ \hline
Ouest, Nord-Ouest & Café arabica, haricot, pomme de terre, maïs, chou, carotte \\ \hline
Sud, Est, Centre & Cacao, bois, café robusta, plantain, manioc, huile de palme \\ \hline
Littoral (Douala), Centre (Yaoundé) & Activités industrielles, tertiaires, ports, aéroports, grands marchés de consommation \\ \hline
\end{tabularx}
\end{center}

\textbf{Document 2 :} Carte simplifiée des axes de transport (description)
Les axes routiers principaux (N1, N2, N3, N10) relient Douala et Yaoundé aux régions productrices. Le rail (Camrail) suit l'axe Douala--Yaoundé--Ngaoundéré. Le port de Douala exporte le coton, le cacao, le bois, l'aluminium. Le port de Kribi exporte le fer (Mbalam), le gaz (GNL), le coton, le cacao.

\emph{Consignes :}
\begin{enumerate}
    \item En vous appuyant sur les deux documents, citez deux flux de marchandises majeurs qui illustrent l'interdépendance entre le Nord (Adamaoua, Nord, Extrême-Nord) et le Sud (Littoral, Centre, Sud, Est).
    \item Expliquez pourquoi les villes de Douala et Yaoundé sont les pôles majeurs de cette interdépendance.
    \item Quel rôle joue le chemin de fer dans l'évacuation de la production cotonnière du Nord ? Citez une limite de ce mode de transport.
\end{enumerate}
\end{exercice}

\begin{exercice}[Étude de cas : Le marché périodique de Bafia (région du Centre)]
On observe le marché hebdomadaire de Bafia (jour de marché : mardi). Les vendeurs arrivent dès l'aube : paysannes des villages environnants (plantain, macabo, légumes, fruits), \cle{bayam sellam} venues de Yaoundé (120~km) avec des camions chargés de manioc, maïs, haricots achetés la veille dans le Nord-Ouest, bouchers vendant du bétail venu de l'Adamaoua, revendeurs de produits manufacturés (savon, huile, tissus, quincaillerie) venus de Douala.
\emph{Consignes :}
\begin{enumerate}
    \item Identifiez trois types d'acteurs du commerce intérieur présents sur ce marché et précisez leur fonction.
    \item Montrez que ce marché périodique est un nœud de l'interdépendance régionale (produits, origines, destinations).
    \item Pourquoi la périodicité (hebdomadaire) est-elle adaptée à ce type d'espace rural ?
\end{enumerate}
\end{exercice}

\begin{exercice}[Construction graphique : La formation des prix du plantain]
Le tableau ci-dessous donne le prix moyen du régime de plantain (environ 15~kg) à différents stades de la filière, en FCFA (données indicatives 2023).
\begin{center}
\begin{tabularx}{\linewidth}{|l|X|X|}\hline
\textbf{Stade} & \textbf{Prix (FCFA)} & \textbf{Acteur principal} \\ \hline
Producteur (champ, Ouest) & 1~500 & Paysan \\ \hline
Collecteur / Bayam sellam (marché rural) & 2~500 & Bayam sellam \\ \hline
Grossiste (marché Mfoundi, Yaoundé) & 4~000 & Grossiste \\ \hline
Détaillant (étal quartier, Yaoundé) & 6~500 & Détaillante / Vendeuse à la sauvette \\ \hline
Consommateur final & 7~000 & Ménage \\ \hline
\end{tabularx}
\end{center}
\emph{Consignes :}
\begin{enumerate}
    \item Construisez un diagramme en bâtons (barres verticales) représentant l'évolution du prix du champ au consommateur. Soignez le titre, les axes, l'échelle et les étiquettes.
    \item Calculez le coefficient multiplicateur total entre le prix producteur et le prix consommateur.
    \item Expliquez les deux causes principales de cette hausse des prix le long de la chaîne (marges, coûts de transport, pertes, taxes...).
\end{enumerate}
\end{exercice}

\begin{exercice}[Analyse de photographies : Formes du commerce intérieur]
Trois photographies (décrites) sont soumises à votre analyse :
\begin{itemize}
    \item \textbf{Photo A :} Un grand espace découvert en plein air, sous des hangars en tôle ou à même le sol. Des centaines de personnes. Des tas de manioc, maïs, bananes, légumes. Des femmes assises derrière leurs produits. Animations, bruit. Jour de marché hebdomadaire à Bafoussam.
    \item \textbf{Photo B :} Une femme debout à l'arrière d'un camion bâché garé en bordure de route ou dans une cour de marché. Elle négocie des sacs de haricots/maïs avec un paysan. Un carnet de notes à la main. Elle charge la marchandise pour Yaoundé.
    \item \textbf{Photo C :} Un trottoir bitumé en centre-ville (Douala, carrefour Wouri). De jeunes hommes et femmes étalent sur des nattes ou des cartons : chaussures d'occasion, téléphones, accessoires, fruits coupés. Pas de local fixe. Présence occasionnelle de la police municipale.
\end{itemize}
\emph{Consignes :}
\begin{enumerate}
    \item Pour chaque photo, identifiez la forme de commerce intérieur illustrée (marché périodique, bayam sellam, vente à la sauvette).
    \item Précisez le rôle économique et social de chacune de ces trois formes dans le circuit de distribution au Cameroun.
    \item Quelle forme relève de l'informel pur et quelles en sont les conséquences pour la collectivité (recettes fiscales, hygiène, embouteillages) ?
\end{enumerate}
\end{exercice}

\courssub{Je me prépare au Bac}

\begin{exercice}[Sujet de synthèse type Bac : Dynamisme et contraintes]
\textbf{Sujet :} \emph{« Les échanges inter-régionaux au Cameroun : dynamisme et contraintes. »}
\emph{Consignes :} Traitez ce sujet sous la forme d'une dissertation géographique structurée (Introduction problématisée, Deux parties équilibrées avec sous-parties, Conclusion ouverte). Vous mobiliserez vos connaissances du cours, des données chiffrées précises (longueurs de routes, tonnages portuaires, productions régionales) et des exemples localisés (villes, régions, axes, acteurs). La rédaction soignée et la pertinence des croquis mentaux attendus seront valorisées.
\end{exercice}

\begin{exercice}[Étude de documents guidée : L'interdépendance régionale à l'épreuve des infrastructures]
\textbf{Document 1 : Tableau -- État des infrastructures de transport (estimation 2023)}
\begin{center}
\begin{tabularx}{\linewidth}{|l|X|X|}\hline
\textbf{Mode} & \textbf{Caractéristiques principales} & \textbf{Contraintes majeures} \\ \hline
Route (bitumée $\approx$ 6~500~km) & Réseau en étoile autour de Yaoundé/Douala. Axes structurants : N1, N2, N3, N10, N6. & Entretien irrégulier, dégradation rapide (saison des pluies), rackets, insécurité (Grand Nord), coûts de transport élevés. \\ \hline
Chemin de fer (Camrail $\approx$ 1~100~km) & Ligne unique Douala--Yaoundé--Ngaoundéré. Fret : coton, hydrocarbures, bois, minerais. & Matériel vétuste, vitesse faible ($\approx$ 30~km/h), saturation, pas de liaison vers l'Est (Bertoua) ni le Nord-Ouest. \\ \hline
Portuaire (Douala, Kribi, Limbé) & Douala : 1er port CEMAC ($\approx$ 12~Mt/an). Kribi : eau profonde, vrac, conteneurs. & Douala : congestion, accès routier/ferré saturé. Kribi : manque de connectivité ferroviaire achevée vers l'intérieur. \\ \hline
Aérien (Camair-Co, aéroports int.) & 10 aéroports internationaux/régionaux. & Camair-Co en redressement judiciaire, desserte intérieure réduite, coûts élevés, fret aérien marginal. \\ \hline
\end{tabularx}
\end{center}

\textbf{Document 2 : Croquis mental -- Flux inter-régionaux majeurs (description)}
Flèche épaisse Nord (Ngaoundéré, Garoua, Maroua) $\rightarrow$ Sud (Douala, Yaoundé) : Coton, Bétail, Oignon.
Flèche épaisse Ouest (Bafoussam, Bamenda) $\rightarrow$ Yaoundé/Douala : Vivriers (haricot, pomme de terre, chou), Café arabica.
Flèche épaisse Sud/Est (Bertoua, Ebolowa, Sangmélima) $\rightarrow$ Douala/Kribi : Bois, Cacao, Café robusta.
Flèche inverse (Sud $\rightarrow$ Nord/Est) : Produits manufacturés, carburants, engrais, médicaments, riz importé.

\emph{Questions :}
\begin{enumerate}
    \item À partir du Document 1, dressez un bilan contrasté de l'état des infrastructures de transport au Cameroun (atouts / limites) en 4 lignes maximum.
    \item En croisant les Documents 1 et 2, expliquez comment l'état des infrastructures (route, rail, port) influence la compétitivité des flux d'exportation (coton, bois, cacao) et le coût de la vie dans les villes du Nord.
    \item Proposez deux mesures concrètes, réalistes et prioritaires pour améliorer la fluidité des échanges inter-régionaux. Justifiez chaque mesure.
\end{enumerate}
\end{exercice}

\begin{exercice}[Travail Pratique noté : Construction de la carte des flux commerciaux inter-régionaux]
\textbf{Consigne :} À partir des données ci-dessous, construisez et légendez une carte schématique des flux commerciaux inter-régionaux au Cameroun sur le fond de carte fourni (contour du pays, capitales régionales, axes routiers principaux N1, N2, N3, N10, voie ferrée, ports de Douala et Kribi). Respectez la méthode cartographique (propreté, hiérarchie des symboles, légende organisée, titre, échelle, orientation).

\textbf{Données à représenter :}
\begin{itemize}
    \item \textbf{Zones de production (surfaces colorées / hachurées) :}
    \begin{itemize}
        \item Zone coton/bétail/oignon : Adamaoua, Nord, Extrême-Nord (hachures obliques \textcolor{poporange}{orange}).
        \item Zone vivriers divers/café arabica : Ouest, Nord-Ouest (hachures \textcolor{popgreen}{vert}).
        \item Zone cacao/bois/café robusta/plantain : Centre, Sud, Est, Littoral hors villes (hachures \textcolor{poporange!70!popdark}{marron}).
    \end{itemize}
    \item \textbf{Zones de consommation / commercialisation majeure (symboles de surface) :}
    \item Douala et Yaoundé : Cercles rouges proportionnels (grands marchés, ports, industries).
    \item Capitales régionales secondaires (Bafoussam, Bamenda, Garoua, Maroua, Ngaoundéré, Bertoua, Ebolowa) : Cercles rouges moyens.
    \item \textbf{Flux de marchandises (flèches d'épaisseur proportionnelle) :}
    \begin{itemize}
        \item Flux d'exportation (vers ports) : Coton (Nord $\rightarrow$ Douala/Kribi), Bois (Est/Sud $\rightarrow$ Douala/Kribi), Cacao (Centre/Sud/Ouest $\rightarrow$ Douala/Kribi), Aluminium (Ngaoundal $\rightarrow$ Douala). Flèches \textcolor{popblue}{bleues} épaisses.
        \item Flux vivriers inter-régionaux (vers villes) : Haricots/pommes de terre (Ouest $\rightarrow$ Yaoundé/Douala), Bétail/Viande (Nord/Adamaoua $\rightarrow$ Yaoundé/Douala), Oignon (Nord $\rightarrow$ tout le pays), Plantain/Manioc (Centre/Sud $\rightarrow$ Yaoundé/Douala/Nord). Flèches \textcolor{popgreen}{vertes} moyennes.
        \item Flux de produits manufacturés/importés (des ports/villes vers l'intérieur) : Carburants, riz, engrais, biens de consommation (Douala/Yaoundé $\rightarrow$ Nord, Est, Ouest). Flèches \textcolor{poporange}{orange} fines en pointillés.
    \end{itemize}
    \item \textbf{Infrastructures clés à symboliser :} Ports (ancre), Aéroports internationaux (avion), Barrages (éclair), Pipeline (trait discontinu).
\end{itemize}

\textbf{Attendus de la légende (organisée) :}
\begin{enumerate}
    \item Titre de la carte.
    \item Zones de production (3 items).
    \item Pôles de commercialisation / consommation (2 items).
    \item Flux (3 catégories : export, vivriers, manufacturés).
    \item Infrastructures (4 items).
    \item Échelle graphique, Flèche Nord, Auteur, Source, Date.
\end{enumerate}
\end{exercice}

\newpage

\coursec{Corrigés des exercices}

\coursec{Entraînement et évaluation : Les échanges inter-régionaux au Cameroun}

\courssub{Exercices d'application variés}

\begin{corrige}[Exercice 1 — QCM : Infrastructures et acteurs du commerce]
\textbf{Énoncé :} Choisir les bonnes affirmations parmi les propositions suivantes :\\
1. Le réseau routier bitumé dépasse 10~000~km et est régulièrement entretenu.\\
2. Camrail exploite la ligne Douala--Yaoundé--Ngaoundéré ; le fret est sa fonction la plus rentable.\\
3. Les ports de Douala et Kribi sont bien connectés à l'intérieur du pays.\\
4. Les bayam selam assurent la jonction entre marchés ruraux de production et marchés urbains de consommation.\\
5. Les vendeurs à la sauvette relèvent du secteur formel.
\tcblower
\textbf{Solution détaillée :} Réponses exactes : \textbf{2 et 4}.

\begin{enumerate}
    \item \textbf{Fausse.} Le réseau bitumé atteint environ 6~500~km (ordre de grandeur 2023), loin des 10~000~km annoncés, et surtout son entretien est \emph{irrégulier} : de nombreux tronçons se dégradent rapidement, notamment en saison des pluies.
    \item \textbf{Exacte.} Camrail exploite la ligne Douala--Yaoundé--Ngaoundéré (environ 1~100~km). Le fret (coton, hydrocarbures, bois, produits pétroliers) est sa fonction la plus rentable, mais le matériel roulant est vétuste et la vitesse commerciale faible (environ 30~km/h).
    \item \textbf{Fausse.} C'est précisément une des faiblesses majeures : les ports de Douala et de Kribi sont \emph{mal connectés à l'intérieur}. Douala souffre d'accès routier et ferré saturés ; Kribi n'est pas encore relié par une voie ferrée achevée vers l'arrière-pays (hinterland).
    \item \textbf{Exacte.} Les \cle{bayam selam} (« ceux qui achètent et vendent ») achètent les vivres dans les marchés ruraux de production (Ouest, Nord-Ouest, Nord) et les revendent sur les marchés urbains de consommation (Yaoundé, Douala). Ce sont des maillons essentiels de l'interdépendance régionale.
    \item \textbf{Fausse.} Les vendeurs à la sauvette relèvent du secteur \emph{informel} : ils n'ont pas de local fixe, ne paient généralement pas d'impôts, occupent trottoirs et abords de marchés, et sont exposés aux expulsions par la police municipale.
\end{enumerate}

\textbf{Points de vigilance :} ne pas confondre « augmentation du réseau bitumé » (réelle) et « entretien régulier » (faux) ; bien distinguer secteur formel et informel ; orthographe officielle : \textbf{bayam selam}.
\end{corrige}

\begin{corrige}[Exercice 2 — Vrai ou Faux justifié]
\textbf{Énoncé :} Dire si les affirmations suivantes sont vraies ou fausses et justifier brièvement.\\
1. L'interdépendance régionale ne porte que sur les produits vivriers du Nord vers le Sud.\\
2. L'enclavement du Nord et de l'Est est totalement résolu grâce à l'extension du réseau bitumé.\\
3. Les marchés périodiques hebdomadaires sont un rouage central du commerce intérieur.\\
4. La compagnie aérienne nationale Camair-Co assure efficacement le fret intérieur.
\tcblower
\textbf{Solution détaillée :}
\begin{enumerate}
    \item \textbf{Faux.} L'interdépendance porte aussi sur les produits manufacturés, les carburants, les engrais, les médicaments et les produits importés (riz) qui circulent du Sud (ports, villes industrielles) vers le Nord, l'Est et l'Ouest. Les flux sont donc \emph{bidirectionnels} et portent sur des marchandises variées.
    \item \textbf{Faux.} Le réseau bitumé a progressé (environ 4~500~km en 2010 contre 6~500~km en 2023), mais l'entretien irrégulier, les coupures en saison des pluies et l'absence de dessertes modernes maintiennent l'enclavement de nombreuses zones de production du Nord et de l'Est (Bertoua n'est pas reliée au rail, par exemple).
    \item \textbf{Vrai.} Les marchés périodiques hebdomadaires (Bafia, Bafoussam, etc.) concentrent l'offre rurale, permettent la collecte par les bayam selam et redistribuent les produits vivriers vers les villes. Ils sont un rouage central du commerce intérieur.
    \item \textbf{Faux.} Camair-Co est en difficultés financières (redressement) ; sa desserte intérieure est réduite et coûteuse, et le fret aérien reste marginal dans les échanges intérieurs.
\end{enumerate}

\textbf{Points de vigilance :} une justification sans appui sur une donnée du cours (chiffre, exemple localisé) est insuffisante ; éviter les formulations absolues (« totalement », « uniquement ») qui sont presque toujours fausses en géographie ; orthographe : \textbf{bayam selam}.
\end{corrige}

\begin{corrige}[Exercice 3 — Définitions et concepts clés]
\textbf{Énoncé :} Définir précisément les termes suivants : échanges inter-régionaux, commerce intérieur, interdépendance régionale, enclavement, bayam selam.
\tcblower
\textbf{Solution détaillée :}
\begin{enumerate}
    \item \textbf{Échanges inter-régionaux :} ensemble des flux de marchandises (produits agricoles, manufacturés, importés) et de services qui circulent entre les régions du Cameroun, en fonction de leurs spécialisations productives et de leurs besoins complémentaires.
    \item \textbf{Commerce intérieur :} ensemble des opérations d'achat, de transport, de stockage et de revente de marchandises à l'intérieur des frontières nationales, assuré par des acteurs formels (grossistes, centres commerciaux) et informels (bayam selam, vendeurs à la sauvette).
    \item \textbf{Interdépendance régionale :} situation dans laquelle les régions, spécialisées dans des productions différentes (coton et bétail au Nord, vivriers à l'Ouest, cacao et bois au Sud-Est, industrie et services au Littoral-Centre), dépendent les unes des autres pour leur approvisionnement et l'écoulement de leurs produits.
    \item \textbf{Enclavement :} difficulté d'accès d'un espace (zone de production, région) aux réseaux de transport et aux marchés, due à l'absence ou à la dégradation des infrastructures ; il renchérit les coûts et isole les producteurs.
    \item \textbf{Bayam selam :} commerçant(e) du secteur informel qui achète des vivres en gros dans les marchés ruraux de production et les revend sur les marchés urbains de consommation, assurant le lien entre campagnes et villes.
\end{enumerate}

\textbf{Points de vigilance :} une définition doit comporter un terme générique (« ensemble des flux... », « commerçant qui... ») et un élément différenciant ; éviter de définir un mot par lui-même ; respecter l'orthographe \textbf{bayam selam}.
\end{corrige}

\begin{corrige}[Exercice 4 — Calcul direct : Taux de variation du réseau routier]
\textbf{Énoncé :} Le linéaire des routes bitumées est passé d'environ 4~500~km en 2010 à 6~500~km en 2023. Calculer l'augmentation absolue, le taux de variation sur la période et l'augmentation moyenne annuelle.
\tcblower
\textbf{Solution détaillée :}
Données : $L_{2010} = 4~500$ km ; $L_{2023} = 6~500$ km.

\begin{enumerate}
    \item \textbf{Augmentation absolue :}
    \[ \Delta L = 6~500 - 4~500 = 2~000~\text{km}. \]
    Le linéaire bitumé a augmenté de \textbf{2~000~km} entre 2010 et 2023.

    \item \textbf{Taux de variation :}
    \[ T = \frac{L_{2023} - L_{2010}}{L_{2010}} \times 100 = \frac{2~000}{4~500} \times 100 \approx 44,4~\%. \]
    Le réseau bitumé a progressé d'environ \textbf{44,4~\%} sur la période.

    \item \textbf{Augmentation moyenne annuelle :}
    \[ \bar{a} = \frac{2~000}{2023 - 2010} = \frac{2~000}{13} \approx 153,8~\text{km/an}. \]
    Soit environ \textbf{154 km de route bitumée supplémentaire par an} en moyenne.
\end{enumerate}

\textbf{Vérification :} $4~500 \times 1{,}444 \approx 6~500$ \checkmark ; $154 \times 13 = 2~002 \approx 2~000$ \checkmark.

\textbf{Points de vigilance :} le taux de variation se calcule toujours par rapport à la valeur \emph{initiale} (2010), pas à la valeur finale ; ne pas confondre augmentation absolue (en km) et taux (en \%) ; arrondir au dixième de pourcent ou au km près selon l'énoncé.
\end{corrige}

\begin{corrige}[Exercice 5 — Analyse de documents : L'interdépendance Nord-Sud]
\textbf{Énoncé :} À partir de la carte des flux (document fourni), citer deux flux majeurs Nord $\rightarrow$ Sud, expliquer pourquoi Douala et Yaoundé sont les pôles majeurs, et décrire le rôle et les limites du chemin de fer.
\tcblower
\textbf{Solution détaillée :}
\begin{enumerate}
    \item \textbf{Deux flux majeurs Nord $\rightarrow$ Sud :}
    \begin{itemize}
        \item le \textbf{coton} de l'Adamaoua, du Nord et de l'Extrême-Nord vers Douala (transformation et exportation par le port) ;
        \item le \textbf{bétail} (viande, cuir) et l'\textbf{oignon} du Grand Nord vers les grands marchés de consommation de Yaoundé et Douala.
    \end{itemize}
    (On pouvait aussi citer le flux inverse : produits manufacturés, carburants et riz importé du Sud vers le Nord.)

    \item \textbf{Douala et Yaoundé, pôles majeurs :} Douala concentre le premier port de la CEMAC (environ 12~Mt/an, ordre de grandeur), les activités industrielles et les grands marchés ; Yaoundé, capitale et nœud du réseau en étoile (routes N1, N2, N3, rail), est le premier marché de consommation intérieur. Toutes deux polarisent les flux grâce à leurs fonctions portuaires, industrielles, tertiaires et à leur position de tête des axes de transport.

    \item \textbf{Rôle du rail :} la ligne Camrail Douala--Yaoundé--Ngaoundéré évacue la fibre de coton du bassin de production (via Ngaoundéré, terminus nord) vers l'usine et le port de Douala, à moindre coût que la route pour les marchandises lourdes et volumineuses. \textbf{Limite :} le matériel roulant est vétuste, la vitesse commerciale faible (environ 30~km/h), la ligne est unique et saturée, et elle ne dessert ni l'Est ni le Nord-Ouest.
\end{enumerate}

\textbf{Points de vigilance :} citer des flux \emph{précis} (produit + origine + destination), pas des généralités ; une « limite » doit être un fait observable (vétusté, vitesse, saturation), pas un jugement vague.
\end{corrige}

\begin{corrige}[Exercice 6 — Étude de cas : Le marché périodique de Bafia]
\textbf{Énoncé :} À partir de la description du marché de Bafia (jour de marché, acteurs, produits), identifier trois types d'acteurs, montrer que ce marché est un nœud d'interdépendance, et expliquer la pertinence de la périodicité hebdomadaire.
\tcblower
\textbf{Solution détaillée :}
\begin{enumerate}
    \item \textbf{Trois types d'acteurs :}
    \begin{itemize}
        \item les \textbf{paysannes-productrices} des villages environnants : elles vendent directement leur production (plantain, macabo, légumes) — commerce de détail de proximité ;
        \item les \textbf{bayam selam} venues de Yaoundé : elles collectent en gros les vivres dans les zones de production (ici le Nord-Ouest) et les revendent sur les marchés de consommation — fonction d'intermédiation et de transport ;
        \item les \textbf{revendeurs de produits manufacturés} venus de Douala : ils diffusent dans l'espace rural les biens industriels et importés (savon, huile, tissus, quincaillerie).
    \end{itemize}
    (Les bouchers revendant le bétail de l'Adamaoua illustrent un quatrième type : intermédiaires spécialisés.)

    \item \textbf{Un nœud d'interdépendance :} le marché de Bafia croise des produits d'origines régionales très diverses — vivriers du Nord-Ouest (manioc, maïs, haricots), viande de l'Adamaoua, produits manufacturés de Douala — et des productions locales du Centre (plantain, macabo). Il reçoit des flux du Nord, de l'Ouest et du Littoral, et redistribue vers Yaoundé (120~km) et les campagnes environnantes : c'est un point de rencontre entre espaces de production et espaces de consommation.

    \item \textbf{Pertinence de la périodicité hebdomadaire :} dans un espace rural à faible densité et à pouvoir d'achat limité, une demande quotidienne serait insuffisante pour rentabiliser les déplacements des commerçants. La concentration hebdomadaire de l'offre et de la demande en un seul jour permet d'atteindre un seuil de rentabilité, d'organiser les tournées des bayam selam et de rythmer la vie économique et sociale de toute une aire d'influence.
\end{enumerate}

\textbf{Points de vigilance :} bien distinguer la \emph{fonction} de chaque acteur (produire, collecter-transporter, redistribuer) ; pour la question 2, nommer explicitement produits, origines et destinations ; orthographe : \textbf{bayam selam}.
\end{corrige}

\begin{corrige}[Exercice 7 — Construction graphique : La formation des prix du plantain]
\textbf{Énoncé :} À partir du tableau de prix (Producteur 1500, Collecteur 2500, Grossiste 4000, Détaillant 6500, Consommateur 7000 FCFA/régime de 15 kg), construire un diagramme en bâtons, calculer le coefficient multiplicateur total et expliquer deux causes principales de la hausse.
\tcblower
\textbf{Solution détaillée :}
\begin{enumerate}
    \item \textbf{Diagramme en bâtons :}

\begin{popfigure}
\begin{center}
\begin{tikzpicture}
\begin{axis}[
    width=0.95\linewidth, height=7cm,
    ybar, bar width=0.9cm,
    ymin=0, ymax=8000,
    ylabel={Prix (FCFA / régime de 15 kg)},
    symbolic x coords={Producteur, Collecteur, Grossiste, D\'etailant, Consommateur},
    xtick=data,
    x tick label style={rotate=20, anchor=east, font=\small},
    ytick={0,1000,2000,3000,4000,5000,6000,7000,8000},
    grid=major, grid style={popline!40},
    nodes near coords, nodes near coords style={font=\small, color=popdark},
    every node near coord/.append style={/pgf/number format/fixed},
]
\addplot[fill=popgreen, draw=popdark] coordinates {
    (Producteur,1500) (Collecteur,2500) (Grossiste,4000) (D\'etailant,6500) (Consommateur,7000)
};
\end{axis}
\end{tikzpicture}
\end{center}
\legende{Formation du prix d'un régime de plantain ($\approx$ 15 kg) du champ (Ouest) au consommateur (Yaoundé), en FCFA. Données indicatives 2023. Chaque bâton correspond à un stade de la filière ; la valeur est reportée au sommet.}
\end{popfigure}

    \item \textbf{Coefficient multiplicateur total :}
    \[ k = \frac{P_{\text{consommateur}}}{P_{\text{producteur}}} = \frac{7~000}{1~500} \approx 4,67. \]
    Le prix final est environ \textbf{4,7 fois} le prix payé au producteur (soit une hausse d'environ $367~\%$).

    \item \textbf{Deux causes principales de la hausse :}
    \begin{itemize}
        \item les \textbf{coûts de transport et de logistique} : carburant, état des routes, péages et tracasseries, chargement/déchargement entre le champ, les marchés ruraux et les marchés urbains ;
        \item les \textbf{marges des intermédiaires} successifs (bayam selam, grossiste, détaillant), auxquelles s'ajoutent les \textbf{pertes} (pourrissement des vivres périssables) et les \textbf{taxes} de marché.
    \end{itemize}
\end{enumerate}

\textbf{Vérification :} $1~500 \times 4{,}67 \approx 7~005 \approx 7~000$ \checkmark.

\textbf{Points de vigilance :} le coefficient multiplicateur est un rapport (sans unité), ne pas le confondre avec le taux de variation en \% ; le diagramme doit comporter titre, axes nommés avec unité, et échelle régulière ; dans le code TikZ/pgfplots, les nombres décimaux s'écrivent avec un \textbf{point} (ex: 4.67) mais ici les coordonnées sont des entiers.
\end{corrige}

\begin{corrige}[Exercice 8 — Analyse de photographies : Formes du commerce intérieur]
\textbf{Énoncé :} Identifier les trois formes de commerce illustrées (Photos A, B, C), décrire leurs rôles économiques et sociaux, et préciser laquelle relève de l'informel pur avec ses conséquences.
\tcblower
\textbf{Solution détaillée :}
\begin{enumerate}
    \item \textbf{Identification :}
    \begin{itemize}
        \item \textbf{Photo A :} le \textbf{marché périodique} (jour de marché hebdomadaire à Bafoussam) — grand espace découvert, affluence concentrée un jour donné, produits vivriers en tas ;
        \item \textbf{Photo B :} la \textbf{bayam selam} — achat en gros au producteur, chargement d'un camion pour la ville (Yaoundé), carnet de commandes ;
        \item \textbf{Photo C :} la \textbf{vente à la sauvette} — étalage sur trottoir, sans local fixe, exposition aux interventions de la police municipale.
    \end{itemize}

    \item \textbf{Rôles économiques et sociaux :}
    \begin{itemize}
        \item le marché périodique concentre l'offre rurale, fixe les prix par confrontation de l'offre et de la demande, et constitue un lieu de vie sociale et d'information ;
        \item la bayam selam assure la \emph{jonction campagne-ville} : elle collecte, transporte et approvisionne les marchés urbains en vivres, offrant un débouché aux producteurs et un revenu (souvent féminin) ;
        \item la vente à la sauvette offre aux citadins des produits accessibles à bas prix et de proximité, et procure un emploi de survie à une main-d'œuvre jeune et peu qualifiée.
    \end{itemize}

    \item \textbf{Informel pur :} la vente à la sauvette (photo C), qui échappe à l'enregistrement et à la fiscalité. Conséquences pour la collectivité : \textbf{perte de recettes fiscales} pour l'État et les communes, problèmes d'\textbf{hygiène} (produits non contrôlés, déchets), \textbf{embouteillages} et occupation illégale de l'espace public, conflits récurrents avec les autorités municipales. (Les bayam selam relèvent aussi largement de l'informel, mais elles paient souvent des taxes de marché.)
\end{enumerate}

\textbf{Points de vigilance :} ne pas réduire l'informel à un « problème » : il assure l'essentiel de la distribution vivrière et de l'emploi urbain ; nuancer le jugement est valorisé ; orthographe : \textbf{bayam selam}.
\end{corrige}

\courssub{Sujet de synthèse type Baccalauréat}

\begin{corrige}[Exercice 9 — Sujet de synthèse : Dynamisme et contraintes des échanges inter-régionaux]
\textbf{Énoncé :} « Les échanges inter-régionaux au Cameroun sont dynamiques mais contraints. » Montrer la validité de cette affirmation en structurant votre réponse en deux parties équilibrées. Un croquis mental légendé des flux est attendu.
\tcblower
\textbf{Corrigé structuré (plan détaillé attendu).}

\textbf{Introduction.} Amorce : au marché Mokolo de Yaoundé, l'oignon de Maroua côtoie le haricot de Bamenda et le riz importé par Douala. Définition des termes : les échanges inter-régionaux sont les flux de marchandises entre régions aux spécialisations complémentaires. Problématique : dans quelle mesure le dynamisme des échanges inter-régionaux, moteur de l'intégration nationale, est-il freiné par des contraintes structurelles ? Annonce du plan en deux parties.

\textbf{I. Des échanges dynamiques, fondés sur la complémentarité des régions.}
\begin{itemize}
    \item \emph{I.1. Une interdépendance régionale forte :} le Grand Nord fournit coton, bétail, oignon ; l'Ouest et le Nord-Ouest, les vivriers et le café arabica ; le Centre-Sud-Est, le cacao, le bois, le plantain ; Douala et Yaoundé concentrent industrie, services et consommation, et renvoient produits manufacturés, carburants et biens importés.
    \item \emph{I.2. Un commerce intérieur bien organisé :} réseau de marchés périodiques hebdomadaires (Bafia, Bafoussam), rôle central des bayam selam entre campagnes et villes, vente à la sauvette et centres commerciaux modernes dans les grandes villes.
    \item \emph{I.3. Des infrastructures en progrès :} réseau bitumé passé d'environ 4~500~km (2010) à 6~500~km (2023), soit +44~\% ; port de Kribi en eau profonde ; axes structurants N1, N2, N3, N10.
\end{itemize}

\textbf{II. Des contraintes qui freinent les échanges.}
\begin{itemize}
    \item \emph{II.1. Des infrastructures insuffisantes et dégradées :} entretien routier irrégulier, chemin de fer vétuste (Camrail, environ 1~100~km, 30~km/h), ports mal connectés à l'intérieur (Kribi sans liaison ferrée achevée), Camair-Co en difficultés.
    \item \emph{II.2. Un enclavement coûteux :} zones de production du Nord et de l'Est mal desservies, renchérissement des prix (le plantain est vendu près de 4,7 fois son prix au champ), pertes de produits périssables.
    \item \emph{II.3. D'autres obstacles :} tracasseries et rackets sur les axes, insécurité dans le Grand Nord, poids de l'informel (pertes fiscales, hygiène, congestion urbaine).
\end{itemize}

\textbf{Conclusion.} Bilan : les échanges inter-régionaux tissent l'intégration nationale mais restent bridés par le déficit d'infrastructures. Ouverture : les projets en cours (autoroute Yaoundé--Douala, extension de Kribi, réhabilitation ferroviaire) peuvent-ils transformer le Cameroun en « grenier et carrefour » de l'Afrique centrale ?

\textbf{Croquis mental légendé (description) :} Fond de carte simplifié du Cameroun. Zones de production hachurées (Nord : coton/bétail/oignon ; Ouest/N-O : vivriers/café ; Sud-Est : cacao/bois). Pôles majeurs (Douala, Yaoundé) en cercles rouges proportionnels. Flèches bleues épaisses (exportations coton/bois/cacao vers ports), flèches vertes (vivriers vers villes), flèches orange pointillées (manufacturés/importés vers intérieur). Axes N1, N3, N10 et rail Camrail tracés. Titre, légende, échelle, Nord.

\textbf{Barème indicatif (sur 20) :} Introduction problématisée : 3 pts ; Partie I (dynamisme, exemples localisés et chiffrés) : 6 pts ; Partie II (contraintes, données précises) : 6 pts ; Conclusion ouverte : 2 pts ; Croquis mental légendé des flux : 2 pts ; Qualité de la langue et de l'organisation : 1 pt.

\textbf{Points de vigilance :} deux parties \emph{équilibrées} ; chaque affirmation appuyée par un exemple localisé (ville, région, axe) ou un chiffre ; le croquis mental (fond de carte simplifié, flèches proportionnelles, légende) est explicitement attendu et valorisé ; éviter le catalogue de notions sans démonstration ; orthographe : \textbf{bayam selam}.
\end{corrige}

\courssub{Étude de documents guidée (Type épreuve)}

\begin{corrige}[Exercice 10 — Étude de documents guidée : Infrastructures et interdépendance]
\textbf{Énoncé :} Document 1 : Carte des infrastructures de transport. Document 2 : Tableau des flux de marchandises Nord/Sud.\\
1. Faire un bilan contrasté des infrastructures de transport (4 lignes max).\\
2. Montrer comment l'état des infrastructures influence les flux et les prix (croiser Doc 1 et Doc 2).\\
3. Proposer deux mesures prioritaires pour améliorer la situation.
\tcblower
\textbf{Solution détaillée :}
\begin{enumerate}
    \item \textbf{Bilan contrasté (4 lignes maximum) :} le Cameroun dispose d'atouts réels — un réseau routier en étoile structuré autour de Yaoundé et Douala (environ 6~500~km bitumés), une ligne ferroviaire Douala--Ngaoundéré, le premier port de la CEMAC à Douala et un port en eau profonde à Kribi — mais ces atouts sont minés par l'entretien irrégulier des routes, la vétusté du rail (environ 30~km/h), la mauvaise connexion des ports à l'arrière-pays et les difficultés de Camair-Co.

    \item \textbf{Influence des infrastructures sur les flux et les prix :}
    \begin{itemize}
        \item \emph{Compétitivité des exportations :} le coton du Nord transite par le rail vétuste et lent jusqu'à Douala, ce qui allonge les délais et renchérit les coûts ; le bois de l'Est (Bertoua) n'a pas de liaison ferrée et part par route dégradée ; la congestion de Douala retarde l'embarquement du cacao. Kribi, malgré son eau profonde, ne peut drainer pleinement les flux faute de voie ferrée achevée vers l'intérieur. Résultat : des produits camerounais plus chers et moins compétitifs sur le marché mondial.
        \item \emph{Coût de la vie dans les villes du Nord :} les flux inverses (carburants, riz importé, engrais, produits manufacturés, Document 2) empruntent la route N1, longue, dégradée et ponctuée de tracasseries ; chaque surcoût de transport se répercute sur les prix finaux à Garoua, Maroua ou Ngaoundéré, où le coût de la vie des produits manufacturés est sensiblement plus élevé qu'à Douala.
    \end{itemize}

    \item \textbf{Deux mesures prioritaires (exemples de réponses) :}
    \begin{itemize}
        \item \textbf{Réhabiliter et moderniser la voie ferrée} (renouvellement du matériel roulant de Camrail, extension vers Bertoua et connexion ferroviaire du port de Kribi) : le rail est le mode le plus économique pour les produits lourds (coton, bois, minerais) et désengorgerait la route N1.
        \item \textbf{Programmer un entretien routier régulier et pérenne} (financement dédié, entretien des axes N1 et N10 avant la saison des pluies) : une route entretenue réduit la durée des trajets, les pertes de vivres périssables et donc les prix payés par les consommateurs.
    \end{itemize}
    (Autres mesures acceptables si justifiées : désengorger le port de Douala, sécuriser les axes du Grand Nord, formaliser progressivement les circuits de distribution.)
\end{enumerate}

\textbf{Barème indicatif (sur 20) :} Q1 : 4 pts (bilan équilibré atouts/limites, concision) ; Q2 : 10 pts (croisement explicite des deux documents, mécanismes économiques expliqués, exemples localisés) ; Q3 : 6 pts (3 pts par mesure : réalisme + justification).

\textbf{Points de vigilance :} la question 2 exige de \emph{croiser} les documents (citer les deux), pas de les résumer séparément ; chaque mesure proposée doit être justifiée par un effet attendu précis ; respecter la contrainte de longueur de la question 1.
\end{corrige}

\coursec{TP 4 : Carte des flux commerciaux inter-régionaux au Cameroun}

\begin{corrige}[Exercice 11 — TP noté : Carte des flux commerciaux inter-régionaux]
\textbf{Consignes du TP :} Sur un fond de carte simplifié du Cameroun, représenter : les zones de production (coton/bétail/oignon au Nord, vivriers/café à l'Ouest/Nord-Ouest, cacao/bois/plantain au Centre-Sud-Est) ; les pôles de commercialisation (Douala, Yaoundé, capitales régionales) ; les flux d'exportation, les flux vivriers, les flux de produits manufacturés/importés ; les axes routiers principaux (N1, N3, N10) et la voie ferrée. Légender soigneusement.
\tcblower
\textbf{Corrigé-type :} la carte ci-dessous constitue le modèle attendu. Le fond de carte est une carte générale du Cameroun ; ce qui est évalué, c'est la \emph{méthode} : hiérarchie des symboles, flèches proportionnelles, légende complète et organisée.

\begin{popfigure}
\begin{center}
\includegraphics[width=\linewidth,height=9.5cm,keepaspectratio]{N01/cmr_map.jpg}
\end{center}
\legende{Fond de carte du Cameroun (régions, villes, routes N1, N3, N10 et voie ferrée) servant de modèle. La carte attendue y situe : les zones de production (Grand Nord : coton, bétail, oignon ; Ouest et Nord-Ouest : vivriers, café arabica ; Centre--Sud--Est : cacao, bois, plantain), Douala et Yaoundé comme pôles majeurs et les capitales régionales comme pôles secondaires, puis les flux (exportations : coton, bois, cacao vers le port ; vivriers vers les villes ; produits manufacturés et importés vers l'intérieur) avec une légende complète. \\ Source : Wikimedia Commons, CIA, domaine public.}
\end{popfigure}

\textbf{Barème indicatif (sur 20) :}
\begin{itemize}
    \item Fond de carte correct (contour, villes bien placées, axes, ports) : 4 pts ;
    \item Zones de production hachurées et localisées correctement : 3 pts ;
    \item Pôles de commercialisation hiérarchisés (cercles proportionnels) : 2 pts ;
    \item Flux correctement orientés, épaisseurs proportionnelles, couleurs respectées : 5 pts ;
    \item Légende complète et organisée (titre, 3 zones, 2 pôles, 3 catégories de flux, infrastructures, échelle, nord, source, date) : 4 pts ;
    \item Propreté et soin général : 2 pts.
\end{itemize}

\textbf{Points de vigilance (erreurs fréquentes à signaler aux élèves) :}
\begin{itemize}
    \item \textbf{Sens des flèches :} l'exportation va \emph{vers} les ports ; les produits manufacturés partent \emph{des} ports et grandes villes vers l'intérieur. Une flèche inversée coûte des points.
    \item \textbf{Proportionnalité :} les flèches d'exportation doivent être visiblement plus épaisses que les flux de produits manufacturés.
    \item \textbf{Légende :} tout symbole utilisé sur la carte doit figurer dans la légende, et réciproquement ; ne jamais oublier titre, échelle et orientation.
    \item \textbf{Localisation :} Garoua et Maroua au nord de Ngaoundéré ; Bafoussam à l'ouest, entre Douala et Bamenda ; Bertoua à l'est de Yaoundé. Vérifier la position relative des villes avant de tracer les flux.
    \item \textbf{Orthographe :} \textbf{bayam selam} (et non sellam).
\end{itemize}
\end{corrige}

\newpage

\coursec{Fiche bilan}

\begin{popfigure}
\begin{tikzpicture}[mindmap, grow cyclic, every node/.style={concept, font=\small}, concept color=popblue!60, level 1 concept/.append style={level distance=4.6cm, sibling angle=51.5, font=\small}, level 2 concept/.append style={level distance=2.9cm, font=\scriptsize}]
\node[concept, concept color=popblue!80, text=white, font=\bfseries\small]{Échanges inter-régionaux au Cameroun}
  child[concept color=popgreen!70]{ node {Facteurs de développement}
    child { node {Routes bitumées en hausse, entretien irrégulier} }
    child { node {Chemin de fer vétuste (Camrail)} }
    child { node {Ports mal connectés à l'intérieur} }
    child { node {Aéroports, Camair-Co en difficulté} }
  }
  child[concept color=poporange!80]{ node {Interdépendance des régions}
    child { node {Complémentarités agro-écologiques} }
    child { node {Villes approvisionnées par les campagnes} }
  }
  child[concept color=poppurple!70]{ node {Acteurs et circuits}
    child { node {Bayam selam (revendeurs de vivres)} }
    child { node {Commerçants grossistes et détaillants} }
    child { node {Transporteurs} }
  }
  child[concept color=poppink!70]{ node {Formes du commerce intérieur}
    child { node {Marchés périodiques} }
    child { node {Vente à la sauvette} }
    child { node {Centres commerciaux} }
  }
  child[concept color=popgold!80]{ node {Difficultés}
    child { node {Enclavement, routes dégradées} }
    child { node {Coût élevé du transport} }
    child { node {Racket, tracasseries, postes de contrôle} }
    child { node {Informalité, pertes post-récolte} }
  }
  child[concept color=popblue!50]{ node {Flux majeurs}
    child { node {Nord $\rightarrow$ Sud : bétail, oignons, céréales} }
    child { node {Ouest $\rightarrow$ villes : vivres, café} }
    child { node {Sud $\rightarrow$ Nord : produits manufacturés, importés} }
  }
  child[concept color=popgreen!45]{ node {Enjeu : intégration nationale}
    child { node {Unité du marché national} }
    child { node {Solidarité économique entre régions} }
  };
\end{tikzpicture}
\legende{Carte mentale de la leçon : les échanges inter-régionaux au Cameroun (facteurs, acteurs, formes, flux, difficultés et enjeux d'intégration nationale).}
\end{popfigure}

\begin{bilan}
\textbf{Définitions clés à connaître par cœur}
\begin{itemize}
  \item \cle{Échange inter-régional} : mouvement de biens, de services (et de personnes) entre les régions du Cameroun, fondé sur leurs complémentarités.
  \item \cle{Commerce intérieur} : ensemble des opérations d'achat et de vente de marchandises à l'intérieur des frontières du pays.
  \item \cle{Marché périodique} : marché qui se tient à dates fixes (hebdomadaires le plus souvent), surtout en milieu rural (ex. marchés du Grand Nord).
  \item \cle{Bayam selam} : revendeur(euse) de vivres frais qui achète en gros dans les zones de production rurales et revend au détail sur les marchés urbains ; maillon essentiel du ravitaillement des villes.
  \item \cle{Vendeur à la sauvette} : commerçant ambulant informel, sans emplacement fixe ni autorisation, actif surtout dans les grandes villes (Douala, Yaoundé).
  \item \cle{Centre commercial} : grand ensemble moderne de distribution (supermarchés, boutiques, services) concentré en un lieu, ex. Douala Grand Mall, carrefours commerciaux de Yaoundé.
  \item \cle{Interdépendance} : situation où chaque région dépend des autres pour se procurer ce qu'elle ne produit pas et écouler ses surplus.
\end{itemize}

\textbf{Repères chiffrés et localisations (ordres de grandeur prudents)}
\begin{center}
\begin{tabularx}{\linewidth}{|l|X|}\hline
\rowcolor{popblueL} \textbf{Élément} & \textbf{Repère} \\ \hline
Réseau routier & environ \num{120000} km au total, dont moins de 10 \% bitumés (ordre de grandeur) \\ \hline
Chemin de fer & environ \num{1000} km (lignes Douala--Yaoundé--Ngaoundéré exploitées par Camrail), infrastructure vétuste \\ \hline
Ports & Douala (principal, plus de 90 \% du trafic), Kribi (port en eau profonde), mal connectés à l'hinterland \\ \hline
Aéroports internationaux & Douala, Yaoundé-Nsimalen, Garoua ; compagnie nationale Camair-Co en difficultés financières \\ \hline
Flux Nord $\rightarrow$ Sud & bétail (Adamaoua, Nord, Extrême-Nord), oignons, céréales, coton \\ \hline
Flux Ouest $\rightarrow$ villes & vivres (pommes de terre, bananes, légumes), café, via les bayam selam \\ \hline
Flux Sud $\rightarrow$ Nord & produits manufacturés et importés (via Douala), sel, tissus, carburant \\ \hline
\end{tabularx}
\end{center}

\textbf{Méthodes express}
\begin{itemize}
  \item \emph{Commenter un document} : identifier la nature du document $\rightarrow$ localiser $\rightarrow$ décrire (du général au particulier) $\rightarrow$ expliquer (facteurs) $\rightarrow$ conclure sur l'enjeu d'intégration nationale.
  \item \emph{Construire un croquis de flux} : fond de carte simplifié du Cameroun $\rightarrow$ zones de production (couleurs) $\rightarrow$ zones de consommation (points/villes) $\rightarrow$ flèches d'épaisseur proportionnelle aux flux $\rightarrow$ titre + légende + nord.
  \item \emph{Analyser une difficulté} : toujours relier la cause (infrastructure, coût, tracasserie) à la conséquence (prix élevés, pertes, enclavement).
\end{itemize}
\end{bilan}

\begin{aretenir}[Checklist avant le Bac]
\begin{itemize}
  \item $\square$ Je sais définir : échange inter-régional, commerce intérieur, marché périodique, bayam selam, vendeur à la sauvette, centre commercial.
  \item $\square$ Je cite les quatre types d'infrastructures et leur état (routes, rail, ports, aérien).
  \item $\square$ J'explique l'interdépendance des régions par les complémentarités agro-écologiques et industrielles.
  \item $\square$ Je donne au moins trois exemples précis de flux inter-régionaux avec leurs produits et leurs directions.
  \item $\square$ Je décris le rôle des bayam selam dans le ravitaillement des villes.
  \item $\square$ Je distingue commerce formel (centres commerciaux, boutiques) et commerce informel (marchés, sauvette).
  \item $\square$ J'énumère au moins quatre difficultés du commerce intérieur et leurs conséquences sur les prix.
  \item $\square$ Je sais construire et légender un croquis de flux commerciaux (titre, légende, nord, flèches proportionnelles).
  \item $\square$ Je relie chaque difficulté à une solution possible (réhabilitation routière et ferroviaire, lutte contre le racket, modernisation des marchés).
  \item $\square$ Je conclus toujours en rattachant les échanges inter-régionaux à l'intégration nationale.
\end{itemize}
\end{aretenir}

\findecours

\end{document}
